% !TeX root = gadamer.tex
\chapter{Het transcenderen van de esthetische dimensie}

\section{De betekenis van de humanistische traditie voor de geesteswetenschappen}

\subsection{Het probleem van de methode}

De \textit{Geisteswissenschaften} ontwikkelden zich in de negentiende eeuw onder sterke invloed van het voorbeeld van de natuurwetenschappen. Wanneer men nadacht over hun logische en wetenschappelijke grondslag, nam men meestal de natuurwetenschappen als model. Dat blijkt ook uit de geschiedenis van het woord \textit{Geisteswissenschaften}. Hoewel deze term pas in de negentiende eeuw zijn huidige betekenis kreeg, werd hij al snel gebruikt om een groep wetenschappen aan te duiden die zich naar het voorbeeld van de natuurwetenschappen wilden begrijpen. Daardoor raakte de oorspronkelijke, idealistische betekenis van \textit{Geist} op de achtergrond. \textit{Geist} verwees immers niet alleen naar empirische kennis, maar ook naar vorming, cultuur en menselijke betekenis.

De verbreiding van het woord \textit{Geisteswissenschaften} hing nauw samen met de Duitse vertaling van John Stuart Mills \textit{Logic}. In een aanvullend deel bij deze vertaling werd uitgelegd hoe de inductieve logica ook kon worden toegepast op de zogenaamde `moral sciences''. Deze `moral sciences'' werden in het Duits aangeduid als \textit{Geisteswissenschaften}.

Bij Mill, en na hem bij Hume, was het niet vanzelfsprekend dat de humanistische wetenschappen een eigen manier van denken of een eigen logica hadden. Integendeel: de inductieve methode, die centraal stond in de experimentele natuurwetenschappen, werd ook voor de menswetenschappen als de juiste methode gezien. Volgens deze opvatting moesten ook de menswetenschappen zoeken naar overeenkomsten, regelmatigheden en wetmatigheden. Op basis daarvan zouden zij individuele verschijnselen en historische processen kunnen verklaren en misschien zelfs voorspellen, net zoals de natuurwetenschappen dat doen.

Dat dit bij sociale en morele verschijnselen moeilijker is dan in de natuurwetenschappen, werd niet gezien als een probleem van de methode zelf. Het lag volgens deze opvatting vooral aan het feit dat men over te weinig gegevens beschikte. De situatie werd bijvoorbeeld vergeleken met de meteorologie: ook daar zijn voorspellingen moeilijk, niet omdat de methode verkeerd is, maar omdat de gegevens complex en onvolledig zijn. De inductieve methode leek daardoor vrij van metafysische aannames. Of men nu wel of niet in de vrije wil geloofde, sociale voorspellingen bleven in principe mogelijk, zolang men maar uitging van vastgestelde regelmatigheden.

Toch doet men de humanistische wetenschappen geen recht wanneer men ze beoordeelt alsof zij vooral wetmatigheden moeten ontdekken. De ervaring van de sociaal-historische wereld kan niet simpelweg tot wetenschap worden gemaakt door de inductieve methode van de natuurwetenschappen toe te passen. Historisch onderzoek maakt natuurlijk wel gebruik van algemene ervaringen en vergelijkingen. Maar het doel ervan is niet om een concreet verschijnsel alleen maar te zien als een voorbeeld van een algemene regel.

Het gaat er in de humanistische wetenschappen juist om een verschijnsel zelf te \textit{begrijpen}, in zijn unieke en historische werkelijkheid. Een individuele gebeurtenis, persoon, volk of staat dient dus niet alleen als bewijs voor een algemene wet. De vraag is eerder: \textit{hoe} is dit specifieke verschijnsel geworden wat het is? Hoe is deze persoon, dit volk of deze staat tot zijn huidige vorm gekomen? En algemener: hoe heeft het zover kunnen komen?

Daarmee ontstaat een belangrijke vraag. Wat voor soort kennis is dit, die begrijpt \textit{waarom} iets is zoals het is, doordat zij begrijpt \textit{hoe} het zo is geworden? En wat betekent het woord `wetenschap'' in dit verband? Zelfs wanneer men erkent dat deze kennis iets anders wil dan de natuurwetenschappen, blijft men vaak geneigd de humanistische wetenschappen negatief te beschrijven. Men noemt ze dan bijvoorbeeld `onnauwkeurige wetenschappen'', omdat ze niet dezelfde strengheid zouden hebben als de natuurwetenschappen.

Een voorbeeld hiervan is Hermann von Helmholtz. In zijn beroemde rede van 1862 benadrukte hij weliswaar het grote menselijke belang van de \textit{Geisteswissenschaften}. Toch bleef hij ze logisch gezien beschrijven vanuit het ideaal van de natuurwetenschappen. Hij maakte onderscheid tussen twee soorten inductie: een logische inductie en een artistiek-intu"itieve inductie. Maar dat onderscheid was volgens Gadamer eerder psychologisch dan werkelijk logisch.

Volgens Helmholtz gebruiken beide soorten wetenschappen inductieve redeneringen. Het verschil is alleen dat de natuurwetenschapper bewust en rationeel tot zijn conclusies komt, terwijl de humanistische wetenschapper vaak onbewust werkt. Die laatste zou afhankelijk zijn van psychologische vermogens zoals tact, een goed geheugen en het vermogen om gezag te erkennen. Daarmee worden de humanistische wetenschappen dus niet vanuit hun eigen aard begrepen, maar slechts als een minder expliciete of minder bewuste vorm van wetenschappelijk redeneren.

Ook Wilhelm Dilthey stond sterk onder invloed van Mills empirisme, al behield hij tegelijk iets van de romantisch-idealistische betekenis van \textit{Geist}. Dilthey zag heel goed dat de historische school iets anders deed dan het natuurwetenschappelijke en natuurrechtelijke denken. Hij meende daarom dat hij boven het Engelse empirisme uitkwam. Toch bleef ook hij zich ori"enteren op het model van de natuurwetenschappen, zelfs wanneer hij de zelfstandigheid van de \textit{Geisteswissenschaften} probeerde te verdedigen.

Dat blijkt bijvoorbeeld uit zijn necrologie op Wilhelm Scherer. Daarin benadrukte Dilthey dat Scherers werkwijze geleid werd door de geest van de natuurwetenschappen. Dilthey beschrijft de moderne wetenschapper als iemand voor wie de wereld van de voorouders niet langer een levende thuisbasis is voor geest en hart, maar een historisch object. Wetenschappelijke kennis lijkt dus te vereisen dat men afstand neemt van het eigen leven en van de eigen geschiedenis. Pas door die afstand kan de geschiedenis tot object van onderzoek worden.

Hoewel Dilthey de zelfstandigheid van de humanistische wetenschappen wilde verdedigen, bleef de moderne wetenschappelijke methode voor hem uiteindelijk de algemene norm. De natuurwetenschappen belichaamden die norm alleen op de duidelijkste manier. De \textit{Geisteswissenschaften} kregen daardoor bij hem geen werkelijk eigen methode.

Precies daar ligt het probleem. Waarin bestaat de wetenschappelijkheid van de humanistische wetenschappen, als die niet ligt in een methode die vergelijkbaar is met die van de natuurwetenschappen? Helmholtz wees al op zaken als geheugen, gezag en tact. Zulke elementen spelen in de natuurwetenschappen nauwelijks een rol, maar lijken in de humanistische wetenschappen juist onmisbaar. Maar waar komt dat tactgevoel vandaan? Hoe leert men het? En vormen juist zulke vermogens misschien de eigenlijke grondslag van de \textit{Geisteswissenschaften}, eerder dan een formele methode?

Deze vragen maken de humanistische wetenschappen tot een filosofisch probleem. Zij passen namelijk niet goed binnen het moderne begrip van wetenschap, dat vooral door de natuurwetenschappen is gevormd.

De antwoorden van Helmholtz en zijn tijdgenoten zijn daarom onvoldoende. Zij gingen uit van het natuurwetenschappelijke model van kennis en probeerden het bijzondere karakter van de humanistische wetenschappen vervolgens te verklaren door te wijzen op een artistiek element: gevoel, intu"itie of artistieke inductie. Maar daarmee wordt niet werkelijk duidelijk wat de humanistische wetenschappen tot wetenschap maakt.

Bovendien had Helmholtz zelf een eenzijdig beeld van de natuurwetenschappen. Hij geloofde niet in `plotselinge inzichten'' of `inspiraties'' en zag wetenschappelijk werk vooral als het bewust trekken van zeer strenge conclusies. Voor hem waren de inductieve natuurwetenschappen het hoogste voorbeeld van wetenschappelijke methode.

Toch erkende Helmholtz dat historische kennis op een andere soort ervaring berust dan het onderzoek naar natuurwetten. Daarom probeerde hij te verklaren waarom inductie in historisch onderzoek onder andere voorwaarden plaatsvindt dan in de natuurwetenschappen. Hij beriep zich daarbij op Kants onderscheid tussen natuur en vrijheid. In de historische wereld, zo meende hij, gelden geen natuurwetten, maar praktische wetten die mensen vrijwillig aanvaarden: geboden en normen. De wereld van de menselijke vrijheid vertoont daarom niet dezelfde regelmaat zonder uitzonderingen als de natuur.

Volgens Gadamer is deze redenering echter niet overtuigend. Helmholtz gebruikt Kants onderscheid tussen natuur en vrijheid op een manier die niet goed past bij Kants eigen bedoeling. Ook past dit onderscheid niet goed bij de logica van inductie zelf. Mill was op dit punt consequenter: hij sloot het probleem van de vrijheid eenvoudigweg methodisch uit.

De humanistische wetenschappen zagen zichzelf overigens helemaal niet als minderwaardig ten opzichte van de natuurwetenschappen. Integendeel: zij waren verbonden met het intellectuele erfgoed van het Duitse classicisme en beschouwden zichzelf juist als de ware erfgenamen van het humanisme. Het Duitse classicisme had niet alleen geleid tot een vernieuwing van literatuur en esthetische kritiek. Het had ook het oudere barokke ideaal van smaak en het rationalisme van de Verlichting achter zich gelaten. Tegelijk gaf het een nieuwe inhoud aan het ideaal van menselijkheid en verlichte rede.

Vooral Herder speelde hierin een belangrijke rol. Hij doorbrak het perfectionisme van de Verlichting met zijn nieuwe ideaal van de ``cultivatie van de mens'' of \textit{Bildung zum Menschen}. Daarmee legde hij de basis voor de bloei van de historische wetenschappen in de negentiende eeuw. Het begrip \textit{Bildung} — dat men kan vertalen als vorming, zelfvorming, opvoeding of cultuur — werd in deze periode van groot belang. Misschien was het zelfs de belangrijkste gedachte van de achttiende eeuw. Het is precies deze gedachtewereld waarin de \textit{Geisteswissenschaften} van de negentiende eeuw leefden, ook al konden zij haar nog niet goed epistemologisch verantwoorden.


Ik heb je vertaling herschreven met eenvoudiger zinsbouw, duidelijkere overgangen en behoud van de filosofische kernbegrippen zoals \textit{Bildung}, \textit{Geist}, tact, herinnering en historisch bewustzijn. Ik baseer me op de tekst die je hebt aangeleverd. 

\subsection{De leidende concepten van het humanisme}

\subsubsection{\textit{Bildung} (Vorming)}

Het begrip \textit{Bildung} is een goed vertrekpunt om de humanistische traditie te begrijpen. Het wijst op een diepe verandering in het denken van de achttiende en negentiende eeuw. Die verandering is zo belangrijk dat de tijd van Goethe ons nog steeds betrekkelijk nabij lijkt, terwijl de barok veel verder van ons af staat.

Veel begrippen die wij nu vanzelfsprekend gebruiken, hebben in deze periode hun moderne betekenis gekregen. Denk aan woorden als `kunst'', `geschiedenis'', `scheppend vermogen'', `wereldbeeld'', `ervaring'', `genie'', `buitenwereld'', `innerlijkheid'', `uitdrukking'', `stijl'' en ``symbool''. Zulke woorden lijken neutraal, maar ze dragen een lange geschiedenis met zich mee. Daarom moeten we hun ontwikkeling onderzoeken. Alleen zo kunnen we onszelf historisch begrijpen, in plaats van ongemerkt door de taal te worden meegesleept.

Het begrip \textit{Bildung} biedt daarvoor een bijzonder geschikt uitgangspunt. De geschiedenis van dit woord is al eerder onderzocht. Het voert terug tot de middeleeuwse mystiek, loopt via de barok en krijgt een religieuze verdieping in Klopstocks \textit{Messias}. Uiteindelijk krijgt het bij Herder de betekenis van een opstijgen tot menselijkheid door cultuur. In de negentiende eeuw blijft deze diepere betekenis aanwezig. Ook onze huidige manier van denken over vorming, cultuur en menselijkheid is daardoor nog steeds mede bepaald.

Oorspronkelijk kon \textit{Bildung} ook een natuurlijke vorm betekenen. Men kon bijvoorbeeld spreken over de vorm van ledematen of over een bergformatie. In de achttiende en negentiende eeuw raakt deze betekenis echter op de achtergrond. \textit{Bildung} gaat dan steeds meer betekenen: de specifiek menselijke ontwikkeling van aanleg, vermogens en karakter. Het begrip wordt dus verbonden met cultuur, opvoeding en zelfvorming.

Tussen Kant en Hegel krijgt Herders opvatting van \textit{Bildung} verder gestalte. Kant gebruikt het woord \textit{Bildung} in deze betekenis nog niet vaak. Hij spreekt eerder over het cultiveren van een vermogen, als een daad van vrijheid door het handelende subject. Hegel daarentegen spreekt al duidelijk over \textit{Sichbilden}, zichzelf vormen, en over \textit{Bildung}. Daarbij sluit hij aan bij Kants gedachte dat de mens ook plichten tegenover zichzelf heeft.

Wilhelm von Humboldt maakte een belangrijk onderscheid tussen \textit{Kultur} en \textit{Bildung}. Bij \textit{Kultur} gaat het vooral om de ontwikkeling van vermogens. Bij \textit{Bildung} gaat het om iets hogers en innerlijkers. Humboldt schrijft dat wij met \textit{Bildung} doelen op een geestesgesteldheid die voortkomt uit kennis en gevoel voor de totale intellectuele en morele inspanning, en die vervolgens doorwerkt in gevoeligheid en karakter.

\textit{Bildung} betekent dus meer dan alleen cultuur of ontwikkeling. Het woord herinnert ook aan de oude mystieke traditie. In die traditie draagt de mens in zijn ziel het goddelijke beeld waarnaar hij gevormd is. Hij moet dat beeld in zichzelf verder tot ontwikkeling brengen.

In andere talen werd hiervoor vaak het Latijnse \textit{formatio} gebruikt, of woorden als `form'' en `formation''. Toch is het niet toevallig dat in het Duits juist \textit{Bildung} de overhand kreeg. In \textit{Bildung} klinkt namelijk het woord \textit{Bild} mee. Dat kan zowel `afbeelding'' als `voorbeeld'' betekenen: zowel \textit{Nachbild} als \textit{Vorbild}. Die dubbele betekenis ontbreekt in het woord ``vorm''.

Net als het moderne woord ``vorming'' verwijst \textit{Bildung} niet alleen naar een proces, maar ook naar het resultaat daarvan. Toch is dat resultaat geen eindproduct dat technisch wordt gemaakt. Het ontstaat van binnenuit. Wie gevormd wordt, blijft bovendien steeds in beweging. \textit{Bildung} is nooit helemaal afgesloten.

Daarom lijkt \textit{Bildung} op de natuur, of \textit{physis}. Zij heeft geen doel buiten zichzelf. Het begrip \textit{Bildungsziel}, het doel van vorming, moet daarom voorzichtig worden gebruikt. \textit{Bildung} kan eigenlijk geen extern doel hebben. Alleen vanuit het perspectief van een opvoeder kan men spreken over een vormingsdoel.

Omdat \textit{Bildung} geen extern doel nastreeft, is zij meer dan het ontwikkelen van bestaande talenten. Een talent oefenen betekent dat men iets wat al gegeven is verder ontwikkelt. Oefening is dan een middel tot een doel. Zo is grammatica leren een middel om een taal beter te beheersen.

Bij \textit{Bildung} ligt dat anders. Wat iemand vormt, wordt werkelijk eigen gemaakt. Alles wat iemand opneemt, blijft bewaard en wordt onderdeel van hemzelf. In \textit{Bildung} gaat niets verloren van wat werkelijk is opgenomen. Daarom is \textit{Bildung} een door en door historisch begrip. Juist omdat zij bewaart wat zij opneemt, is zij van groot belang voor het begrijpen van de humanistische wetenschappen.

Zelfs een korte blik op de geschiedenis van het woord \textit{Bildung} brengt ons bij de ideeënwereld van Hegel. Hegel was de eerste die deze thematiek tot het terrein van de ``eerste filosofie'' rekende. Hij heeft de aard van \textit{Bildung} scherp beschreven. Volgens hem heeft de filosofie, en dat geldt ook voor de \textit{Geisteswissenschaften}, haar bestaansvoorwaarde in \textit{Bildung}. Het wezen van \textit{Geist}, van geest, is namelijk nauw verbonden met het idee van vorming.

De mens is niet zomaar een natuurlijk wezen. Hij wordt gekenmerkt door een breuk met het onmiddellijke en natuurlijke. Zijn intellectuele en rationele aard vraagt van hem dat hij boven zijn directe natuur uitstijgt. Daarom zegt Hegel dat de mens in deze sfeer van nature niet is wat hij zou moeten zijn. Juist daarom heeft hij \textit{Bildung} nodig.

Voor Hegel bestaat de formele aard van \textit{Bildung} in haar gerichtheid op het universele. Vorming betekent dat de mens zich verheft boven het louter bijzondere. Dat geldt niet alleen voor theoretische vorming, maar ook voor praktische vorming. Het gaat om de menselijke rationaliteit als geheel. De mens wordt pas werkelijk geestelijk wanneer hij zich niet opsluit in zijn eigen particuliere belangen, maar zich leert verstaan als een universeel geestelijk wezen.

Wie volledig opgaat in zijn eigen bijzonderheid, is volgens Hegel \textit{ungebildet}, ongevormd. Denk bijvoorbeeld aan iemand die zich overgeeft aan blinde woede, zonder maat of gevoel voor verhouding. Zo iemand mist het vermogen tot abstractie. Hij kan zijn eigen situatie niet bekijken vanuit een ruimer gezichtspunt.

\textit{Bildung} als opstijgen naar het universele is daarom een menselijke opdracht. Zij vraagt dat men zijn eigen bijzonderheid leert relativeren. Negatief gezegd betekent dit dat men zijn begeerten leert beheersen. Men wordt vrij van het object van begeerte en daardoor vrij voor de objectiviteit ervan. Men leert de dingen niet alleen te zien vanuit zijn eigen behoefte, maar zoals zij op zichzelf kunnen verschijnen.

In de \textit{Fenomenologie van de Geest} werkt Hegel dit verder uit aan de hand van arbeid. Arbeid is voor hem niet in de eerste plaats consumeren, maar vormen. Door te werken vormt de mens een ding. Maar in dat vormen van het ding vormt hij ook zichzelf. Arbeid is beteugelde begeerte. De mens stelt de onmiddellijke bevrediging uit en richt zich op iets dat duurzamer en algemener is.

In het gevormde object herkent het werkende bewustzijn zichzelf terug. Door iets te maken, verwerft de mens een vaardigheid, een kwaliteit en een besef van zichzelf. Wat aanvankelijk vreemd of opgelegd leek, wordt door arbeid eigen. Daarom kan men zeggen dat arbeid vormt. Praktische \textit{Bildung} bestaat uit afstand nemen van directe begeerte, persoonlijke behoefte en privébelang, en uit openstaan voor de eis van het universele.

In zijn \textit{Propaedeutik} geeft Hegel voorbeelden van praktische \textit{Bildung}. Zij blijkt bijvoorbeeld uit matiging: men beperkt de bevrediging van behoeften door rekening te houden met een algemener gezichtspunt, zoals gezondheid. Zij blijkt ook uit voorzichtigheid: men kijkt niet alleen naar de onmiddellijke situatie, maar houdt rekening met wat verder nodig kan zijn.

Ook de keuze voor een beroep bevat volgens Hegel een element van praktische \textit{Bildung}. Een beroep heeft altijd iets van een lot of een uiterlijke noodzaak. Het vraagt dat men zich wijdt aan taken die men niet zomaar uit persoonlijk verlangen zou kiezen. Praktische vorming bestaat erin dat men zijn beroep volledig leert vervullen. Men overwint het vreemde element ervan en maakt het tot iets eigens. Zich wijden aan de algemeenheid van een beroep betekent dat men zich leert beperken. Maar wanneer het beroep werkelijk eigen wordt, is die beperking niet langer alleen een beperking.

In deze beschrijving van praktische \textit{Bildung} wordt de grondtrek van de historische geest zichtbaar: de mens vindt zichzelf terug in wat eerst vreemd was. Hij verzoent zich met zichzelf door zich in het andere te herkennen.

Dat wordt nog duidelijker bij theoretische \textit{Bildung}. Een theoretische houding betekent al dat men afstand neemt van het onmiddellijke. Men houdt zich bezig met iets dat niet direct vertrouwd is, iets dat vreemd is, iets dat tot herinnering en denken behoort. Theoretische vorming voert ons voorbij wat wij onmiddellijk weten en ervaren.

Theoretische \textit{Bildung} betekent dat men leert erkennen wat anders is dan men zelf is. Men leert universele gezichtspunten innemen. Daardoor kan men het objectieve in zijn vrijheid begrijpen, zonder het meteen te beoordelen vanuit de eigen belangen of gewoonten. Daarom gaat vorming altijd samen met de ontwikkeling van theoretische belangstelling.

Hegel zag de wereld en taal van de klassieke oudheid als bijzonder geschikt voor deze vorming. Die wereld is vreemd genoeg om ons van onszelf los te maken, maar tegelijk bevat zij volgens hem de wegen waardoor wij tot onszelf kunnen terugkeren. Door ons in de oudheid te verdiepen, raken wij vertrouwd met iets vreemds en vinden wij onszelf op een meer universeel niveau terug.

In deze gedachte herkennen we Hegel als directeur van het \textit{Gymnasium}. Hij deelt het klassieke vooroordeel dat juist in de wereld van de oudheid het universele wezen van de geest het duidelijkst te vinden is. Toch is zijn grondgedachte belangrijk. De fundamentele beweging van de geest is dat hij het eigene herkent in het vreemde. De geest komt tot zichzelf door eerst van zichzelf weg te gaan en via het andere terug te keren.

Daarom is alle theoretische \textit{Bildung}, zelfs het leren van vreemde talen en vreemde begrippenwerelden, een voortzetting van een proces dat al veel eerder begint. Ieder mens die boven zijn natuurlijke bestaan uitstijgt, treft in de taal, gewoonten en instellingen van zijn volk een wereld aan die al gevormd is. Net zoals een kind leert spreken, moet hij zich die wereld eigen maken.

Zo is ieder individu altijd betrokken in een proces van \textit{Bildung}. Hij stijgt boven zichzelf uit doordat hij opgroeit in een wereld die al door taal en gewoonte menselijk gevormd is. Hegel benadrukt dat een volk zichzelf bestaan geeft in zijn wereld. Het werkt uit wat in zichzelf ligt en brengt zo zijn wezen naar buiten.

De kern van \textit{Bildung} is dus niet vervreemding op zichzelf. Het gaat om de terugkeer naar zichzelf, maar die terugkeer loopt via het vreemde. Toch moeten we \textit{Bildung} niet alleen zien als een historisch proces waarin de geest zich verheft tot het universele. \textit{Bildung} is ook het element waarin de gevormde mens zich beweegt.

Daarmee komen de vragen terug die eerder bij Helmholtz opkwamen. Wat is dat element waarin de gevormde mens zich beweegt? Hegels antwoord kan ons niet helemaal bevredigen. Hij ziet \textit{Bildung} uiteindelijk als voltooid in absolute kennis, waarin de geest alle vervreemding heeft opgeheven en de werkelijkheid volledig heeft begrepen. Dat past bij zijn filosofie van de absolute geest.

Toch hoeven wij Hegels systeem niet over te nemen om zijn inzicht te behouden. We kunnen erkennen dat \textit{Bildung} een element van de geest is, zonder Hegels leer van de absolute geest te aanvaarden. Net zo kunnen we de historiciteit van het bewustzijn erkennen zonder Hegels filosofie van de wereldgeschiedenis te volgen.

Het ideaal van volledige \textit{Bildung} blijft belangrijk, ook voor historische wetenschappen die zich van Hegel distantiëren. Want \textit{Bildung} is het element waarin deze wetenschappen zich bewegen. Zelfs wanneer oudere schrijvers over de uiterlijke ``volmaaktheid van vorm'' spraken, bedoelden zij niet alleen het eindpunt van een ontwikkeling. Zij bedoelden een volwassen toestand waarin alle onderdelen vrij en harmonisch kunnen bewegen.

Zo veronderstellen de humanistische wetenschappen dat het geleerde bewustzijn al gevormd is. Zij veronderstellen een bepaald tactgevoel: een juist oordeel dat men niet eenvoudig kan aanleren of nadoen. Dit tactgevoel omgeeft de manier van oordelen en kennen in de humanistische wetenschappen. Het is als het element waarin zij zich bewegen.

Helmholtz beschreef het functioneren van de humanistische wetenschappen met woorden als `kunstzinnig gevoel'' en `tact''. Daarmee wees hij inderdaad op dit element van \textit{Bildung}, waarin de geest een bijzondere bewegelijkheid bezit. Hij sprak bijvoorbeeld over het gemak waarmee zeer uiteenlopende ervaringen in het geheugen van de historicus of filoloog moeten kunnen opkomen.

Op het eerste gezicht lijkt dit een beschrijving van buitenaf. Helmholtz blijft immers denken vanuit het ideaal van de natuurwetenschappen: het bewuste trekken van strenge conclusies. Maar zijn begrip van geheugen is te beperkt. Het verklaart niet werkelijk wat er in de humanistische wetenschappen gebeurt.

Tact is niet zomaar een extra psychische bekwaamheid, alsof iemand met een sterk geheugen daardoor tot inzichten komt die niet streng controleerbaar zijn. Wat tact mogelijk maakt, is geen louter psychologische aanleg. Het is iets dat gevormd wordt. Het behoort tot de manier waarop iemand zich in een historische en culturele wereld beweegt.

Ook geheugen zelf moet anders worden begrepen. Het is niet simpelweg een algemeen vermogen waarmee men informatie opslaat. Onthouden, vergeten en herinneren horen bij het historische bestaan van de mens. Zij maken deel uit van zijn geschiedenis en van zijn \textit{Bildung}.

Wie zijn geheugen alleen gebruikt als techniek, bezit het nog niet werkelijk als iets eigens. Geheugen moet gevormd worden. Men herinnert zich niet alles op dezelfde manier. Men heeft geheugen voor bepaalde dingen en niet voor andere. Men wil sommige dingen bewaren en andere juist verdringen.

Daarom moeten we geheugen bevrijden uit de beperkte opvatting dat het slechts een psychologisch vermogen is. Het is een wezenlijk element van het eindige en historische bestaan van de mens. Ook vergeten hoort daarbij. Vergeten is niet alleen een gebrek of een afwezigheid. Zoals vooral Nietzsche heeft benadrukt, is vergeten ook een voorwaarde voor het leven van de geest. Alleen doordat wij kunnen vergeten, kan de geest zich vernieuwen. Daardoor kunnen wij dingen opnieuw en met frisse ogen zien, zodat het bekende zich met het nieuwe kan verbinden.

Hetzelfde geldt voor het begrip ``tact''. In het gewone taalgebruik bedoelen we met tact een bijzondere gevoeligheid voor situaties en voor de manier waarop men zich daarin moet gedragen. Algemene regels zijn daarvoor niet genoeg. Tact is juist nodig wanneer men niet eenvoudig uit principes kan afleiden wat men moet doen.

Een wezenlijk kenmerk van tact is dat zij zich niet volledig laat uitspreken. Tact bestaat vaak juist in wat men niet zegt. Men kan iets tactvol zeggen, maar dat betekent meestal ook dat men iets anders tactvol voorbijgaat. Het is tactloos om uit te spreken wat beter onuitgesproken kan blijven.

Iets voorbijgaan betekent echter niet dat men het niet ziet. Integendeel: men houdt het juist in het oog, maar zonder het hard of kwetsend aan te raken. Tact bewaart afstand. Zij vermijdt het opdringerige, het kwetsende en het binnendringen in de persoonlijke sfeer van de ander.

De tact waarover Helmholtz spreekt, is niet helemaal hetzelfde als tact in de omgang met mensen. Toch delen beide vormen iets wezenlijks. De tact die in de humanistische wetenschappen werkt, is niet zomaar een gevoel en ook niet alleen onbewust. Zij is tegelijk een manier van kennen en een manier van zijn.

Dat wordt duidelijk vanuit het begrip \textit{Bildung}. Wat Helmholtz tact noemt, veronderstelt vorming. Het is verbonden met zowel esthetische als historische \textit{Bildung}. Wie in de humanistische wetenschappen op zijn tact wil kunnen vertrouwen, moet een gevoel voor het esthetische en het historische hebben ontwikkeld.

Dat gevoel is niet simpelweg een natuurlijke aanleg. Daarom spreken we terecht over esthetisch bewustzijn en historisch bewustzijn, niet alleen over gevoel. Toch heeft dit bewustzijn iets gemeen met de onmiddellijkheid van de zintuigen. Het kan in een concreet geval juiste onderscheidingen maken, zonder altijd precies te kunnen uitleggen waarom.

Iemand met esthetisch gevoel kan bijvoorbeeld onderscheid maken tussen mooi en lelijk, verheven en laag. Iemand met historisch gevoel weet wat in een bepaalde tijd mogelijk was en wat niet. Hij heeft gevoel voor de andersheid van het verleden ten opzichte van het heden.

Als dit alles \textit{Bildung} veronderstelt, dan gaat het niet om een methode of procedure. Het gaat om iets dat in iemand tot stand is gekomen. Men kan niet eenvoudig zeggen: observeer nauwkeuriger, bestudeer de traditie grondiger, en dan ontstaat vanzelf historisch of esthetisch begrip. Daarvoor moet er al ontvankelijkheid zijn voor de andersheid van het kunstwerk of van het verleden.

Dat is wat wij, in aansluiting bij Hegel, als de algemene eigenschap van \textit{Bildung} kunnen beschrijven: openstaan voor wat anders is. De gevormde mens kan zich verplaatsen in andere, meer universele gezichtspunten. Hij heeft gevoel voor verhouding en afstand ten opzichte van zichzelf. Daarom bestaat \textit{Bildung} in een opstijgen boven zichzelf naar universaliteit.

Afstand nemen van zichzelf en van de eigen particuliere doelen betekent dat men ze leert zien zoals anderen ze zouden kunnen zien. Deze universaliteit is echter niet de abstracte universaliteit van een begrip of een regel. Het gaat er niet om dat een bijzonder geval onder een algemene regel wordt gebracht. Er wordt ook niets definitief bewezen.

De universele gezichtspunten waarvoor de gevormde mens openstaat, zijn geen vaste maatstaven die men eenvoudig kan toepassen. Zij zijn eerder gezichtspunten van mogelijke anderen. Daarom heeft gevormd bewustzijn inderdaad iets van een gevoel. Een zintuig, zoals het zien, is ook op een bepaalde manier universeel: het opent een veld waarin onderscheidingen mogelijk worden. Zo opent \textit{Bildung} een veel ruimer veld. Zij is geen afzonderlijk zintuig, maar een universeel gevoel.

Deze gedachte van een universeel en gemeenschappelijk gevoel brengt ons terug naar een brede historische context. Als we nadenken over \textit{Bildung}, zoals die ook op de achtergrond staat bij Helmholtz, komen we in een lange traditie terecht. Die traditie moeten we onderzoeken als we het filosofische probleem van de humanistische wetenschappen willen begrijpen.

Het probleem is namelijk dat de negentiende-eeuwse methodologie de humanistische wetenschappen te nauw heeft opgevat. Het moderne begrip van wetenschap en methode is onvoldoende om te begrijpen wat deze wetenschappen tot wetenschap maakt. Dat wordt duidelijker wanneer we uitgaan van de humanistische traditie van \textit{Bildung}. Juist in haar verzet tegen de aanspraken van de moderne wetenschap krijgt deze traditie opnieuw betekenis.

Het zou een afzonderlijk onderzoek waard zijn om te volgen hoe, vanaf het humanisme, kritiek op de scholastieke wetenschap is ontstaan en telkens is veranderd naargelang haar tegenstander. Oorspronkelijk werden hierbij klassieke motieven opnieuw tot leven gewekt. De humanisten prezen de Griekse taal en de weg van de \textit{eruditio}, de geleerdheid. Dat was meer dan antiquarische belangstelling voor het verleden.

De herleving van de klassieke talen bracht ook een nieuwe waardering van de retorica met zich mee. Zij verzette zich tegen de ``school'', dat wil zeggen tegen de scholastieke wetenschap. Daartegenover stelde zij een ideaal van menselijke wijsheid dat niet in de schoolse methode werd bereikt. Deze tegenstelling is eigenlijk al aanwezig aan het begin van de filosofie.

Plato's kritiek op de sofisten, en vooral zijn dubbelzinnige houding tegenover Isocrates, wijzen op het filosofische probleem dat hier opkomt. Toen in de zeventiende eeuw het nieuwe methodische bewustzijn van de moderne wetenschap ontstond, werd dit oude probleem opnieuw urgent. De nieuwe wetenschap maakte immers exclusieve aanspraken op waarheid. Daardoor werd de vraag dringender of het humanistische begrip \textit{Bildung} misschien een eigen bron van waarheid bevatte.

We zullen zien dat de humanistische wetenschappen van de negentiende eeuw, zonder het zelf duidelijk te erkennen, leefden uit het voortbestaan van dit humanistische ideaal van \textit{Bildung}.

Daarbij is duidelijk dat hier niet de wiskunde centraal staat, maar de humanistische studies. Want wat kon het nieuwe methodische ideaal van de zeventiende eeuw eigenlijk betekenen voor de humanistische wetenschappen? Wie de relevante hoofdstukken leest van de \textit{Logique de Port-Royal}, over de regels van de rede toegepast op historische waarheden, ziet hoe weinig dit methode-ideaal daar oplevert.

De resultaten zijn vaak tamelijk triviaal. Men zegt bijvoorbeeld dat men, om de waarheid van een gebeurtenis te beoordelen, rekening moet houden met de omstandigheden, de \textit{circonstances}. Met zulke argumenten probeerden de Jansenisten een methode te geven om te bepalen welke wonderen geloofwaardig waren. Zij bestreden een kritiekloos geloof in wonderen met de geest van de nieuwe methode. Tegelijk wilden zij daarmee de ware wonderen van de Bijbel en de kerkelijke traditie veiligstellen.

Hier zien we de nieuwe wetenschap in dienst van de oude kerk. Dat deze verhouding niet stabiel kon blijven, is duidelijk. Zodra de christelijke vooronderstellingen zelf ter discussie kwamen te staan, moest de methode tot andere resultaten leiden. Wanneer het natuurwetenschappelijke methode-ideaal werd toegepast op de geloofwaardigheid van de historische getuigenissen van de Schrift, had dat onvermijdelijk ingrijpende gevolgen voor het christendom.

De stap van de Jansenistische kritiek op wonderen naar de historische kritiek op de Bijbel is klein. Spinoza is daarvan een duidelijk voorbeeld. Later zal blijken dat een volledig consequente toepassing van deze methode als enige maatstaf voor waarheid in de humanistische wetenschappen zou leiden tot hun zelfvernietiging.

Hier is een geredigeerde versie van je vertaling. Ik heb de zinnen minder zwaar gemaakt, enkele begrippen consistenter vertaald, en de argumentatie duidelijker geordend zonder de inhoud te verkorten of te versimpelen. Ik baseer me op de door jou aangeleverde tekst. 

\subsubsection{\textit{Sensus communis}}

De humanistische traditie biedt een belangrijke sleutel om de eigen kenniswijze van de menswetenschappen te begrijpen. Een goed uitgangspunt daarvoor is Giambattista Vico's \textit{De nostri temporis studiorum ratione}. Zijn verdediging van het humanisme is geïnspireerd door de jezuïtische pedagogische traditie en richt zich zowel tegen het jansenisme als tegen de cartesiaanse filosofie.

Net als in zijn ontwerp van een ``nieuwe wetenschap'' beroept Vico zich in dit pedagogische manifest op oude waarheden. Daarbij staan vooral twee begrippen centraal: de \textit{sensus communis}, het gemeenschappelijk zintuig of gezonde verstand, en het humanistische ideaal van \textit{eloquentia}, de welsprekendheid. Beide begrippen komen voort uit de klassieke opvatting van wijsheid. In die traditie betekende goed spreken, \textit{eu legein}, niet alleen mooi of overtuigend spreken, maar ook: het juiste zeggen, de waarheid spreken. Het ging dus niet alleen om stijl, maar ook om waarheid en inhoud.

In de oudheid werd dit ideaal zowel door filosofen als door leraren in de retorica verdedigd. De retorica stond vaak tegenover de filosofie, maar zij maakte tegelijk aanspraak op een eigen vorm van wijsheid. Anders dan de speculatieve sofisten wilde zij echte praktische wijsheid onderwijzen. Vico, zelf hoogleraar retorica, plaatst zich nadrukkelijk in deze humanistische traditie, die teruggaat tot de klassieke oudheid.

Deze traditie is van groot belang voor het zelfbegrip van de menswetenschappen. Zij laat namelijk zien dat het retorische ideaal een positieve dubbelzinnigheid bezit. Retorica is niet louter versiering of overreding, maar heeft ook betrekking op waarheid, oordeel en praktische wijsheid. Juist dit element werd niet alleen door Plato bekritiseerd, maar later ook door de anti-retorische methodologie van de moderne tijd naar de achtergrond gedrongen.

Naast het retorische element beroept Vico zich op de \textit{sensus communis}. Ook dit begrip is diep geworteld in de klassieke traditie. Daarbij speelt het onderscheid tussen de geleerde en de wijze mens een centrale rol. Dit onderscheid werd voor het eerst scherp zichtbaar bij de Cynici, in hun opvatting van Socrates. Het berust op het verschil tussen \textit{sophia}, theoretische wijsheid, en \textit{phronesis}, praktische wijsheid.

Aristoteles werkte dit onderscheid verder uit als kritiek op een uitsluitend theoretisch levensideaal. In de hellenistische periode werd het vervolgens belangrijk voor het beeld van de wijze mens. Dat gebeurde vooral toen het Griekse ideaal van \textit{Bildung} samenging met het Romeinse zelfbewustzijn van de heersende klasse. Ook de latere Romeinse rechtswetenschap ontwikkelde zich tegen de achtergrond van een rechtspraktijk die dichter bij het praktische ideaal van \textit{phronesis} stond dan bij het theoretische ideaal van \textit{sophia}.

Met de herleving van de klassieke filosofie en retorica kreeg het beeld van Socrates opnieuw betekenis. Het werd een symbool van verzet tegen een eenzijdig begrip van wetenschap. Dat blijkt onder meer uit de figuur van de \textit{idiota}, de leek, die een nieuwe positie innam tussen de geleerde en de wijze mens.

Vico bekritiseert de stoïcijnen omdat zij de rede beschouwden als de \textit{regula veri}, de maatstaf van de waarheid. Daartegenover prijst hij de Oude Academici, die alleen beweerden te weten dat zij niets wisten, en de Nieuwe Academici, die uitblonken in de kunst van het argumenteren. Ook die argumentatiekunst behoort tot het terrein van de retorica.

Vico's beroep op de \textit{sensus communis} krijgt binnen deze humanistische traditie een eigen kleur. Hij erkent de verdiensten van de moderne kritische wetenschap, maar wijst tegelijk op haar grenzen. Ook met de nieuwe wetenschap en haar wiskundige methoden kunnen wij volgens hem niet zonder de wijsheid van de oudheid. Die wijsheid bestaat vooral in de vorming van \textit{prudentia}, praktische voorzichtigheid, en \textit{eloquentia}, welsprekendheid.

Voor de opvoeding is volgens Vico echter vooral de oefening van de \textit{sensus communis} van belang. Die wordt niet gevoed door absolute waarheden, maar door het waarschijnlijke en het waarachtig lijkende: het \textit{verisimile}. Juist in dit domein van het waarschijnlijke beweegt zich veel van het menselijke leven.

Voor Vico is de \textit{sensus communis} niet slechts een algemeen menselijk vermogen. Het is vooral het gevoel voor wat recht en goed is. Dit gevoel is in alle mensen aanwezig, maar wordt gevormd door het leven in gemeenschap. Het begrip doet denken aan het natuurrecht en aan de stoïcijnse \textit{koinai ennoiai}, de gemeenschappelijke begrippen. Toch grijpt Vico vooral terug op het Romeinse begrip van \textit{sensus communis}, zoals we dat bij de Romeinse klassieken aantreffen.

De Romeinen zetten zich af tegen de Griekse cultuur en hechtten grote waarde aan hun eigen tradities van burgerlijk en sociaal leven. In hun begrip van \textit{sensus communis} klinkt daarom een kritische toon mee tegenover de theoretische speculaties van de filosofen. Vico neemt die toon opnieuw op, maar nu in zijn verzet tegen de moderne wetenschap, de \textit{critica}.

Het ligt voor de hand dat filologische en historische studies, en daarmee de werkwijze van de menswetenschappen, op dit begrip van \textit{sensus communis} kunnen worden gebaseerd. Hun onderwerp is immers het morele en historische bestaan van de mensheid, zoals dat zichtbaar wordt in woorden en daden. Dat bestaan wordt zelf wezenlijk bepaald door de \textit{sensus communis}. Een redenering die alleen uitgaat van algemene principes of van formele bewijsvoering is hier niet voldoende, omdat de concrete omstandigheden beslissend zijn.

Dit is geen louter negatieve vaststelling. De \textit{sensus communis} levert een eigen positieve vorm van kennis op. De kenniswijze van de geschiedenis is niet afdoende beschreven wanneer men zegt dat hier, zoals Tetens deed, `geloof in het getuigenis van anderen'' in de plaats komt van de `zelfbewuste deductie'' waarover Helmholtz sprak. Ook is het niet zo dat historische kennis daarom minder waarheidswaarde zou hebben.

D'Alembert heeft gelijk wanneer hij schrijft dat waarschijnlijkheid vooral een rol speelt bij historische feiten, en in het algemeen bij alle gebeurtenissen uit verleden, heden en toekomst die wij aan het toeval toeschrijven omdat wij hun oorzaken niet volledig kennen. Het deel van deze kennis dat betrekking heeft op heden en verleden, kan, ook al berust het slechts op getuigenissen, vaak een overtuiging teweegbrengen die even sterk is als de overtuiging die uit axioma's voortkomt.

Geschiedenis is dus een bron van waarheid die fundamenteel verschilt van de theoretische rede. Dat bedoelde Cicero toen hij de geschiedenis \textit{vita memoriae} noemde: het leven van het geheugen. Geschiedenis heeft een eigen recht, omdat menselijke hartstochten en handelingen niet eenvoudig door algemene voorschriften van de rede kunnen worden bestuurd. In deze sfeer hebben wij juist overtuigende voorbeelden nodig, en alleen de geschiedenis kan die leveren. Daarom noemt Bacon de \textit{historia}, die zulke voorbeelden biedt, bijna een andere manier van filosoferen: \textit{alia ratio philosophandi}.

Vico's terugkeer naar het Romeinse begrip van \textit{sensus communis} en zijn verdediging van de humanistische retorica tegenover de moderne wetenschap zijn daarom voor Gadamer bijzonder belangrijk. Hier wordt een element van waarheid zichtbaar dat in de negentiende eeuw nauwelijks nog herkenbaar was. De menswetenschappen probeerden zichzelf toen immers te begrijpen aan de hand van de methoden van de moderne wetenschap.

Vico leefde nog in een ononderbroken traditie van retorische en humanistische cultuur. Daarom hoefde hij alleen de blijvende aanspraak van die traditie opnieuw te bevestigen. In die traditie was men zich er altijd van bewust geweest dat rationele bewijsvoering en methodisch onderwijs niet het hele domein van kennis kunnen omvatten.

Wij daarentegen moeten ons moeizaam een weg terug banen naar deze traditie. Dat kan alleen door eerst te laten zien welke moeilijkheden ontstaan wanneer men het moderne methodebegrip toepast op de menswetenschappen. Daarom is het belangrijk te onderzoeken hoe deze traditie verarmde, en hoe de waarheidsaanspraak van de menswetenschappen gemeten ging worden aan een maatstaf die hun vreemd is: het methodische denken van de moderne wetenschap.

Hoewel Vico en de ononderbroken retorische traditie van Italië deze ontwikkeling niet rechtstreeks beïnvloedden — zij werd vooral bepaald door de Duitse historische school — was Vico niet de enige die zich op de \textit{sensus communis} beriep. Een belangrijke parallel vinden we bij Shaftesbury, die in de achttiende eeuw grote invloed had.

Shaftesbury bracht de beoordeling van de sociale betekenis van geestigheid en humor onder bij de \textit{sensus communis}. Daarbij beriep hij zich uitdrukkelijk op de Romeinse klassieken en hun humanistische uitleggers. Het begrip \textit{sensus communis} doet onmiskenbaar denken aan de stoïcijnen en aan het natuurrecht, maar bij Shaftesbury is ook de humanistische interpretatie vanuit de Romeinse klassieken duidelijk aanwezig.

Volgens Shaftesbury verstonden de humanisten onder \textit{sensus communis} niet alleen een gevoel voor het algemeen welzijn, maar ook liefde voor gemeenschap en samenleving, natuurlijke genegenheid, menselijkheid en behulpzaamheid. Zij namen daarbij een woord van Marcus Aurelius over: \textit{koinonoemosune}. Dat is een zeldzaam en kunstmatig woord. Juist daardoor blijkt dat het begrip niet rechtstreeks uit de Griekse filosofie afkomstig is, al klinkt de stoïcijnse opvatting er wel in mee.

De humanist Salmasius omschreef de inhoud van dit woord als een beheerste, gewone en regelmatige denkwijze van een mens die als het ware uitgaat van de gemeenschap. Zo iemand brengt niet alles terug tot zijn eigen voordeel, maar richt zijn aandacht op de zaak waarmee hij te maken heeft en denkt met maat en beheersing over zichzelf.

Voor Shaftesbury is de \textit{sensus communis} dus niet in de eerste plaats een vermogen dat alle mensen van nature bezitten, en ook niet slechts een onderdeel van het natuurrecht. Zij is vooral een sociale deugd: eerder een deugd van het hart dan van het verstand. Wanneer hij geestigheid en humor in verband brengt met de \textit{sensus communis}, sluit hij aan bij oude Romeinse opvattingen over \textit{humanitas}. Daarin hoort ook een verfijnde \textit{savoir vivre}: het vermogen om een grap te begrijpen en te maken, juist omdat men een diepere verbondenheid met de gesprekspartner ervaart.

Hoewel de \textit{sensus communis} hier vooral verschijnt als een deugd van het sociale verkeer, ligt er toch een morele, zelfs metafysische grondslag onder. Shaftesbury denkt daarbij aan de intellectuele en sociale deugd van sympathie. Op die sympathie baseert hij niet alleen de moraal, maar zelfs een hele esthetische metafysica. Zijn opvolgers, vooral Hutcheson en Hume, werken deze gedachte verder uit tot de leer van de \textit{moral sense}. Die leer zou later een tegenhanger vormen van de Kantiaanse ethiek.

In de Schotse filosofie kreeg het begrip ``common sense'' een centrale systematische functie. Deze filosofie verzette zich tegen de metafysica en tegen de sceptische ontbinding daarvan. Zij bouwde een nieuw systeem op vanuit oorspronkelijke en natuurlijke oordelen van het \textit{common sense}, zoals bij Thomas Reid.

Ongetwijfeld speelde hierbij de aristotelische en scholastieke traditie van het begrip \textit{sensus communis} een rol. Het onderzoek naar de zintuigen en hun kenvermogens stamt uit deze traditie. Het was uiteindelijk bedoeld om de overdrijvingen van filosofische speculatie te corrigeren. Toch bleef ook hier de verbinding tussen \textit{common sense} en samenleving behouden. Het gezonde verstand moest ons leiden in de alledaagse zaken van het leven, waar ons redeneervermogen ons in het duister zou laten.

Voor deze denkers was de filosofie van het gezonde verstand, van de \textit{good sense}, dus niet alleen een remedie tegen de ``maanziekte'' van de metafysica. Zij bevatte ook de grondslag van een morele filosofie die recht deed aan het sociale leven.

Het morele element in het begrip \textit{common sense}, of in het Franse \textit{le bon sens}, is tot op vandaag bewaard gebleven. Daardoor verschillen deze begrippen van het Duitse \textit{der gesunde Menschenverstand}, dat eerder een theoretische bijklank heeft gekregen.

Een goed voorbeeld hiervan is Henri Bergsons rede over \textit{le bon sens}, uitgesproken bij de prijsuitreiking aan de Sorbonne in 1895. Bergson bekritiseert daarin de abstracties van de natuurwetenschap, de taal en het juridische denken. Hij doet een hartstochtelijk beroep op de innerlijke energie van een intelligentie die zichzelf telkens opnieuw terugwint, doordat zij gevormde ideeën loslaat ten gunste van ideeën die nog in wording zijn.

In Frankrijk werd dit \textit{le bon sens} genoemd. Natuurlijk bevat het begrip een verwijzing naar de zintuigen. Maar voor Bergson spreekt het vanzelf dat \textit{le bon sens}, anders dan de gewone zintuigen, betrekking heeft op het \textit{milieu social}. Terwijl de andere zintuigen ons verbinden met dingen, bestuurt het gezonde verstand onze relaties met personen.

\textit{Le bon sens} is bij Bergson een soort genie voor het praktische leven. Toch is het minder een aangeboren gave dan een voortdurende opdracht: zich telkens opnieuw aanpassen aan nieuwe situaties. Het gaat om het toepassen van algemene principes op de werkelijkheid, waardoor rechtvaardigheid mogelijk wordt. Bergson noemt het een tact voor praktische waarheid, een juistheid van oordeel die voortkomt uit juistheid van de ziel.

Voor Bergson is \textit{le bon sens}, als gemeenschappelijke bron van denken en willen, een \textit{sens social}. Het vermijdt zowel de fouten van wetenschappelijke dogmatici die naar sociale wetten zoeken, als die van metafysische utopisten. Misschien, zegt hij, bestaat er strikt genomen geen methode, maar eerder een bepaalde wijze van handelen.

Bergson wijst ook op het belang van klassieke studies voor de ontwikkeling van deze \textit{bon sens}. Hij ziet die studies als een poging om het ``ijs van de woorden'' te doorbreken en de vrije stroom van gedachten daaronder te ontdekken. Maar hij stelt niet de omgekeerde vraag: hoe noodzakelijk is \textit{le bon sens} voor de klassieke studies zelf? Met andere woorden, hij spreekt niet over de hermeneutische functie van \textit{le bon sens}. Zijn vraag heeft niet in de eerste plaats betrekking op de wetenschappen, maar op de zelfstandige betekenis van \textit{le bon sens} voor het leven.

Voor Gadamer is vooral belangrijk dat de morele en politieke betekenis van dit begrip voor Bergson en zijn toehoorders vanzelfsprekend was. Dat laat zien hoe levend deze traditie buiten Duitsland gebleven is.

Kenmerkend voor de zelfreflectie van de menswetenschappen in de negentiende eeuw is echter dat zij niet werd beïnvloed door deze traditie van de morele filosofie, waartoe Vico en Shaftesbury behoren. Die traditie leefde vooral voort in Frankrijk, het klassieke land van \textit{le bon sens}. De Duitse menswetenschappen lieten zich daarentegen vooral leiden door de Duitse filosofie uit de tijd van Kant en Goethe.

In Engeland en de Romaanse landen bleef het begrip \textit{sensus communis} tot op heden meer dan een kritische leus. Het bleef verbonden met een algemene burgerdeugd. In Duitsland daarentegen namen de achttiende-eeuwse volgelingen van Shaftesbury en Hutcheson het politieke en sociale element van de \textit{sensus communis} niet werkelijk over.

De schoolmetafysica en de populaire filosofie van de achttiende eeuw bestudeerden en imiteerden weliswaar de leidende landen van de Verlichting, Engeland en Frankrijk. Maar zij konden een idee niet werkelijk opnemen waarvoor de sociale en politieke voorwaarden ontbraken. Het begrip \textit{sensus communis} werd wel overgenomen, maar doordat men het van zijn politieke inhoud ontdeed, verloor het zijn kritische kracht.

In Duitsland werd \textit{sensus communis} vooral opgevat als een theoretisch vermogen: het vermogen tot theoretisch oordeel, naast moreel bewustzijn en smaak. Zo werd het opgenomen in een scholastiek van de grondvermogens. Herder bekritiseerde deze manier van denken in zijn vierde \textit{Kritisches Wäldchen}, gericht tegen Riedel. Daarmee werd hij een voorloper van het historicisme, ook op het gebied van de esthetiek.

Toch is er één belangrijke uitzondering: het piëtisme. Niet alleen voor een wereldman als Shaftesbury was het belangrijk om de aanspraken van de wetenschap, de \textit{demonstratio}, tegenover de ``school'' te begrenzen en zich te beroepen op de \textit{sensus communis}. Dat gold ook voor de prediker die het hart van zijn gemeente wilde bereiken.

Zo beriep de Zwabische piëtist Oetinger zich uitdrukkelijk op Shaftesbury's verdediging van de \textit{sensus communis}. Bij Oetinger wordt \textit{sensus communis} eenvoudig vertaald als ``hart''. Hij omschrijft het als datgene wat te maken heeft met de dingen die alle mensen dagelijks voor zich zien, de dingen die een samenleving bijeenhouden, en met waarheden en uitspraken én met de ordeningen en patronen die in zulke uitspraken besloten liggen.

Oetinger wilde aantonen dat het niet genoeg is om begrippen helder te maken. Helderheid alleen is volgens hem niet voldoende voor levende kennis. Daarvoor zijn ook zekere voorgevoelens en neigingen nodig. Vaders worden bijvoorbeeld zonder bewijs bewogen om voor hun kinderen te zorgen. Liefde bewijst niet, maar trekt het hart aan, soms zelfs tegen de rede in, ook wanneer de geliefde wordt berispt.

Oetingers beroep op de \textit{sensus communis} tegen het rationalisme van de ``school'' is voor Gadamer vooral interessant omdat hij er een uitdrukkelijk hermeneutische toepassing aan geeft. Voor Oetinger, als kerkman, staat het verstaan van de Schrift centraal. Omdat de wiskundige en demonstratieve methode daar tekortschiet, vraagt hij om een andere methode: de generatieve methode. Daarmee bedoelt hij een organische presentatie van de Schrift, zodat de gerechtigheid als een scheut kan worden geplant.

Oetinger maakte het begrip \textit{sensus communis} ook tot onderwerp van een uitvoerig en geleerd onderzoek, eveneens gericht tegen het rationalisme. Hij ziet daarin de bron van alle waarheden, de \textit{ars inveniendi} zelf, de kunst van het vinden. Daarmee keert hij zich tegen Leibniz, die alles baseert op een louter \textit{calculus metaphysicus}, met uitsluiting van elke innerlijke smaak: \textit{excluso omni gusto interno}.

Volgens Oetinger ligt de ware basis van de \textit{sensus communis} in het begrip \textit{vita}, leven: \textit{sensus communis vitae gaudens}. Tegenover het gewelddadig ontleden van de natuur door experiment en berekening stelt hij de natuurlijke ontwikkeling van het eenvoudige naar het complexe. Dat ziet hij als de universele groe wet van de goddelijke schepping, en ook van de menselijke geest.

Voor zijn gedachte dat alle kennis in de \textit{sensus communis} ontstaat, beroept Oetinger zich op Wolff, Bernoulli, Pascal, Maupertuis' onderzoek naar de oorsprong van de taal, Bacon, Fénelon en anderen. Hij definieert de \textit{sensus communis} als de levendige en doordringende waarneming van voorwerpen die voor alle mensen evident zijn, doordat zij onmiddellijk worden aangeraakt en aanschouwd, en die volstrekt eenvoudig zijn.

Uit deze definitie blijkt dat Oetinger de humanistische en politieke betekenis van \textit{sensus communis} verbindt met het peripatetische begrip ervan. Zijn omschrijving doet soms denken aan Aristoteles' leer van de \textit{nous}. Ook neemt hij de aristotelische vraag over naar de \textit{koinè dunamis}, het gemeenschappelijke vermogen dat zien, horen en de andere zintuigen verbindt.

Voor Oetinger bevestigt dit gemeenschappelijke vermogen het werkelijk goddelijke mysterie van het leven. Dat mysterie bestaat in eenvoud. De mens heeft die eenvoud door de zondeval verloren, maar kan door Gods genade terugkeren tot eenheid en eenvoud. De werking van de \textit{logos}, dat wil zeggen de aanwezigheid van God, integreert de veelheid in de eenheid.

De aanwezigheid van God bestaat volgens Oetinger juist in het leven zelf, in dit gemeenschappelijke zintuig dat alle levende wezens van dode dingen onderscheidt. Het is dan ook geen toeval dat hij voorbeelden noemt als de poliep en de zeester, die zich zelfs na in stukken te zijn gesneden kunnen herstellen en nieuwe individuen vormen.

Bij de mens werkt dezelfde goddelijke kracht als instinct en innerlijke aandrang om de sporen van God te ontdekken en te herkennen wat het meest verbonden is met menselijk geluk en leven. Oetinger maakt daarom een duidelijk onderscheid tussen rationele waarheden en ontvankelijkheid voor algemene waarheden: zintuiglijke waarheden die nuttig zijn voor alle mensen, op alle tijden en plaatsen.

Het gemeenschappelijke zintuig is volgens hem een geheel van instincten. Daarmee bedoelt hij een natuurlijke drang naar datgene waarvan het ware geluk van het leven afhangt. In die zin is de \textit{sensus communis} een werking van Gods aanwezigheid. Instincten mogen, anders dan Leibniz meende, niet worden opgevat als affecten of verwarde voorstellingen, \textit{confusae repraesentationes}. Zij zijn niet vluchtig, maar diepgewortelde neigingen. Zij bezitten een dwingende, goddelijke en onweerstaanbare kracht.

Juist omdat deze instincten gaven van God zijn, is de \textit{sensus communis} van bijzonder belang voor onze kennis. Oetinger schrijft dat de \textit{ratio} zichzelf bestuurt door regels, vaak zelfs zonder God, maar dat het zintuig altijd met God verbonden blijft. Zoals de natuur verschilt van de kunst, zo verschilt het zintuig van de ratio. God werkt door de natuur in een gelijktijdige groei die zich regelmatig over het geheel verspreidt. Kunst daarentegen begint met een afzonderlijk deel. Het zintuig volgt de natuur na; de ratio volgt de kunst na.

Opmerkelijk is dat deze uitspraak uit een hermeneutische context voortkomt. In Oetingers geleerde werk vormt de \textit{Sapientia Salomonis}, de wijsheid van Salomo, het hoogste voorbeeld van kennis. De passage komt uit het hoofdstuk over het gebruik, de \textit{usus}, van de \textit{sensus communis}. Daar valt Oetinger de hermeneutische theorie van de Wolffiaanse school aan.

Belangrijker dan alle hermeneutische regels is volgens hem het \textit{sensu plenus} zijn: vol zijn van zintuig, gevoel of innerlijke zin. Deze gedachte is spiritualistisch en extreem geformuleerd, maar zij heeft wel een eigen logische grondslag in het begrip \textit{vita}, leven, of \textit{sensus communis}.

De hermeneutische betekenis ervan blijkt uit de gedachte dat de ideeën die in de Schrift en in de werken van God worden aangetroffen, des te vruchtbaarder en zuiverder zijn naarmate elk afzonderlijk idee in het geheel kan worden gezien en het geheel in elk afzonderlijk idee. Wat men in de negentiende en twintigste eeuw vaak ``intuïtie'' zou noemen, wordt hier dus teruggevoerd op een metafysische grondslag: de structuur van het levende, organische zijn, waarin het geheel in elk individu aanwezig is.

Volgens Oetinger heeft het geheel van het leven zijn centrum in het hart, dat door middel van de \textit{sensus communis} talloze dingen tegelijk begrijpt. Nog dieper dan kennis van hermeneutische regels ligt daarom de toepassing op zichzelf. Oetinger zegt: pas de regels vooral op uzelf toe, en dan zult u de sleutel hebben tot het verstaan van Salomo's spreuken.

Op deze basis kan Oetinger zijn gedachten verbinden met die van Shaftesbury, die volgens hem als enige onder deze titel over \textit{sensus communis} heeft geschreven. Maar hij verwijst ook naar anderen die de eenzijdigheid van de rationele methode hebben opgemerkt, zoals Pascal met zijn onderscheid tussen \textit{esprit géométrique} en \textit{esprit de finesse}. Toch krijgt het begrip \textit{sensus communis} bij deze Zwabische piëtist vooral een theologische betekenis, eerder dan een politieke of sociale.

Ook andere piëtistische theologen, zoals Rambach, benadrukten de toepassing van dit begrip tegen het heersende rationalisme. Dat blijkt bijvoorbeeld uit Rambachs invloedrijke hermeneutiek, waarin ook de toepassing een plaats kreeg. Maar toen de piëtistische tendensen later in de achttiende eeuw verdwenen, werd de hermeneutische functie van de \textit{sensus communis} gereduceerd tot een louter correctief. Wat in strijd is met de consensus van gevoelens, oordelen en conclusies — dat wil zeggen met de \textit{sensus communis} — kan niet juist zijn.

In vergelijking met de betekenis die Shaftesbury aan de \textit{sensus communis} gaf voor samenleving en staat, laat deze negatieve functie zien hoezeer het begrip in de Duitse Verlichting was uitgehold en geïntellectualiseerd.

\subsubsection{Oordeel}

De ontwikkeling van het begrip \textit{sensus communis} in het achttiende-eeuwse Duitsland verklaart wellicht waarom dit begrip zo nauw verbonden raakte met het begrip oordeel. \textit{Gesunder Menschenverstand}, gezond verstand, soms ook \textit{gemeiner Verstand}, gemeenschappelijk verstand, wordt immers wezenlijk gekenmerkt door het vermogen om te oordelen.

Het verschil tussen een dwaze en een verstandige mens is dat de eerste een gebrek aan oordeelsvermogen heeft. Hij is niet in staat om een bijzonder geval op de juiste manier onder een algemene regel te brengen, en kan daardoor niet goed toepassen wat hij geleerd heeft of weet. Het woord ``oordeel'' werd in de achttiende eeuw gebruikt als vertaling van \textit{judicium}, dat gold als een fundamentele intellectuele deugd.

De Engelse morele filosofen benadrukten dat morele en esthetische oordelen niet eenvoudig door de rede worden geleid, maar het karakter hebben van gevoel of smaak. Op vergelijkbare wijze beschouwde Tetens, een vertegenwoordiger van de Duitse Verlichting, de \textit{sensus communis} als een \textit{judicium} zonder reflectie.

De logische grondslag van het oordeel ligt in de subsumptie: het onderbrengen van een bijzonder geval onder een algemene regel, of het herkennen van iets als voorbeeld van een regel. Maar juist deze toepassing kan zelf niet worden bewezen. Het oordeel heeft dus een leidend principe nodig om regels te kunnen toepassen. Maar om dat principe toe te passen, zou opnieuw oordeelsvermogen nodig zijn. Kant heeft dit probleem scherp gezien.

Daarom kan oordeelsvermogen niet abstract worden aangeleerd. Het kan alleen van geval tot geval worden geoefend. In die zin lijkt het meer op een vaardigheid of zelfs op een zintuig. Men kan het niet eenvoudig onderwijzen door middel van begrippen of bewijzen, omdat geen enkele demonstratie op basis van begrippen volledig kan bepalen hoe regels in concrete gevallen moeten worden toegepast.

De Duitse Verlichtingsfilosofie rekende het oordeel daarom niet tot de hogere, maar tot de lagere vermogens van de geest. Daarmee week zij duidelijk af van de oorspronkelijke Romeinse betekenis van \textit{sensus communis}, terwijl zij tegelijk de scholastieke traditie voortzette. Dit werd vooral belangrijk voor de esthetica.

Baumgarten was er bijvoorbeeld van overtuigd dat het oordeel betrekking heeft op het waarneembare individuele ding. Wat het oordeel in dat individuele ding beoordeelt, is de volmaaktheid of onvolmaaktheid ervan. Opmerkelijk is dat oordeel hier niet simpelweg betekent dat men een vooraf gegeven begrip op een ding toepast. Het individuele ding wordt veeleer op zichzelf begrepen, voor zover het de overeenstemming van het vele met het ene laat zien.

Niet de toepassing van een algemeen begrip staat dus centraal, maar de innerlijke samenhang van het concrete verschijnsel zelf. Hier vinden we al wat Kant later het ``reflecterende oordeel'' zal noemen. Kant begrijpt dit als een oordeel volgens reële en formele doelmatigheid. Er is dan geen begrip vooraf gegeven. In plaats daarvan wordt het individuele object immanent beoordeeld.

Kant noemt dit een esthetisch oordeel. Zoals Baumgarten het \textit{iudicium sensitivum}, het zintuiglijke oordeel, als \textit{gustus}, smaak, beschreef, zo zegt ook Kant: ``Een zintuiglijk oordeel over volmaaktheid wordt smaak genoemd.''

We zullen later zien dat deze esthetische ontwikkeling van het begrip \textit{iudicium}, waarvoor vooral Gottsched in de achttiende eeuw belangrijk was, voor Kant een systematische betekenis kreeg. Tegelijk zal blijken dat Kants onderscheid tussen bepalend en reflecterend oordeel niet zonder problemen is.

Bovendien kan de betekenis van \textit{sensus communis} moeilijk worden teruggebracht tot het esthetische oordeel alleen. Uit het gebruik dat Vico en Shaftesbury van dit begrip maken, blijkt juist dat \textit{sensus communis} niet in de eerste plaats een formeel vermogen is dat men eenvoudig kan inzetten. Het begrip omvat al een geheel van oordelen en maatstaven voor oordeel, die de inhoud ervan mede bepalen.

De \textit{sensus communis} komt vooral tot uitdrukking in oordelen over goed en kwaad, gepast en ongepast. Wie over een gezond oordeelsvermogen beschikt, is niet alleen in staat om bijzondere gevallen onder algemene gezichtspunten te brengen. Hij weet vooral wat werkelijk belangrijk is. Hij ziet de dingen vanuit juiste en gezonde gezichtspunten.

Een bedrieger kan menselijke zwakheden heel goed inschatten en in zijn misleidingen steeds de juiste zet doen. Toch bezit hij daarmee nog geen ``gezond oordeel'' in de hogere zin van het woord. De algemeenheid die aan het oordeelsvermogen wordt toegeschreven, is daarom niet zo algemeen of alledaags als Kant lijkt te denken.

Oordeel is niet zozeer een vermogen waarover men eenvoudig beschikt, maar een eis die aan iedereen gesteld mag worden. Van ieder mens mag worden verwacht dat hij voldoende \textit{gemeinen Sinn}, zin voor het gemeenschappelijke, bezit om \textit{Gemeinsinn}, echte gemeenschapszin, te tonen. Dat wil zeggen: hij moet kunnen oordelen over goed en kwaad en zorg dragen voor het algemeen welzijn.

Juist daarom is Vico's beroep op de humanistische traditie zo indrukwekkend. Tegenover de intellectualisering van het begrip gemeenschapszin bewaart hij de rijkdom aan betekenissen die in de Romeinse traditie van het woord aanwezig was. Die rijkdom is tot op de dag van vandaag kenmerkend gebleven voor de Romaanse wereld.

Ook bij Shaftesbury, zoals we hebben gezien, bleef het begrip verbonden met de politieke en sociale traditie van het humanisme. De \textit{sensus communis} is een element van sociaal en moreel zijn. Zelfs wanneer het begrip later werd gebruikt in een polemiek tegen de metafysica, zoals in het piëtisme en de Schotse filosofie, behield het iets van zijn oorspronkelijke kritische functie.

Bij Kant krijgt dit idee in de \textit{Kritiek van het oordeelsvermogen} echter een heel andere nadruk. Voor de fundamentele morele betekenis van het begrip is bij hem geen systematische plaats meer. Zoals bekend ontwikkelde Kant zijn morele filosofie juist in uitdrukkelijke afwijzing van de leer van het ``morele gevoel'', zoals die in de Engelse filosofie was uitgewerkt. Daarom sluit hij het begrip \textit{sensus communis} uit de morele filosofie uit.

Wat in de onvoorwaardelijkheid van de morele imperatief naar voren komt, kan volgens Kant niet op gevoel worden gebaseerd. Dat geldt zelfs wanneer men met gevoel niet een individueel gevoel bedoelt, maar een algemene morele gevoeligheid. De imperatief die in de moraliteit zelf ligt besloten, sluit immers elke vergelijkende afweging met anderen uit.

Dat betekent niet dat het morele bewustzijn star of ongevoelig zou moeten oordelen over anderen. Integendeel: het is moreel verplicht zich los te maken van de subjectieve en particuliere voorwaarden van het eigen oordeel, en het standpunt van de ander in te nemen. Maar de onvoorwaardelijkheid van de morele imperatief betekent ook dat het morele bewustzijn zich niet kan baseren op het oordeel van anderen.

De verplichting van de imperatief is universeel in een strengere zin dan enige algemeenheid van gevoel ooit kan zijn. De toepassing van de morele wet op de wil is weliswaar een zaak van het oordeel. Maar omdat het hier gaat om een oordeel onder de wetten van de zuivere praktische rede, bestaat zijn taak er juist in ons te beschermen tegen het `empirisme van de praktische rede''. Dat empirisme zou de praktische begrippen van goed en kwaad alleen baseren op empirische gevolgen. Kant wil dit voorkomen door de zogenaamde `typiek'' van de zuivere praktische rede.

Voor Kant is er daarnaast nog een andere vraag: hoe kan de strenge wet van de zuivere praktische rede in de menselijke geest worden ingeplant? Deze vraag behandelt hij in de ``Methodologie van de zuivere praktische rede''. Daar probeert hij in grote lijnen te laten zien hoe echte morele gezindheden kunnen worden opgewekt en gevormd.

Daarvoor beroept Kant zich inderdaad op de gewone menselijke rede. Hij wil het praktische oordeelsvermogen oefenen en ontwikkelen, en daarbij spelen ook esthetische elementen een rol. Maar dat moreel gevoel kan worden gekweekt, behoort volgens Kant eigenlijk niet tot de grondslagen van de morele filosofie. Voor die grondslagen is het in elk geval niet beslissend.

Kant eist immers dat onze wil uitsluitend wordt bepaald door motieven die voortkomen uit de zelfwetgeving van de zuivere praktische rede. Dat kan niet worden gebaseerd op een louter gemeenschappelijkheid van gevoel. Het kan alleen worden gebaseerd op een weliswaar duistere, maar toch veilig leidende praktische daad van de wil. De taak van de \textit{Kritiek van de praktische rede} is om deze daad te verhelderen en te versterken.

Ook in logische zin speelt de \textit{sensus communis} bij Kant geen centrale rol. Wat Kant behandelt in de transcendentale leer van het oordeel — de leer van het schematisme en de grondbeginselen — heeft niet meer met de \textit{sensus communis} te maken. Daar gaat het om begrippen die \textit{a priori} op hun objecten moeten worden betrokken, niet om het gewone vermogen om het bijzondere onder het algemene te brengen.

Wanneer we werkelijk te maken hebben met het vermogen om een bijzonder geval als voorbeeld van een algemene regel te begrijpen, en wanneer we dit gezond verstand noemen, dan is dat volgens Kant iets wat in de ware zin van het woord ``gemeenschappelijk'' is: iets dat overal voorkomt. Maar juist daarom is het bezit ervan geen bijzondere verdienste.

De enige betekenis van dit gezonde verstand is dat het een voorstadium vormt van ontwikkelde en verlichte rede. Het werkt in een duistere vorm van oordeel die men gevoel noemt. Toch oordeelt het nog altijd volgens begrippen, al zijn die begrippen meestal slechts vaag voorgesteld. Het kan daarom volgens Kant niet worden beschouwd als een bijzondere \textit{Gemeinsinn}, een echt gevoel voor gemeenschap. Het algemene logische gebruik van het oordeel, dat men met de \textit{sensus communis} in verband kan brengen, bezit bij Kant geen eigen principe.

Van het hele gebied dat men een zintuiglijk oordeelsvermogen zou kunnen noemen, blijft bij Kant uiteindelijk alleen het oordeel van de esthetische smaak over. Alleen hier mag men volgens hem spreken van een werkelijk \textit{Gemeinsinn}. Hoe twijfelachtig het ook is of men bij esthetische smaak van kennis kan spreken, en hoe zeker het ook is dat esthetische oordelen niet volgens begrippen worden geveld, toch veronderstelt het smaakoordeel noodzakelijkerwijs algemene instemming. Die instemming is zintuiglijk en niet begripsmatig. Daarom is volgens Kant smaak het ware gevoel voor gemeenschap.

Deze formulering is paradoxaal. De achttiende eeuw discussieerde immers juist graag over de verscheidenheid van menselijke smaak. Zelfs wanneer men uit die verschillen geen sceptische of relativistische conclusie trekt, maar vasthoudt aan het idee van goede smaak, blijft het vreemd om juist ``goede smaak'' een gevoel voor gemeenschap te noemen. Goede smaak was immers vaak een onderscheiding waarmee de leden van een gecultiveerde samenleving zich van anderen onderscheidden.

Als empirische bewering zou dit inderdaad absurd zijn. Later zal blijken in hoeverre Kants beschrijving zin heeft binnen zijn transcendentale doelstelling: namelijk als een \textit{a priori} rechtvaardiging van een kritiek van de smaak. Maar we moeten ook vragen wat er gebeurt met de waarheidsaanspraak die oorspronkelijk in de \textit{sensus communis} besloten lag, wanneer het begrip wordt beperkt tot het smaakoordeel over het schone. Even belangrijk is de vraag hoe Kants subjectieve \textit{a priori} van de smaak het zelfbegrip van de menswetenschappen heeft beïnvloed.

Hier is een geredigeerde versie van je vertaling. Ik heb vooral gewerkt aan vloeiender Nederlands, correcties van enkele kleine verschrijvingen, terminologische consistentie en duidelijkere zinsbouw. Ik baseer me op de door jou aangeleverde tekst. 

\subsubsection{Smaak}

Om de ontwikkeling van het begrip smaak in zijn volle omvang te begrijpen, moeten we verder teruggaan in de tijd. Het gaat hier niet alleen om de beperking van het begrip ``gemeenschapszin'' tot smaak, maar ook om de beperking van het begrip smaak zelf. De lange geschiedenis van dit idee vóór Kant — een idee dat hij tot grondslag maakte van zijn \textit{Kritiek van het oordeelsvermogen} — laat zien dat smaak oorspronkelijk eerder een moreel dan een esthetisch begrip was. Het beschreef een ideaal van ware menselijkheid en kreeg zijn betekenis in de poging om kritisch afstand te nemen van het dogmatisme van de `'`school''. Pas later werd het begrip beperkt tot het ``esthetische''.

Aan het begin van deze geschiedenis staat Balthasar Gracián. Hij gaat uit van de gedachte dat de smaakzin — de meest dierlijke en tegelijk meest innerlijke van onze zintuigen — toch al de kiem bevat van de intellectuele onderscheidingen die wij maken wanneer wij dingen beoordelen. De zintuiglijke onderscheiding van smaak, die op de meest directe manier aanvaardt of verwerpt, is dus niet louter instinct. Zij vormt een tussengebied tussen zintuiglijke aandrift en intellectuele vrijheid.

De smaakzin is in staat de afstand te scheppen die nodig is om keuzes te maken en te oordelen over wat tot de meest dringende levensbehoeften behoort. Gracián ziet in smaak daarom al een ``vergeestelijking van de dierlijkheid''. Hij wijst terecht op het bestaan van cultuur, \textit{cultura}, niet alleen van de geest, \textit{ingenio}, maar ook van smaak, \textit{gusto}. Dat geldt natuurlijk ook voor de zintuiglijke smaak: er zijn mensen met ``een goede tong'', fijnproevers die hun smaak ontwikkelen en verfijnen.

Dit idee van \textit{gusto} vormt het uitgangspunt voor Graciáns ideaal van sociale vorming. Zijn ideaal van de gevormde mens, \textit{el discreto}, is de mens die, als \textit{hombre en su punto}, als iemand die tot volle rijpheid is gekomen, de juiste vrije afstand bereikt tot alle zaken van het leven en de samenleving. Daardoor kan hij bewust en reflecterend onderscheiden en kiezen.

Graciáns ideaal van \textit{Bildung} vormde een volledig nieuw vertrekpunt. Het kwam in de plaats van het ideaal van de christelijke hoveling, zoals bij Castiglione, en is opmerkelijk in de geschiedenis van westerse vormingsidealen omdat het niet gebonden was aan sociale klasse. Het stelde het ideaal van een door \textit{Bildung} gevormde samenleving voorop.

Dit ideaal van sociale vorming lijkt overal op te komen in het spoor van het absolutisme en de terugdringing van de erfelijke aristocratie. Zo volgt de geschiedenis van het idee smaak de geschiedenis van het absolutisme: van Spanje naar Frankrijk en Engeland. Zij is nauw verbonden met de voorlopers van de derde stand. Smaak is niet alleen een ideaal dat door een nieuwe samenleving wordt voortgebracht; omgekeerd brengt dit ideaal van ``goede smaak'' ook voort wat later de ``goede samenleving'' werd genoemd. Die samenleving legitimeert zich niet langer door afkomst of rang, maar door de gemeenschappelijkheid van haar oordeel — of beter: door het vermogen om boven enge belangen en persoonlijke voorkeuren uit te stijgen tot het niveau van oordeel.

Het begrip smaak impliceert ongetwijfeld een vorm van kennen. Goede smaak wordt gekenmerkt door het vermogen afstand te nemen van zichzelf en van de eigen voorkeuren. Smaak is daarom in wezen niet privé, maar een sociaal verschijnsel van de eerste orde. Zij kan zelfs ingaan tegen de persoonlijke neigingen van het individu, alsof zij in naam van een bedoelde en vertegenwoordigde universaliteit als een rechtbank optreedt. Men kan iets mooi vinden dat de eigen smaak toch afkeurt.

Het smaakoordeel is op een opvallende manier beslissend. Zoals men zegt: \textit{de gustibus non est disputandum}, over smaak valt niet te twisten. Kant zegt terecht dat er in zaken van smaak wel meningsverschillen kunnen zijn, maar geen bewijzende discussies. Dat komt niet simpelweg doordat er geen universele, begripsmatige criteria bestaan die iedereen moet accepteren. Het komt ook doordat men naar zulke criteria niet zoekt en het zelfs ongepast zou vinden als ze bestonden.

Smaak moet men hebben. Men kan haar niet leren door bewijsvoering, en zij kan ook niet worden vervangen door louter imitatie. Toch is smaak geen zuiver privébezit, want zij streeft altijd naar goede smaak. De beslissendheid van het smaakoordeel omvat daarom een aanspraak op geldigheid. Goede smaak is zeker van haar oordeel. Zij is wezenlijk zekere smaak: een aanvaarden en afwijzen zonder aarzeling, zonder heimelijk naar anderen te kijken en zonder naar redenen te zoeken.

Smaak is daarom iets als een zintuig. In haar werking kent zij geen redenen. Wanneer smaak iets afwijst, kan zij vaak niet precies zeggen waarom. Toch ervaart zij die afwijzing met grote zekerheid. Zekerheid van smaak betekent daarom vooral: zekerheid tegenover het smakeloze.

Opmerkelijk is dat wij vooral gevoelig zijn voor het negatieve in de oordelen van smaak. Het positieve tegendeel is niet zozeer wat ``smakelijk'' is, maar wat niet tegen de smaak ingaat. Dat is vooral wat smaak beoordeelt. Smaak wordt juist gekenmerkt door het feit dat zij zich stoort aan het smakeloze en het daarom vermijdt, zoals men alles vermijdt wat letsel dreigt toe te brengen. Het tegendeel van goede smaak is dan ook niet slechte smaak, maar veeleer het ontbreken van smaak. Goede smaak is een gevoeligheid die het smakeloze zo vanzelfsprekend vermijdt dat haar reactie voor iemand zonder smaak volledig onbegrijpelijk kan zijn.

Een verschijnsel dat nauw met smaak verbonden is, is mode. In mode wordt het element van sociale veralgemening, dat in het begrip smaak besloten ligt, een bepalende werkelijkheid. Maar juist het verschil met mode laat zien dat de universaliteit van smaak een andere grondslag heeft en niet hetzelfde is als empirische algemeenheid. Dat is voor Kant het beslissende punt.

Het woord ``mode'', \textit{Mode}, wijst er al op dat het gaat om een veranderlijke regel, een \textit{modus}, binnen een betrekkelijk vast geheel van omgangsvormen. Wat louter mode is, heeft geen andere norm dan wat iedereen doet. Mode reguleert naar willekeur dingen die evengoed anders hadden kunnen zijn. Zij wordt bepaald door empirische algemeenheid: door rekening te houden met anderen, door vergelijking en door de dingen te bekijken vanuit het algemene standpunt van wat men nu eenmaal doet. Zo schept mode een sociale afhankelijkheid die moeilijk af te schudden is. Kant heeft gelijk wanneer hij zegt dat het beter is een dwaas in de mode te zijn dan een dwaas tegen de mode in — ook al blijft het dwaas om mode te serieus te nemen.

Smaak daarentegen is een geestelijk vermogen om te onderscheiden. Smaak functioneert wel binnen een gemeenschap, maar is daaraan niet onderworpen. Integendeel: goede smaak blijkt juist uit het vermogen zich te verhouden tot de smaakrichting die door de mode wordt vertegenwoordigd. Zij kan zich daaraan aanpassen, maar ook omgekeerd datgene wat de mode voorschrijft aanpassen aan de eigen goede smaak.

Bij smaak hoort daarom ook maat houden, zelfs in de mode. Men volgt niet klakkeloos de wisselende bevelen van de mode, maar gebruikt het eigen oordeel. Men handhaaft zijn eigen ``stijl''. Dat wil zeggen: men brengt de eisen van de mode in verband met een geheel dat de eigen smaak voor ogen heeft, en aanvaardt alleen wat daarmee harmonieert en samenpast.

Smaak herkent dus niet alleen dit of dat als mooi, maar heeft ook oog voor het geheel waarin alles wat mooi is moet passen. Daarom is smaak geen sociaal gevoel in de zin van afhankelijkheid van empirische algemeenheid, alsof het zou gaan om de feitelijke eenstemmigheid van de oordelen van anderen. Smaak zegt niet dat iedereen het met ons oordeel eens zal zijn, maar dat men het ermee eens zou moeten zijn, zoals Kant zegt.

Tegenover de tirannie van de mode bewaart zekere smaak een eigen vrijheid en superioriteit. Daarin ligt haar bijzondere normatieve kracht: zij weet zich zeker van de instemming van een ideale gemeenschap. In tegenstelling tot een smaak die door mode wordt beheerst, toont zich hier de idealiteit van goede smaak. Daaruit volgt dat smaak iets weet, maar op een manier die niet losgemaakt kan worden van het concrete moment waarin het object verschijnt en die niet kan worden herleid tot regels en begrippen.

Juist dit geeft het begrip smaak zijn oorspronkelijke breedte: smaak vormt een bijzondere wijze van kennen. Net als het reflecterende oordeel behoort zij tot het domein waarin men in het individuele object het universele begrijpt waaronder het zou moeten vallen. Zowel smaak als oordeel beoordelen een object in verhouding tot een geheel, om te zien of het daarin past — dat wil zeggen: of het passend is. Daarvoor moet men een gevoel hebben; het kan niet worden gedemonstreerd.

Deze manier van voelen is duidelijk nodig waar een geheel bedoeld is, maar niet als geheel gegeven is — dat wil zeggen: niet in doelbegrippen kan worden vastgelegd. Daarom is smaak geenszins beperkt tot het schone in natuur en kunst, alsof zij alleen de decoratieve kwaliteit daarvan beoordeelt. Smaak omvat ook het hele gebied van moraliteit en omgangsvormen.

Ook morele begrippen zijn nooit als geheel gegeven of normatief volledig bepaald. Integendeel: de ordening van het leven door regels van recht en moraal blijft onvolledig en vraagt om productieve aanvulling. Oordeel is nodig om het concrete geval juist te beoordelen. Deze functie van het oordeel kennen we vooral uit de rechtswetenschap, waar de aanvullende functie van de hermeneutiek bestaat in het concretiseren van de wet.

Het gaat daarbij altijd om meer dan de juiste toepassing van algemene principes. Ook onze kennis van recht en moraal wordt steeds aangevuld door het individuele geval, en zelfs mede door dat geval bepaald. De rechter past de wet niet alleen concreet toe, maar draagt door zijn oordeel ook bij aan de ontwikkeling van de wet, zoals in het \textit{judge-made law}. Net als het recht ontwikkelt ook de moraliteit zich voortdurend door de vruchtbaarheid van het individuele geval.

Daarom is oordeel, als beoordeling van het schone en het verhevene, allerminst alleen productief op het gebied van natuur en kunst. Men kan zelfs niet, met Kant, zeggen dat de productiviteit van het oordeel vooral op dat gebied te vinden is. Integendeel: het schone in natuur en kunst moet worden aangevuld door de hele oceaan van het schone die verspreid ligt over de morele werkelijkheid van de mensheid.

Alleen met betrekking tot de uitoefening van zuivere theoretische en praktische rede kan men spreken van het onderbrengen van het individuele onder een gegeven universeel, dus van Kants bepalende oordeel. Maar zelfs hier is een esthetisch element betrokken. Kant geeft dit indirect toe wanneer hij de waarde van voorbeelden erkent voor het scherpen van het oordeel. Wel voegt hij daaraan toe dat voorbeelden de juistheid en precisie van het intellectuele inzicht meestal enigszins belemmeren, omdat zij zelden volledig aan de eisen van de regel voldoen als \textit{casus in terminis}.

Maar de andere kant van deze opmerking is duidelijk: het geval dat als voorbeeld functioneert, is in werkelijkheid iets anders dan slechts een geval van de regel. Daarom bevat het recht doen aan het concrete geval — zelfs wanneer het alleen om technisch of praktisch oordeel gaat — altijd een esthetisch element.

In zoverre is het onderscheid tussen bepalend en reflecterend oordeel, waarop Kant zijn kritiek van het oordeel baseert, niet absoluut. Het gaat duidelijk niet alleen om logisch, maar ook om esthetisch oordeel. Het individuele geval waarop het oordeel betrekking heeft, is nooit louter een geval. Het wordt niet uitgeput door een bijzonder voorbeeld van een algemene wet of een algemeen begrip te zijn.

Integendeel: het is altijd een individueel geval. Het is veelzeggend dat wij spreken van een bijzonder geval, omdat de regel het nooit volledig omvat. Elk oordeel over iets dat in zijn concrete individualiteit bedoeld is — bijvoorbeeld het oordeel dat nodig is in een situatie waarin gehandeld moet worden — is, strikt genomen, een oordeel over een bijzonder geval.

Dat betekent dat het beoordelen van een geval niet alleen inhoudt dat het algemene principe wordt toegepast volgens welk het wordt beoordeeld. Het betekent ook dat dit principe mede wordt bepaald, aangevuld en gecorrigeerd. Hieruit volgt uiteindelijk dat alle morele beslissingen smaak vereisen. Dat betekent niet dat deze uiterst individuele afweging de enige grondslag van de beslissing is, maar wel dat zij een onmisbaar element vormt.

Het is werkelijk een prestatie van niet te demonstreren tact om het juiste punt te treffen en de toepassing van het universele — de morele wet, bij Kant — zo te sturen dat de rede dit niet alleen kan doen. Smaak is daarom niet de grondslag, maar wel de hoogste voltooiing van het morele oordeel. De mens die ervaart dat het slechte tegen zijn smaak ingaat, bezit de grootste zekerheid in het aanvaarden van het goede en het afwijzen van het slechte — een zekerheid die vergelijkbaar is met die van onze meest vitale zintuigen, die voedsel kiezen of verwerpen.

De opkomst van het begrip smaak in de zeventiende eeuw, en de sociale en maatschappelijk verbindende functie die wij hierboven hebben aangeduid, heeft banden met de morele filosofie die teruggaan tot de oudheid. In de christelijke morele filosofie werkt een humanistisch en uiteindelijk Grieks element door. De Griekse ethiek — de ethiek van de maat bij de Pythagoreeërs en Plato, en de ethiek van het midden, \textit{mesotès}, bij Aristoteles — is in diepe en brede zin een ethiek van de goede smaak.

Zo'n stelling klinkt ons misschien vreemd in de oren. Dat komt deels doordat wij het normatieve element in het begrip smaak meestal niet meer herkennen en nog steeds beïnvloed zijn door het relativistisch-sceptische argument dat smaken verschillen. Maar vooral komt het door Kants prestatie in de morele filosofie, die de ethiek heeft gezuiverd van alle esthetiek en gevoel.

Als wij nu de betekenis van Kants \textit{Kritiek van het oordeelsvermogen} voor de geschiedenis van de menswetenschappen willen onderzoeken, moeten we zeggen dat zijn transcendentale fundering van de esthetiek grote gevolgen had en een keerpunt vormde. Zij betekende het einde van een traditie, maar ook het begin van een nieuwe ontwikkeling.

Kant beperkte het idee van smaak tot een gebied waarin zij, als bijzonder principe van oordeel, zelfstandige geldigheid kon opeisen. Dat deed hij door het begrip kennis te beperken tot het theoretische en praktische gebruik van de rede. Het tot het schone en verhevene beperkte verschijnsel van oordeel was voldoende voor zijn transcendentale doel. Maar daardoor raakten de bredere ervaring van smaak en de werking van esthetisch oordeel in recht en moraal uit het centrum van de filosofie.

Het belang hiervan kan nauwelijks worden overschat. Wat hier werd prijsgegeven, was precies het element waarin filologische en historische studies leefden. Toen deze studies zich later, onder de naam ``menswetenschappen'', methodologisch probeerden te funderen naast de natuurwetenschappen, was juist dit element de enige mogelijke bron van hun volledige zelfbegrip. Kants transcendentale analyse maakte het echter onmogelijk om de waarheidsaanspraak van de overgeleverde teksten en tradities waaraan deze studies zich wijdden, werkelijk te erkennen. Daarmee verloor ook de methodologische eigenheid van de menswetenschappen haar legitimiteit.

In zijn kritiek van het esthetische oordeel legitimeerde Kant wat hij wilde legitimeren: de subjectieve universaliteit van de esthetische smaak, waarin geen kennis van het object meer ligt, en op het gebied van de schone kunsten de superioriteit van het genie boven elke op regels gebaseerde esthetiek.

Daarmee vonden de romantische hermeneutiek en geschiedschrijving een aanknopingspunt voor hun zelfbegrip in het begrip genie, dat door Kants esthetiek was gelegitimeerd. Dat was de andere kant van Kants invloed. De transcendentale rechtvaardiging van het esthetische oordeel werd de grondslag van de autonomie van het esthetische bewustzijn, en op overeenkomstige wijze moest ook het historische bewustzijn worden gelegitimeerd.

De radicale subjectivering die met Kants nieuwe fundering van de esthetica gepaard ging, was in strikte zin baanbrekend voor een nieuw tijdperk. Doordat zij elke vorm van theoretische kennis buiten de natuurwetenschap diskwalificeerde, dwong zij de menswetenschappen ertoe zichzelf te begrijpen naar het model van de natuurwetenschappelijke methode. Tegelijk maakte zij deze afhankelijkheid gemakkelijker door het ``artistieke element'', ``gevoel'' en ``inleving'' als aanvullende elementen aan te bieden. Helmholtz' beschrijving van de menswetenschappen, die hierboven is besproken, is in beide opzichten een goed voorbeeld van deze Kantiaanse invloed.

Als wij willen aantonen wat onvoldoende is aan deze zelfinterpretatie van de menswetenschappen en betere mogelijkheden willen openen, moeten wij verdergaan met de problemen van de esthetiek. De transcendentale functie die Kant aan het esthetische oordeel toekent, is voldoende om het te onderscheiden van begripsmatige kennis en zo de verschijnselen van het schone en de kunst te bepalen. Maar is het juist om het begrip waarheid voor te behouden aan begripsmatige kennis? Moeten wij niet ook erkennen dat het kunstwerk waarheid bezit? We zullen zien dat de erkenning hiervan niet alleen het verschijnsel kunst, maar ook het verschijnsel geschiedenis in een nieuw licht stelt.



\section{De subjectivering van de esthetiek door de Kantiaanse kritiek}

\subsection{Kants leer van smaak en genie}

\subsubsection{De transcendentale eigenheid van smaak}

Bij het onderzoeken van de grondslagen van smaak verraste Kant zichzelf door een \textit{a priori} element aan te treffen dat verder ging dan empirische universaliteit. Dit inzicht gaf de aanzet tot de \textit{Kritiek van het oordeelsvermogen}. Het is niet langer louter een kritiek op smaak in de zin dat smaak het voorwerp is van een kritisch oordeel door een waarnemer. Het is een kritiek \textit{van} de kritiek; dat wil zeggen, het gaat om de legitimiteit van een dergelijke kritiek op het gebied van smaak. De kwestie is niet langer de empirische principes die een wijdverspreide en dominante smaak zouden moeten rechtvaardigen — zoals bijvoorbeeld in de oude discussie over de oorsprong van verschillen in smaak — maar het gaat om een echte \textit{a priori} dat op zichzelf de mogelijkheid van kritiek volledig zou kunnen rechtvaardigen. Wat zou een dergelijke rechtvaardiging kunnen vormen?

Het is duidelijk dat de geldigheid van een esthetisch oordeel niet kan worden afgeleid of bewezen uit een universeel principe. Niemand veronderstelt dat kwesties van smaak door argumenten en bewijzen kunnen worden beslecht. Even duidelijk is dat goede smaak nooit werkelijk empirische universaliteit zal bereiken, en dat een beroep doen op de heersende smaak de ware aard van smaak mist. In het begrip smaak zelf ligt besloten dat het zich niet klakkeloos onderwerpt aan populaire waarden en voorkeursmodellen, en deze eenvoudigweg imiteert. In het domein van de esthetische smaak hebben modellen en patronen zeker een bevoorrechte functie; maar, zoals Kant terecht opmerkt, zijn zij niet bedoeld om te worden geïmiteerd, maar om te worden gevolgd. Het model en het voorbeeld moedigen smaak aan om zijn eigen weg te gaan, maar zij doen het werk van smaak niet voor hem. Want smaak moet men zelf bezitten.

Aan de andere kant heeft onze schets van de geschiedenis van het begrip smaak duidelijk genoeg aangetoond dat niet de bijzondere voorkeuren bepalend zijn; maar dat in het geval van een esthetisch oordeel een supra-empirische norm werkzaam is. We zullen zien dat Kants fundering van de esthetiek op het oordeel van smaak recht doet aan beide aspecten van het fenomeen: de empirische niet-universaliteit ervan en de \textit{a priori} claim op universaliteit.

Maar de prijs die hij betaalt voor deze legitimering van kritiek op het gebied van smaak is dat hij smaak elke betekenis als kennis ontzegt. Hij reduceert de \textit{sensus communis} tot een subjectief principe. In smaak wordt niets geweten van de objecten die als mooi worden beoordeeld; er wordt slechts gesteld dat er een gevoel van genoegen met hen verbonden is \textit{a priori} in het subjectieve bewustzijn. Zoals bekend, ziet Kant dit gevoel als gebaseerd op het feit dat de voorstelling van het object geschikt (\textit{zweckmäßig}) is voor ons vermogen tot kennis. Het is een vrij spel van verbeelding en verstand, een subjectieve relatie die geheel passend is bij kennis en die de reden voor het genoegen in het object aantoont. Deze geschiktheid voor het subject is in principe voor iedereen hetzelfde — dat wil zeggen, universeel mededeelbaar — en grondt zo de claim dat het oordeel van smaak universele geldigheid bezit.

Dit is het principe dat Kant in het esthetisch oordeel ontdekt. Het is diens eigen wet. Zo is het een \textit{a priori} effect van het schone, gelegen halverwege tussen een louter zintuiglijke, empirische overeenstemming in zaken van smaak en de rationalistische universaliteit van een regel. Geef toe: als men de relatie tot het \textit{Lebensgefühl} als enige grondslag neemt, kan men smaak niet langer een \textit{cognitio sensitiva} noemen. Het verschaft geen kennis van het object, maar het is ook niet louter een kwestie van een subjectieve reactie, zoals die wordt opgeroepen door wat aangenaam is voor de zintuigen. Smaak is reflecterend.

Wanneer Kant smaak dan ook het ware gezond verstand noemt, denkt hij niet langer aan de grote morele en politieke traditie van het begrip \textit{sensus communis} die we hierboven hebben geschetst. Integendeel, hij ziet dit idee als samengesteld uit twee elementen: ten eerste de universaliteit van smaak voor zover deze het resultaat is van het vrije spel van al onze kenvermogens en niet beperkt is tot een specifiek gebied, zoals een uitwendige zintuiglijke waarneming; ten tweede de gemeenschappelijke kwaliteit van smaak, voor zover deze, volgens Kant, abstractie maakt van alle subjectieve, private voorwaarden zoals aantrekkingskracht en emotie. In beide opzichten wordt de universaliteit van deze zin dus negatief gedefinieerd door het contrast met datgene waarvan hij abstractie maakt, en niet positief door wat gemeenschappelijkheid schept en gemeenschap vormt.

Toch blijft voor Kant de oude verbinding tussen smaak en sociabiliteit geldig. Maar de cultuur van smaak wordt slechts als een aanhangsel behandeld onder de titel De methodologie van de smaak. Daar wordt de \textit{humaniora}, zoals vertegenwoordigd door het Griekse model, gedefinieerd als de sociabiliteit die past bij de menselijkheid, en het cultiveren van moreel gevoel wordt aangeduid als de weg waarop echte smaak een bepaalde onveranderlijke vorm aanneemt. Zo zijn de specifieke inhoudelijke aspecten van smaak irrelevant voor zijn transcendentale functie. Kant is alleen geïnteresseerd voor zover er sprake is van een bijzonder principe van esthetisch oordeel, en daarom is hij alleen geïnteresseerd in het zuivere oordeel van smaak.

In overeenstemming met zijn transcendentale bedoeling put de ``Analytiek van de smaak'' zijn voorbeelden van esthetisch genoegen tamelijk onverschillig uit natuurlijke schoonheid, het decoratieve en artistieke voorstellingen. Het type object waarvan de idee genoegen verschaft, beïnvloedt de essentie van het esthetisch oordeel niet. De kritiek van het esthetisch oordeel streeft er niet naar een filosofie van de kunst te zijn — hoezeer kunst ook een voorwerp van dit oordeel is. Het begrip zuiver esthetisch oordeel van smaak is een methodologische abstractie die slechts zijdelings gerelateerd is aan het onderscheid tussen natuur en kunst. Daarom is het, bij een nauwkeuriger onderzoek van Kants esthetiek, noodzakelijk om die interpretaties weer in het juiste perspectief te plaatsen die zijn esthetiek lezen als een filosofie van de kunst, interpretaties die zich met name baseren op het begrip genie. Daartoe zullen we Kants opmerkelijke en omstreden leer van vrije en afhankelijke schoonheid bekijken.

\subsubsection{De leer van vrije en afhankelijke schoonheid}

Kant bespreekt hier het onderscheid tussen het ``zuivere'' en het ``geïntellectualiseerde'' oordeel van smaak, wat overeenkomt met het contrast tussen ``vrije'' en ``afhankelijke'' schoonheid (dat wil zeggen, afhankelijk van een begrip). Dit is een bijzonder riskante leer voor het begrip van kunst, want de vrije schoonheid van de natuur en — in de sfeer van de kunst — de versiering verschijnen als de schoonheid die eigen is aan het zuivere oordeel van smaak, omdat deze dingen ``in zichzelf'' mooi zijn. Waar een begrip wordt ingebracht — en dat is het geval niet alleen in de poëzie, maar in alle afbeeldende kunst — lijkt de situatie hetzelfde als in de voorbeelden van ``afhankelijke'' schoonheid die Kant noemt.

Zijn voorbeelden — mens, dier, gebouw — zijn natuurlijke dingen zoals die voorkomen in een wereld die wordt gedomineerd door menselijke doelen, of dingen die door de mens zijn vervaardigd voor menselijke doelen. In elk geval beperkt het feit dat het ding een bepaald doel dient het esthetische genoegen dat het kan bieden. Zo is voor Kant tatoeëren, het versieren van het menselijk lichaam, afkeurenswaardig, ook al kan het ``onbemiddeld'' genoegen opwekken. Zeker, Kant spreekt hier niet over kunst als zodanig (de ``mooie voorstelling van een ding''), maar nadrukkelijker over mooie dingen (van de natuur of architectuur).

Het onderscheid tussen natuurlijke en artistieke schoonheid, dat hij zelf later bespreekt (§48), is hier niet van belang; maar wanneer hij onder de voorbeelden van vrije schoonheid, naast bloemen, ook een tapijt met arabesken en muziek noemt (``zonder thema'' of zelfs ``zonder tekst''), dan wijst dat indirect op alle dingen die vallen onder ``objecten die onder een bepaald begrip vallen'' en derhalve moeten worden ingedeeld onder voorwaardelijke, onvrije schoonheid: het hele rijk van de poëzie, de beeldende kunsten en de architectuur, alsmede alle objecten van de natuur die wij niet louter beoordelen op hun schoonheid, zoals wij dat wel doen bij decoratieve bloemen. In al deze gevallen wordt het oordeel van smaak vertroebeld en beperkt. Het lijkt onmogelijk om recht te doen aan kunst als de esthetiek is gebaseerd op het ``zuivere oordeel van smaak'' — tenzij het criterium van smaak slechts als voorwaarde wordt gesteld.

De introductie van het begrip genie in de latere delen van de \textit{Kritiek van het oordeelsvermogen} kan aldus worden begrepen. Maar dat zou betekenen dat er achteraf een verschuiving in het zwaartepunt plaatsvindt. Want dit is in eerste instantie niet de kwestie. Hier (in §16) is het standpunt van smaak zo ver verwijderd van louter een voorwaarde te zijn, dat het integendeel claimt de aard van het esthetische oordeel volledig te omvatten en het te beschermen tegen beperkingen door ``intellectuele'' criteria. En hoewel Kant ziet dat hetzelfde object vanuit twee verschillende gezichtspunten kan worden beoordeeld — die van vrije en afhankelijke schoonheid — lijkt de ideale rechter van smaak toch degene te zijn die oordeelt volgens ``wat hij voor zijn zintuigen heeft'' en niet volgens ``wat hij voor zijn gedachten heeft''. Ware schoonheid is die van bloemen en versieringen, die in onze door doelen gedomineerde wereld zich als schoonheden onmiddellijk en in zichzelf presenteren, en die derhalve niet vereisen dat enig begrip of doel bewust wordt genegeerd.

Als men echter iets nauwkeuriger kijkt, past deze opvatting noch bij Kants woorden noch bij zijn onderwerp. De veronderstelde verschuiving in Kants standpunt van smaak naar genie vindt niet plaats; men hoeft slechts te leren herkennen hoe in het begin al een verborgen voorbereiding wordt getroffen voor wat later wordt ontwikkeld. Er is geen twijfel dat Kant de beperkingen die een mens verbieden getatoeëerd te worden of een kerk met een bepaalde versiering te decoreren, niet betreurt, maar juist eist; Kant ziet de daaruit voortvloeiende vermindering van esthetisch genoegen als een winst vanuit moreel oogpunt. De voorbeelden van vrije schoonheid zijn kennelijk niet bedoeld om de eigenlijke schoonheid te tonen, maar slechts om ervoor te zorgen dat genoegen als zodanig geen oordeel is over de volmaaktheid van het object.

En hoewel Kant aan het einde van deze sectie (§16) gelooft dat het onderscheid tussen de twee soorten schoonheid — of liever tussen de twee verhoudingen tot het schone — hem in staat stelt veel kritische geschillen over schoonheid te beslechten, is deze mogelijkheid om smaakgeschillen te beslechten slechts, zo men wil, een gevolg van de samenwerking tussen de twee benaderingen. Inderdaad zullen deze twee benaderingen het vaakst verenigd zijn.

Deze eenheid zal altijd gegeven zijn waar ``het kijken naar een begrip'' de vrijheid van de verbeelding niet afschaft. Zonder zichzelf te weerspreken, kan Kant beschrijven dat het een legitieme voorwaarde van esthetisch genoegen is dat er geen conflict is met doelmatige elementen. En zoals het kunstmatig was om schoonheden die in zichzelf vrij bestaan te isoleren (``smaak'', in elk geval, lijkt zich het meest te bewijzen waar niet alleen het juiste wordt gekozen, maar het juiste voor de juiste plaats), zo kan en moet men ook verder gaan dan het standpunt van het zuivere oordeel van smaak door te zeggen dat men zeker niet over schoonheid kan spreken wanneer een bepaald begrip van het verstand schematisch door de verbeelding wordt geïllustreerd, maar alleen wanneer de verbeelding in vrije harmonie is met het verstand — dat wil zeggen, waar deze productief kan zijn. Deze productiviteit van de verbeelding is niet het rijkst waar deze louter vrij is, zoals in de krullen van de arabesken, maar veeleer in een speelveld waar de wens van het verstand naar eenheid de verbeelding niet zoveel beperkt als wel prikkels tot spel aanreikt.

\subsubsection{De leer van het schoonheidsideaal}

Deze laatste opmerkingen gaan verder dan wat daadwerkelijk in Kants tekst staat, maar de gang van zijn gedachten (§17) rechtvaardigt deze interpretatie. De balans in deze sectie wordt pas duidelijk na een zorgvuldige bestudering.

Het normatieve idee van schoonheid dat daar uitvoerig wordt besproken, is niet het hoofdpunt en vertegenwoordigt niet het ideaal van schoonheid waar smaak van nature naar streeft. Er is alleen een ideaal van schoonheid als het gaat om de menselijke gedaante, in de ``uitdrukking van het morele'', ``zonder welke het object niet universeel bevallig zou kunnen zijn''. Een oordeel volgens een ideaal van schoonheid is, zoals Kant zegt, dan ook geen louter smaakoordeel. Het belangrijke gevolg van deze leer zal blijken te zijn dat iets meer moet zijn dan louter smaakvol bevallig om als kunstwerk te bevallen.

Dit is werkelijk verbazingwekkend. Hoewel we zojuist zagen dat echte schoonheid leek uit te sluiten dat het gebonden is aan doelen, wordt hier het tegendeel gesteld voor een mooi huis, een mooie boom, een mooie tuin, enzovoort — namelijk dat we ons geen ideaal van deze dingen kunnen voorstellen, ``omdat deze doelen \textit{niet voldoende} [cursief van mij] door hun begrip zijn bepaald en vastgelegd; en bijgevolg hun doelmatigheid (\textit{Zweckmäßigkeit}) bijna net zo vrij is als in het geval van schoonheid die geheel ongebonden is''. Er is alleen een ideaal van schoonheid voor de menselijke gedaante, precies omdat deze alleen in staat is tot een schoonheid die door een doelbegrip is vastgelegd! Deze leer, die door Winckelmann en Lessing naar voren is gebracht, neemt een sleutelpositie in in Kants fundering van de esthetiek. En deze these toont duidelijk aan hoe weinig een formele esthetiek van smaak (een arabeske esthetiek) overeenkomt met het Kantiaanse idee.

De leer van het ideaal van schoonheid is gebaseerd op het onderscheid tussen het normatieve idee en het rationele idee of ideaal van schoonheid. Het esthetische normatieve idee komt voor in alle natuurlijke geslachten. De manier waarop een mooi dier (bijvoorbeeld een koe: Myron) eruit zou moeten zien, is de maatstaf waaraan het individuele voorbeeld wordt beoordeeld. Dit normatieve idee is dus een enkele intuïtie van de verbeelding als ``het beeld van het geslacht dat zweeft tussen alle enkelvoudige individuen''. De voorstelling van zo'n normatief idee wekt echter geen genoegen op door zijn schoonheid, maar slechts ``omdat het geen voorwaarde tegenspreekt waaronder alleen een ding dat tot dit geslacht behoort, mooi kan zijn''. Het is niet het prototype van schoonheid, maar slechts van juistheid.

Dit geldt ook voor het normatieve idee van de menselijke gedaante. Maar er is een waar ideaal van de schoonheid van de menselijke gedaante in de ``uitdrukking van het morele''. Uitdrukking van het morele: als we dat combineren met de latere leer van esthetische ideeën en van schoonheid als symbool van moraliteit, dan zien we dat de leer van het ideaal van schoonheid ook een plaats bereidt voor de essentie van kunst. De toepassing op de kunsttheorie in de geest van Winckelmanns classicisme is evident. Het is duidelijk dat Kant bedoelt dat in de voorstelling van de menselijke gedaante het voorgestelde object samenvalt met de artistieke betekenis die tot ons spreekt in de voorstelling. Er kan geen andere betekenis in deze voorstelling zijn dan die al tot uiting komt in de gedaante en verschijning van wat wordt voorgesteld. In Kantiaanse termen: het geïntellectualiseerde en geïnteresseerde genoegen in dit voorgestelde ideaal van schoonheid leidt ons niet af van het esthetische genoegen, maar is er juist één mee. Alleen in de voorstelling van de menselijke gedaante spreekt de hele inhoud van het werk tegelijkertijd tot ons als een uitdrukking van zijn object.

De aard van alle kunst, zoals Hegel het formuleerde, is dat deze ``de mens met zichzelf confronteert''. Andere natuurlijke objecten — niet alleen de menselijke gedaante — kunnen morele ideeën uitdrukken in een artistieke voorstelling. Alle artistieke voorstellingen, of het nu gaat om landschappen, stillevens of zelfs een inspirerend uitzicht op de natuur, bereiken dit. Hier heeft Kant echter gelijk: de uitdrukking van morele waarde is dan geleend.

Maar de mens drukt deze ideeën uit in zijn eigen wezen, en omdat hij is wat hij is. Een boom die door ongunstige groeiomstandigheden gekrompen is, kan ons ellendig lijken, maar de boom voelt zich niet ellendig of drukt deze ellende uit, en vanuit het oogpunt van het ideaal van de boom is gekrompen zijn geen ``ellende''. De ellendige mens is echter ellendig, gemeten aan het menselijke morele ideaal zelf (en niet alleen omdat we van hem eisen dat hij zich onderwerpt aan een menselijk ideaal dat eenvoudigweg niet voor hem geldt, waaraan gemeten hij ellende voor ons zou uitdrukken zonder zelf ellendig te zijn). Hegel begreep dat perfect in zijn colleges over esthetiek toen hij de uitdrukking van het morele beschreef als de ``glans van het geestelijke''.

Zo leidt het formalisme van ``droog genoegen'' tot een beslissende afbraak, niet alleen van het rationalisme in de esthetiek, maar ook van elke algemene (kosmologische) leer van schoonheid. Met precies dat classicistische onderscheid tussen een normatief idee en het ideaal van schoonheid ondergraaft Immanuel Kant de grondslagen van een esthetica van volmaaktheid, die de unieke en onvergelijkbare schoonheid van alles zoekt in haar volledige aanwezigheid voor de zintuigen.
Pas vanaf dit moment kan ``kunst'' een autonoom verschijnsel worden. Haar taak is niet langer om de idealen van de natuur weer te geven, maar om de mens in staat te stellen zichzelf te ontmoeten in de natuur én in de menselijke, historische wereld.
Kants stelling dat het schone behaagt zonder begrip, neemt niet weg dat alleen datgene wat zowel mooi is als betekenisvol tot ons spreekt, onze volle belangstelling wekt. Juist de erkenning dat smaak niet in begrippen te vangen is, voert ons voorbij een esthetica die louter op smaak berust.

\subsubsection{De interesse die door natuurlijke en artistieke schoonheid wordt opgewekt}

Wanneer Kant de vraag stelt naar de interesse die in het schone niet empirisch, maar \textit{a priori} wordt genomen, roept deze vraag naar de interesse in het schone — in tegenstelling tot wat hij stelt over de fundamentele onpartijdigheid van esthetisch genoegen — een nieuw probleem op en voltooit de overgang van het standpunt van smaak naar dat van genie. Het is dezelfde leer die in verband wordt gebracht met beide verschijnselen. Bij het leggen van de grondslagen is het belangrijk om de ``kritiek van de smaak'' te bevrijden van sensualistische en rationalistische vooroordelen. Het is volstrekt passend dat Kant niet onderzoekt hoe het object dat esthetisch wordt beoordeeld, bestaat (en dus de hele kwestie van de relatie tussen de schoonheid van de natuur en die van de kunst). Maar deze dimensie van de vraag wordt noodzakelijkerwijs geopend als men het standpunt van smaak doordenkt — wat betekent dat men daar voorbij gaat.
Het fundamentele probleem dat Kants esthetiek drijft, is dat het schone onze interesse wekt. Het doet dit anders in de natuur dan in de kunst, en de vergelijking tussen natuurlijke en artistieke schoonheid opent deze problematiek.

Hier vinden we Kants meest kenmerkende overtuigingen. In tegenstelling tot wat we misschien zouden verwachten, is het niet omwille van de kunst dat Kant voorbijgaat aan ``onpartijdig genoegen'' en de interesse in het schone onderzoekt. Uit de leer van het ideaal van schoonheid leidden we een voordeel van kunst boven natuurlijke schoonheid af: het voordeel van een directere uitdrukking van het morele. Kant benadrukt echter vooral (§42) het voordeel van natuurlijke boven artistieke schoonheid. Het is niet alleen voor het zuivere esthetische oordeel dat natuurlijke schoonheid een voordeel heeft, namelijk om duidelijk te maken dat het schone afhangt van de doelmatigheid (\textit{Zweckmäßigkeit}) van het voorgestelde voor ons kenvermogen. Dit is zo duidelijk het geval bij natuurlijke schoonheid omdat deze geen inhoudelijke betekenis bezit en daardoor het smaakoordeel in zijn niet-geïntellectualiseerde zuiverheid manifesteert.

Maar het heeft niet alleen dit methodologische voordeel; volgens Kant heeft het ook een inhoudelijk voordeel, en hij hecht blijkbaar veel waarde aan dit punt van zijn leer. Mooie natuur is in staat om een directe interesse op te wekken, namelijk een morele. Wanneer we de mooie vormen van de natuur mooi vinden, wijst deze ontdekking voorbij zichzelf naar de gedachte ``dat de natuur die schoonheid heeft voortgebracht''. Waar deze gedachte interesse wekt, is er cultivering van het morele gevoel. Hoewel Kant, geïnspireerd door Rousseau, weigert om een algemene redenering te maken van de verfijning van smaak voor het schone naar moreel gevoel, is het gevoel voor de schoonheid van de natuur voor Kant een bijzonder geval. Dat de natuur mooi is, wekt alleen interesse op bij iemand die ``zijn interesse al diep in het moreel goede heeft gevestigd''. Vandaar dat de interesse in natuurlijke schoonheid ``verwant is aan het morele''. Door de onbedoelde harmonie van de natuur met ons volledig onpartijdige genoegen te observeren — dat wil zeggen, de wonderbaarlijke doelmatigheid (\textit{Zweckmäßigkeit}) van de natuur voor ons — wijst deze naar ons als het uiteindelijke doel van de schepping, naar ``de morele kant van ons bestaan''.

Hier past de afwijzing van de volmaaktheidsesthetiek perfect bij de morele betekenis van natuurlijke schoonheid. Juist omdat we in de natuur geen doelen op zich vinden en toch schoonheid aantreffen — dat wil zeggen, een doelmatigheid (\textit{Zweckmäßigkeit}) voor het doel (\textit{Zweck}) van ons genoegen — geeft de natuur ons een ``hint'' dat wij inderdaad het uiteindelijke doel, het einddoel van de schepping zijn. De ontbinding van het oude kosmologische denken, dat de mens zijn plaats toekende in de totale structuur van het zijn en elk wezen zijn doel van volmaaktheid toekende, geeft de wereld, die ophoudt mooi te zijn als een structuur van absolute doelen, een nieuwe schoonheid: die van doelmatig te zijn voor ons. Het wordt ``natuur'', wier onschuld erin bestaat dat deze niets weet van de mens of zijn sociale ondeugden. Toch heeft het ons iets te zeggen. Als mooi vindt de natuur een taal die ons een begrijpelijke idee brengt van wat de mensheid zou moeten zijn.

Natuurlijk hangt de betekenis van kunst ook af van het feit dat deze tot ons spreekt, dat deze de mens confronteert met zichzelf in zijn moreel bepaalde bestaan. Maar de producten van kunst bestaan alleen om ons op deze manier aan te spreken — natuurlijke objecten echter bestaan niet om ons op deze manier aan te spreken. Dit is de betekenisvolle interesse van het natuurlijk mooie: dat het de mens nog steeds in staat stelt om zichzelf te ontmoeten in zijn moreel bepaalde bestaan. Kunst kan ons deze zelfontdekking van de mens in een werkelijkheid die niet de bedoeling heeft om dit te doen, niet meedelen. Wanneer de mens zichzelf in kunst tegenkomt, is dit niet de bevestiging van zichzelf door een ander.

Dat klopt, tot op zekere hoogte. De overtuigingskracht van Kants argument is indrukwekkend, maar hij past niet de juiste criteria toe voor het verschijnsel van kunst. Men kan een tegenargument maken. Het voordeel van natuurlijke schoonheid boven artistieke schoonheid is slechts de andere kant van het onvermogen van natuurlijke schoonheid om iets specifieks uit te drukken. Omgekeerd kan men zien dat het voordeel van kunst boven natuurlijke schoonheid is dat de taal van kunst haar eisen doet gelden, en zich niet vrij en onbepaald aanbiedt voor interpretatie naar believen, maar op een betekenisvolle en bepaalde manier tot ons spreekt. En het wonderlijke en mysterieuze van kunst is dat deze bepaaldheid geenszins een belemmering is voor onze geest, maar juist ruimte opent voor spel, voor het vrije spel van onze kenvermogens. Kant heeft gelijk wanneer hij zegt dat kunst in staat moet zijn om ``als natuur'' te worden beschouwd — dat wil zeggen, te behagen zonder de dwang van regels te verraden. We meten niet de bedoelde overeenstemming tussen wat wordt voorgesteld en de werkelijkheid die we kennen, we kijken niet wat het lijkt, we meten de claim op betekenis niet aan een criterium dat we al goed kennen, maar integendeel, dit criterium — ``het begrip'' — wordt op een onbeperkte manier ``esthetisch uitgebreid''.

Kants definitie van kunst als de ``mooie voorstelling van een ding'' houdt hier rekening mee, voor zover zelfs het lelijke mooi is in artistieke voorstelling. Toch komt de eigenlijke aard van kunst slecht naar voren uit de tegenstelling met natuurlijke schoonheid. Als het idee van een ding alleen maar op een mooie manier werd voorgesteld, dan zou dat een louter ``academische'' voorstelling zijn, en zou het slechts aan de minimale eis van alle schoonheid voldoen. Maar voor Kant is kunst meer dan de ``mooie voorstelling van een ding'': het is de presentatie van esthetische ideeën — dat wil zeggen, van iets dat voorbij alle begrippen ligt. Het begrip genie probeert dit inzicht van Kant onder woorden te brengen.

Het kan niet worden ontkend dat de leer van esthetische ideeën, door wier voorstelling de kunstenaar het gegeven begrip oneindig uitbreidt en het vrije spel van de geestelijke vermogens aanmoedigt, voor een moderne lezer iets onbevredigends heeft. Het lijkt alsof deze ideeën worden verbonden met het al dominante begrip, zoals de attributen van een godheid met zijn vorm. De traditionele superioriteit van het rationele begrip boven de onuitsprekelijke esthetische voorstelling is zo sterk dat zelfs bij Kant de schijn wordt gewekt alsof het begrip voorgaat aan de esthetische idee, terwijl het helemaal niet het verstand is, maar de verbeelding die de leiding neemt onder de vermogens in spel. De estheticus zal veel andere uitspraken vinden in het licht waarvan het voor Kant moeilijk is, zonder de superioriteit van het begrip te claimen, vast te houden aan zijn leidende inzicht dat het schone zonder begrip wordt begrepen en toch tegelijkertijd een bindende kracht heeft.

Maar de hoofdlijnen van zijn denken zijn vrij van deze fouten en vertonen een indrukwekkende logische consistentie, die zijn hoogtepunt bereikt in zijn uiteenzetting over genie als grondslag van kunst. Zelfs zonder in te gaan op een gedetailleerdere interpretatie van deze ``bekwaamheid om esthetische ideeën voor te stellen'', kan worden opgemerkt dat Kant hier niet wordt afgeleid van transcendentale vraagstelling en in een doodlopende steeg van een psychologie van artistieke schepping wordt geduwd. Integendeel, de irrationaliteit van genie benadrukt één element in de scheppende productie van regels dat zowel bij schepper als ontvanger duidelijk is, namelijk dat er geen andere manier is om de inhoud van een kunstwerk te begrijpen dan door de unieke vorm van het werk en in het mysterie van zijn indruk, die nooit volledig in woorden kan worden uitgedrukt. Vandaar dat het begrip genie overeenkomt met wat Kant ziet als het cruciale aan esthetische smaak, namelijk dat deze het spel van de geestelijke vermogens bevordert, de levendigheid vergroot die voortkomt uit de harmonie tussen verbeelding en verstand, en uitnodigt om te vertoeven bij het schone. Genie is uiteindelijk een uiting van deze levendige geest, want in tegenstelling tot de starre navolging van regels door de pedant, toont genie een vrije vlucht van uitvinding en daardoor de originaliteit die nieuwe modellen schept.

\subsubsection{De relatie tussen smaak en genie}

In deze context rijst de vraag hoe Kant de wederzijdse relatie tussen smaak en genie ziet. Kant behoudt de bevoorrechte positie van smaak, voor zover kunstwerken (dat wil zeggen, de kunst van genie) vanaf het leidende standpunt van schoonheid moeten worden beoordeeld. Men kan het spijtig vinden dat smaak verbeteringen oplegt aan de uitvindingen van genie, maar smaak is een noodzakelijke discipline voor genie. In gevallen van conflict vindt Kant dan ook dat smaak moet zegevieren. Maar dit is geen belangrijke kwestie, want in wezen delen smaak en genie een gemeenschappelijke basis. De kunst van genie dient ertoe het vrije spel van de geestelijke vermogens mededeelbaar te maken. Dit wordt bereikt door de esthetische ideeën die het uitvindt. Maar ook het esthetische genoegen van smaak werd gekenmerkt door de mededeelbaarheid van een gemoedstoestand — genoegen. Smaak is een vermogen van oordeel, en derhalve reflecterend, maar waar het over reflecteert, is slechts die gemoedstoestand: de levendigheid van de kenvermogens die zowel door natuurlijke als door artistieke schoonheid ontstaat. Zo is de systematische betekenis van het begrip genie beperkt doordat het een bijzonder geval is van het artistiek mooie, terwijl het begrip smaak juist universeel is.

Dat Kant het begrip genie volledig in dienst stelt van zijn transcendentale onderzoek en niet afzakt naar empirische psychologie, blijkt duidelijk uit zijn beperking van het begrip genie tot artistieke schepping. Wanneer hij deze naam onthoudt aan grote uitvinders en onderzoekers op het gebied van wetenschap en technologie, is dit, vanuit het oogpunt van empirische psychologie, volledig ongerechtvaardigd. Waar men iets moet ``ontdekken'' dat niet louter door leren en methodisch werk kan worden gevonden — dat wil zeggen, waar er sprake is van \textit{inventio}, waar iets te danken is aan inspiratie en niet aan methodische berekening — daar is \textit{ingenium}, genie, van belang. Toch is Kants bedoeling juist: alleen het kunstwerk is van binnenuit zo bepaald dat het slechts door genie kan worden geschapen. Alleen in het geval van de kunstenaar blijft zijn ``uitvinding'' — het werk — van nature verbonden met de geest: de geest die schept, maar ook die oordeelt en geniet. Alleen dergelijke uitvindingen kunnen niet worden nageaapt, en daarom is het — vanuit transcendentale invalshoek — terecht dat Kant (alleen hier) spreekt van genie en kunst definieert als de kunst van genie. Alle andere prestaties en uitvindingen van genie, hoe geniaal ze ook mogen zijn, worden in hun wezen niet door genie bepaald.

Ik ben van mening dat voor Kant het begrip genie in feite slechts een aanvulling was op wat hem ``om transcendentale redenen'' in het esthetische oordeel interesseerde. We mogen niet vergeten dat het tweede deel van de \textit{Kritiek van het oordeelsvermogen} uitsluitend gaat over de natuur (en over het beoordelen daarvan aan de hand van doeleinden) en helemaal niet over kunst. Voor de systematische bedoeling van het geheel is het toepassen van het esthetische oordeel op het schone en het verhevene in de natuur dan ook belangrijker dan de transcendentale fundering van kunst. De ``doelmatigheid van de natuur voor onze kenvermogens'' — die als transcendentale grond van het esthetische oordeel alleen op natuurlijke schoonheid (en niet op kunst) betrekking heeft — fungeert tegelijkertijd als voorbereiding voor het verstand om het begrip doel toe te passen op de natuur. Dit is Kants filosofische bedoeling: het legitimeren van teleologie, wier constitutieve claim als principe van oordeel in de kennis van de natuur door de \textit{Kritiek van de zuivere rede} was vernietigd. Deze bedoeling brengt zijn hele filosofie tot een systematisch besluit. Het oordeel vormt de brug tussen verstand en rede. Het intelligibele waar smaak naar wijst, het bovenzintuiglijke substraat in de mens, bevat tegelijkertijd de bemiddeling tussen de begrippen van natuur en vrijheid. Dit is de systematische betekenis die het probleem van natuurlijke schoonheid voor Kant heeft: het grondt de centrale positie van teleologie. Alleen natuurlijke schoonheid, niet kunst, kan helpen om het begrip doel in het oordeel over de natuur te legitimeren. Om deze systematische reden alleen al biedt het ``zuivere'' smaakoordeel de onmisbare basis voor de derde \textit{Kritiek}.

Maar zelfs binnen de ``kritiek van het esthetische oordeel'' is het duidelijk dat het standpunt van genie uiteindelijk dat van smaak verdringt. Men hoeft slechts te kijken naar de manier waarop Kant genie beschrijft: de genie is een gunsteling van de natuur — net zoals natuurlijke schoonheid wordt gezien als een gunst van de natuur. We moeten kunst kunnen beschouwen alsof het natuur is. Door genie geeft de natuur de kunst haar regels. In al deze uitspraken is het begrip natuur het onbetwiste criterium. Wat het begrip genie bereikt, is slechts dat de producten van kunst esthetisch op één lijn worden gesteld met natuurlijke schoonheid. Ook kunst wordt esthetisch beoordeeld — dat wil zeggen, ook kunst vraagt om een reflecterend oordeel. Wat met opzet is geproduceerd, en derhalve doelmatig is, mag niet in verband worden gebracht met een begrip, maar zoekt erkenning louter door beoordeeld te worden — net als natuurlijke schoonheid. ``Kunst is kunst geschapen door genie'' betekent dat ook voor artistieke schoonheid er geen ander principe van oordeel is, geen criterium van begrip en kennis, dan dat van de geschiktheid om het gevoel van vrijheid in het spel van onze kenvermogens te bevorderen. Of het nu in de natuur of in de kunst is, schoonheid heeft hetzelfde \textit{a priori} principe, dat volledig binnen de subjectiviteit ligt. De autonomie van het esthetische oordeel betekent niet dat er een autonoom geldigheidsdomein is voor mooie objecten. Kants transcendentale reflectie over het \textit{a priori} van het oordeel rechtvaardigt de claim van het esthetische oordeel, maar in wezen staat het een filosofische esthetiek in de zin van een filosofie van de kunst niet toe (Kant zegt zelf dat deze \textit{Kritiek} niet bedoeld is als een leer of als een metafysisch systeem.)

\subsection{De esthetiek van het genie en het begrip ervaring \textit{Erlebnis}}

\subsubsection{De dominantie van het begrip genie}

Het baseren van het esthetische oordeel op het \textit{a priori} van de subjectiviteit zou een heel nieuwe betekenis krijgen toen de implicaties van de transcendentale filosofische reflectie veranderden met Kants opvolgers. Als de metafysische achtergrond — die ten grondslag ligt aan de voorrang van natuurlijke schoonheid bij Kant en die het begrip genie terugvoert naar de natuur — niet meer bestaat, rijst het probleem van de kunst op een nieuwe manier. Zelfs de wijze waarop Schiller Kants \textit{Kritiek van het oordeelsvermogen} oppakte en al zijn morele en pedagogische inzet achter het idee van een ``esthetische opvoeding'' zette, gaf niet het standpunt van smaak en oordeel (zoals bij Kant), maar dat van de kunst een centrale plaats.

Vanuit het standpunt van de kunst wisselen de Kantiaanse ideeën van smaak en genie volledig van plaats. Genie moest het meer omvattende begrip worden, en omgekeerd moest het verschijnsel smaak worden gedevalueerd. Toch zijn er bij Kant zelf al aanknopingspunten voor zo’n omkering van waarden. Ook volgens Kant is het voor het oordelend vermogen van smaak van enig belang dat kunst het werk is van genie. Een van de dingen die smaak beoordeelt, is of een kunstwerk geest heeft of geesteloos is. Kant zegt over artistieke schoonheid dat ``men bij het beoordelen van zo’n object rekening moet houden met de mogelijkheid van geest — en dus van genie — daarin'', en op een andere plaats maakt hij de voor de hand liggende opmerking dat zonder genie niet alleen kunst, maar ook een juiste, onafhankelijke smaak bij het beoordelen daarvan onmogelijk is. Daardoor gaat het standpunt van smaak, voor zover het wordt toegepast op zijn belangrijkste object, kunst, onvermijdelijk over in het standpunt van genie. Genie in het begrip correspondeert met genie in de schepping. Kant drukt het niet zo uit, maar het begrip geest dat hij hier gebruikt, is in beide gevallen evenzeer van toepassing. Hierop moet later verder worden voortgebouwd.

Het is inderdaad duidelijk dat het begrip smaak zijn betekenis verliest wanneer het verschijnsel kunst op de voorgrond treedt. Het standpunt van smaak is ondergeschikt aan het kunstwerk. De gevoeligheid bij het selecteren die smaak kenmerkt, heeft vaak een nivellerend effect in tegenstelling tot de originaliteit van het artistieke werk van genie. Smaak mijdt het ongebruikelijke en het monsterlijke. Het heeft betrekking op het oppervlak der dingen; het bekommert zich niet om wat origineel is aan een artistieke schepping. Al in de beginjaren van het idee van genie in de achttiende eeuw vinden we een polemische toon tegen het begrip smaak. Deze was gericht tegen de classicistische esthetiek en eiste dat het ideaal van smaak van het Franse classicisme plaats moest maken voor Shakespeare (Lessing). In die zin is Kant ouderwets en neemt hij een tussenpositie in, doordat hij om transcendentale redenen het begrip smaak hardnekkig handhaafde, terwijl de \textit{Sturm und Drang} dit niet alleen fel verwierp, maar ook fel afbrak.

Maar wanneer Kant overgaat van het leggen van algemene grondslagen naar de specifieke problemen van de filosofie van de kunst, wijst hij zelf voorbij het standpunt van smaak en spreekt hij over een volmaaktheid van smaak. Maar wat is dat? Het normatieve karakter van smaak impliceert de mogelijkheid dat deze kan worden gekweekt en volmaakt. Volmaakte smaak, die het belangrijk is te bereiken, zal volgens Kant een bepaalde onveranderlijke vorm aannemen. Dat is heel logisch, hoe absurd het ook in onze oren mag klinken. Want als smaak goede smaak moet zijn, dan doet dit de hele relativiteit van smaak, die door esthetisch scepticisme wordt verondersteld, teniet. Het zou alle kunstwerken omvatten die ``kwaliteit'' hebben, en dus uiteraard alle werken die door genie zijn geschapen.

We zien dus dat het idee van volmaakte smaak waar Kant over spreekt, beter gedefinieerd zou kunnen worden door het begrip genie. Het is duidelijk onmogelijk om het idee van volmaakte smaak toe te passen binnen de sfeer van natuurlijke schoonheid. Het zou misschien acceptabel zijn in het geval van tuinieren; maar in overeenstemming met zijn betoog, wijst Kant tuinieren toe aan de sfeer van het artistiek mooie. Maar geconfronteerd met natuurlijke schoonheid — bij voorbeeld de schoonheid van een landschap — is het idee van een volmaakte smaak volledig misplaatst. Zou het erin bestaan om elke natuurlijke schoonheid naar verdienste te beoordelen? Kan er in deze sfeer keuze zijn? Is er een hiërarchie van verdienste? Is een zonnig landschap mooier dan een landschap in de regen? Is er iets lelijks in de natuur? Of is er alleen maar verschillend aantrekkelijks in verschillende stemmingen, verschillend bevalligs voor verschillende smaken? Kant heeft misschien gelijk wanneer hij het moreel betekenisvol vindt dat iemand genoegen kan scheppen in de natuur. Maar is het zinvol om een onderscheid te maken tussen goede en slechte smaak in relatie daartoe? Waar dit onderscheid onbetwistbaar van toepassing is — namelijk in relatie tot kunst en kunstmatigheid — is smaak, zoals we hebben gezien, slechts een beperking van het schone en bevat het geen eigen principe. Het idee van een volmaakte smaak is dus twijfelachtig, zowel in relatie tot de natuur als tot de kunst. Men doet geweld aan het begrip smaak als men de veranderlijkheid ervan niet accepteert. Smaak is, als het al iets is, een getuigenis van de veranderlijkheid van alle menselijke dingen en de relativiteit van alle menselijke waarden.

Kants fundering van de esthetiek op het begrip smaak is niet geheel bevredigend. Het begrip genie, dat Kant ontwikkelt als een transcendentale grond voor artistieke schoonheid, lijkt veel beter geschikt om een universeel esthetisch principe te zijn. Het voldoet immers veel beter dan het begrip smaak aan de eis om onveranderlijk te zijn in de stroom der tijd. Het wonder van de kunst — die raadselachtige volmaaktheid die succesvolle artistieke scheppingen bezitten — is in alle tijden zichtbaar. Het lijkt mogelijk om smaak ondergeschikt te maken aan de transcendentale verklaring van kunst en onder smaak het zekere gevoel voor genie in de kunst te verstaan. Kants uitspraak ``Schone kunst is de kunst van genie'' wordt dan een transcendentale grond voor de esthetiek in het algemeen. Esthetiek is uiteindelijk alleen mogelijk als filosofie van de kunst.

Het Duitse idealisme trok deze conclusie. In navolging van Kants leer van de transcendentale verbeelding maakten Fichte en Schelling in hun esthetiek op deze en andere punten nieuw gebruik van dit idee. In tegenstelling tot Kant zagen zij het standpunt van de kunst (als de onbewuste schepping van genie) als alomvattend — inclusief de natuur, die wordt begrepen als een product van de geest.

Maar nu is de basis van de esthetiek verschoven. Net als het begrip smaak is ook het begrip natuurlijke schoonheid gedevalueerd of anders geïnterpreteerd. Het morele belang van natuurlijke schoonheid, dat Kant zo enthousiast had geschetst, trekt zich nu terug achter de zelfontmoeting van de mens in kunstwerken. In Hegels meesterwerk \textit{Esthetiek} bestaat natuurlijke schoonheid alleen als een ``weerspiegeling van de geest''. Er is inderdaad geen onafhankelijk element meer in het systematische geheel van de esthetiek.

De onbepaaldheid waarmee natuurlijke schoonheid zich aan de interpreterende en begrijpende geest presenteert, rechtvaardigt het dat we met Hegel zeggen dat ``haar substantie [in de geest] is bevat''. Esthetisch gesproken trekt Hegel hier een volstrekt correcte conclusie; ik ben daar al op ingegaan toen ik sprak over de ongeschiktheid van het toepassen van het idee van smaak op de natuur. Want oordelen over de schoonheid van een landschap hangen ongetwijfeld af van de artistieke smaak van de tijd. Men hoeft alleen maar te denken aan het Alpenlandschap dat in de achttiende eeuw nog als lelijk werd beschreven — een effect, zoals we weten, van de geest van kunstmatige symmetrie die de eeuw van het absolutisme domineert. Zo is Hegels esthetiek stevig geworteld in het standpunt van de kunst. In de kunst ontmoet de mens zichzelf, geest ontmoet geest.

Beslissend voor de ontwikkeling van de moderne esthetiek is dat ook hier, net als in de hele sfeer van de systematische filosofie, het speculatieve idealisme een effect had dat ver uitstijgt boven zijn erkende betekenis. De felle afwijzing van het dogmatische schematisme van de Hegeliaanse school in het midden van de negentiende eeuw leidde tot de eis van een vernieuwing van de kritiek onder de vlag ``terug naar Kant''. Hetzelfde gold voor de esthetiek. Hoe briljant kunst ook werd gebruikt voor het schrijven van de geschiedenis van wereldbeelden, zoals in Hegels \textit{Esthetiek}, deze methode van \textit{a priori} geschiedschrijving, die vaak door de Hegeliaanse school (Rosenkranz, Schasler, etc.) werd toegepast, raakte snel in diskrediet. De oproep om terug te keren naar Kant, die hiertegen in opstand kwam, kon nu echter geen echte terugkeer en herstel van de horizon van Kants kritieken zijn. Integendeel, het verschijnsel kunst en het begrip genie bleven in het centrum van de esthetiek staan; het probleem van natuurlijke schoonheid en het begrip smaak werden naar de rand gedrongen.

Dit blijkt ook uit het taalgebruik. Kants beperking van het begrip genie tot de kunstenaar (die ik hierboven heb onderzocht) heeft niet standgehouden; integendeel, in de negentiende eeuw steeg het begrip genie op tot de status van een universeel waardebegrip en bereikte het — samen met het begrip van het scheppende — een ware apotheose, mede dankzij de romantische en idealistische opvatting van onbewuste schepping en door Schopenhauers filosofie van het onbewuste, die enorme populariteit verwierf. Ik heb laten zien dat deze soort systematische dominantie van het begrip genie boven het begrip smaak niet Kantiaans is. Kants hoofdzaak was echter om de esthetiek een autonome basis te geven, bevrijd van het criterium van het begrip, en niet om de vraag naar waarheid in de sfeer van de kunst op te werpen, maar om het esthetische oordeel te baseren op het subjectieve \textit{a priori} van ons \textit{Le\-bensgefühl}, de harmonie van ons vermogen tot ``kennis in het algemeen'', die de kern vormt van zowel smaak als genie. Al dit alles paste naadloos bij de negentiende-eeuwse irrationalisme en de verering van genie. Kants leer van de ``verhoging van het \textit{Lebensgefühl}'' in het esthetische genoegen hielp het idee van ``genie'' om zich te ontwikkelen tot een omvattend levensbegrip (\textit{Leben}), met name nadat Fichte genie en wat genie schept tot een universele transcendentale positie had verheven. Daardoor verklaarde het neokantianisme, door alle objectieve geldigheid af te leiden uit transcendentale subjectiviteit, het begrip \textit{Erlebnis} tot de eigenlijke stof van het bewustzijn.

\subsubsection{Over de geschiedenis van het woord \textit{Erlebnis}}

Het is verrassend om te constateren dat, in tegenstelling tot het werkwoord \textit{erleben}, het zelfstandig naamwoord \textit{Erlebnis} pas in de jaren 1870 algemene bekendheid verwierf. In de achttiende eeuw komt het woord helemaal niet voor, en zelfs Schiller en Goethe kennen het niet. De eerste vermelding schijnt voor te komen in een van Hegels brieven. Maar ook in de jaren dertig en veertig van de negentiende eeuw ken ik slechts sporadische gevallen (bij Tieck, Alexis en Gutzkow). Het woord komt in de jaren vijftig en zestig eveneens zelden voor en duikt plotseling vaker op in de jaren zeventig. Kennelijk dringt het woord tegelijkertijd door in het algemene taalgebruik als het in biografische geschriften gaat worden gebruikt.

Aangezien \textit{Erlebnis} een afleiding is van het oudere werkwoord \textit{erleben}, dat in de tijd van Goethe vaak voorkomt, moeten we de betekenis van \textit{erleben} analyseren om te begrijpen waarom het nieuwe woord is gevormd. \textit{Erleben} betekent in de eerste plaats ``nog in leven zijn wanneer iets gebeurt''. Het woord suggereert dus de onmiddellijkheid waarmee iets reëels wordt begrepen — in tegenstelling tot iets wat men veronderstelt te weten, maar wat niet door eigen ervaring is bevestigd, of het nu is overgenomen van anderen, afkomstig is uit mondelinge overlevering, of is afgeleid, verondersteld of verbeeld. Wat ervaren wordt, is altijd wat men zelf heeft meegemaakt.

Tegelijkertijd wordt de vorm ``das Erlebte'' gebruikt om de blijvende inhoud van wat wordt ervaren aan te duiden. Deze inhoud is als een opbrengst of resultaat dat uit de vluchtigheid van het ervaren blijvende gewicht en betekenis verkrijgt. Beide betekenissen liggen kennelijk ten grondslag aan de vorming van \textit{Erlebnis}: zowel de onmiddellijkheid, die aan alle interpretatie, bewerking en mededeling voorafgaat en slechts een uitgangspunt voor interpretatie biedt — stof die gevormd moet worden — als de ontdekte opbrengst ervan, het blijvende resultaat.

Overeenkomstig de dubbele betekenis van het werkwoord \textit{erleben} is het feit dat het woord \textit{Erlebnis} wortel schiet via de biografische literatuur. De kern van een biografie, met name negentiende-eeuwse biografieën van kunstenaars en dichters, is om de werken te begrijpen aan de hand van het leven. Hun prestatie bestaat juist in het bemiddelen tussen de twee betekenissen die we in het woord \textit{Erlebnis} hebben onderscheiden en in het zien van deze betekenissen als een productieve eenheid: iets wordt een ``ervaring'' niet alleen voor zover het wordt meegemaakt, maar voor zover het meemaken een bijzondere indruk achterlaat die het een blijvende betekenis geeft. Een ``ervaring'' van deze aard verkrijgt een geheel nieuwe status wanneer deze in kunst tot uiting wordt gebracht. Diltheys bekende titel \textit{Das Erlebnis und die Dichtung} (\textit{Ervaring en poëzie}) vat deze associatie bondig samen. Inderdaad was Dilthey de eerste die het woord een conceptuele functie gaf die al snel zo populair werd en een zo vanzelfsprekend waardebegrip aanduidde, dat veel Europese talen het als leenwoord overnamen. Maar het is redelijk om aan te nemen dat Diltheys gebruik van de term slechts benadrukte wat zich in werkelijkheid in het taalgebruik had voorgedaan.

Bij Dilthey kunnen we gemakkelijk de diverse elementen isoleren die in het taalkundig en conceptueel nieuwe woord \textit{Erlebnis} werkzaam zijn. De titel \textit{Das Erlebnis und die Dichtung} is laat genoeg (1905). De eerste versie van het essay over Goethe dat daarin is opgenomen, en dat Dilthey in 1877 publiceerde, gebruikt het woord \textit{Erlebnis} tot op zekere hoogte, maar vertoont niets van de latere terminologische precisie van het begrip. De vroege vormen van de later conceptueel vastgestelde betekenis van \textit{Erlebnis} verdienen nader onderzoek. Het lijkt meer dan toeval dat het woord plotseling vaker opduikt in een biografie van Goethe (en in een essay over dat onderwerp). Goethe nodigt meer dan wie ook uit om dit woord te munten, omdat zijn poëzie op een heel nieuwe manier begrijpelijk wordt door wat hij heeft meegemaakt. Hij zei zelf dat al zijn poëzie het karakter had van een grote belijdenis. Hermann Grimms biografie van Goethe neemt deze uitspraak als methodologisch uitgangspunt en gebruikt daardoor vaak het meervoud \textit{Erlebnisse}.

Diltheys essay over Goethe laat ons een blik werpen op de onbewuste voorgeschiedenis van het woord, omdat dit essay voorafgaat aan de versie van 1877 en de latere bewerking in \textit{Das Erlebnis und die Dichtung} (1905). In dit essay vergelijkt Dilthey Goethe met Rousseau, en om het nieuwe soort schrijven te beschrijven dat Rousseau baseerde op de wereld van zijn innerlijke ervaringen, gebruikt hij de uitdrukking \textit{das Erleben}. In zijn paraphrase van Rousseau vinden we ook de uitdrukking ``die Erlebnisse früherer Tage'' (de ervaringen van vroegere dagen).

Toch is zelfs bij de vroege Dilthey de betekenis van het woord \textit{Erlebnis} nog tamelijk onzeker. Dit blijkt duidelijk uit een passage waaruit Dilthey in latere edities het woord \textit{Erlebnis} heeft geschrapt: ``Overeenkomstig zowel met wat hij had meegemaakt als met wat hij, gelet op zijn onwetendheid over de wereld, zich had voorgesteld en als ervaring (\textit{Erlebnis}) had behandeld...'' Opnieuw heeft hij het over Rousseau. Maar een verbeelde ervaring past niet bij de oorspronkelijke betekenis van \textit{erleben}, en zelfs niet bij Diltheys eigen latere technische gebruik, waar \textit{Erlebnis} betekent wat onmiddellijk gegeven is, de uiteindelijke stof voor alle scheppende verbeelding. Het gevormde woord \textit{Erlebnis} drukt natuurlijk kritiek uit op het rationalisme van de Verlichting, dat, in navolging van Rousseau, het begrip leven (\textit{Leben}) benadrukte. Waarschijnlijk was het Rousseau’s invloed op het Duitse classicisme die het criterium van \textit{Erlebtsein} (meegemaakt zijn) introduceerde en daardoor de vorming van het woord \textit{Erlebnis} mogelijk maakte. Maar het begrip leven vormt ook de metafysische achtergrond voor het Duitse speculatieve idealisme en speelt een fundamentele rol bij Fichte, Hegel en zelfs Schleiermacher. In tegenstelling tot de abstractie van het verstand en de particulariteit van de waarneming of voorstelling, impliceert dit begrip een verbinding met totaliteit, met oneindigheid. Dit is duidelijk hoorbaar in de klank die het woord \textit{Erlebnis} zelfs vandaag de dag nog heeft.

Schleiermachers oproep tot levend gevoel tegen het koude rationalisme van de Verlichting, Schillers oproep tot esthetische vrijheid tegen de mechanistische samenleving, Hegels tegenstelling tussen leven (later: geest) en ``positiviteit'', waren de voorlopers van het protest tegen de moderne industriële samenleving, dat aan het begin van onze eeuw ertoe leidde dat de woorden \textit{Erlebnis} en \textit{Erleben} bijna heilige strijdkreten werden. De opstand van de \textit{Jugendbewegung} (Jeugdbeweging) tegen de burgerlijke cultuur en haar instellingen werd geïnspireerd door deze ideeën; de invloed van Friedrich Nietzsche en Henri Bergson speelde hierin een rol, maar ook een ``geestelijke beweging'' als die rond Stefan George, en niet in de laatste plaats de seismografische nauwkeurigheid waarmee de filosofie van Georg Simmel op deze gebeurtenissen reageerde, maken allemaal deel uit van hetzelfde verschijnsel. De levensfilosofie van onze eigen tijd bouwt voort op haar romantische voorgangers. De afwijzing van de mechanisering van het leven in de hedendaagse massasamenleving maakt het woord zo vanzelfsprekend dat de conceptuele implicaties ervan volledig verborgen blijven.

We moeten Diltheys munting van het begrip dan ook begrijpen in het licht van de eerdere geschiedenis van het woord bij de romantici, en niet vergeten dat Dilthey de biograaf van Schleiermacher was. Het is waar dat we het woord \textit{Erlebnis} bij Schleiermacher niet aantreffen, en kennelijk ook niet eens het werkwoord \textit{erleben}. Maar er is geen gebrek aan synoniemen die het betekenisbereik van \textit{Erlebnis} dekken, en de pantheïstische achtergrond is altijd duidelijk aanwezig. Elke daad, als element van het leven, blijft verbonden met de oneindigheid van het leven die zich daarin manifesteert. Al het eindige is een uitdrukking, een representatie van het oneindige.

Inderdaad vinden we in Diltheys biografie van Schleiermacher, in de beschrijving van religieuze contemplatie, een bijzonder treffend gebruik van het woord \textit{Erlebnis}, dat al zijn conceptuele inhoud aankondigt: ``Elk van zijn ervaringen (\textit{Erlebnisse}), die op zichzelf bestaan, is een afzonderlijk beeld van het universum, losgemaakt uit de verklarende context.''

\subsubsection{Het begrip \textit{Erlebnis}}

Nu we de geschiedenis van het woord hebben bestudeerd, laten we de geschiedenis van het begrip \textit{Erlebnis} onderzoeken. Uit het voorgaande weten we dat Diltheys concept van \textit{Erlebnis} duidelijk twee elementen bevat: het pantheïstische en het positivistische, namelijk de ervaring (\textit{Erlebnis}) en nog meer het resultaat ervan (\textit{Erlebnis}). Dit is geen toeval, maar het gevolg van zijn eigen tussenpositie tussen speculatie en empirisme, die we later moeten bekijken. Omdat hij erop gericht is de humanistische wetenschappen epistemologisch te legitimeren, wordt hij doorlopend gedomineerd door de vraag wat werkelijk gegeven is. Zijn begrippen worden dus gemotiveerd door dit epistemologische doel, of liever door de behoeften van de epistemologie zelf — behoeften die weerspiegeld worden in het taalkundige proces dat we hierboven hebben geanalyseerd.
Net zoals de afstand tot en het verlangen naar ervaring — veroorzaakt door de ontreddering over de ingewikkelde werking van de door de Industriële Revolutie veranderde beschaving — het woord \textit{Erlebnis} algemeen in gebruik bracht, zo geeft ook de nieuwe, afstandelijke houding die het historisch bewustzijn tegenover de traditie inneemt, het begrip \textit{Erlebnis} zijn epistemologische functie. Kenmerkend voor de ontwikkeling van de humanistische wetenschappen in de negentiende eeuw is dat zij niet alleen de natuurwetenschappen als een extern model erkennen, maar dat zij, komend uit dezelfde achtergrond als de moderne wetenschap, hetzelfde gevoel voor experiment en onderzoek ontwikkelen. Net zoals de tijd van de mechanica zich vervreemd voelde van de natuur, opgevat als de natuurlijke wereld, en deze vervreemding epistemologisch uitdrukte in het begrip zelfbewustzijn en in de regel — ontwikkeld tot een methode — dat alleen ``duidelijke en distincte waarnemingen'' zeker zijn, zo voelden ook de humanistische wetenschappen van de negentiende eeuw een vergelijkbare vervreemding van de wereld van de geschiedenis. De geestelijke scheppingen van het verleden, kunst en geschiedenis, behoren niet langer vanzelfsprekend tot het heden; in plaats daarvan worden zij overgeleverd aan onderzoek, zij zijn gegevens of feiten waaruit een verleden tegenwoordig kan worden gemaakt. Zo is het begrip van het gegeven ook belangrijk in Diltheys formulering van het begrip \textit{Erlebnis}.

Wat Dilthey met het begrip ``ervaring'' probeert te vatten, is de speciale aard van het gegeven in de humanistische wetenschappen. In navolging van Descartes' formulering van de \textit{res cogitans} definieert hij het begrip ervaring door reflexiviteit, door innerlijkheid, en op basis van deze speciale wijze van gegeven zijn probeert hij een epistemologische rechtvaardiging te construeren voor kennis van de historische wereld. De primaire gegevens, waarnaar de interpretatie van historische objecten teruggaat, zijn geen gegevens van experiment en meting, maar eenheden van betekenis. Dat is wat het begrip ervaring stelt: de betekenissenstructuren die we in de humanistische wetenschappen tegenkomen, hoe vreemd en onbegrijpelijk ze ons ook mogen lijken, kunnen worden herleid tot uiteindelijke eenheden van wat in het bewustzijn gegeven is, eenheden die zelf niets vreemds, objectiefs of behulp van interpretatie bevatten. Deze eenheden van ervaring zijn zelf eenheden van betekenis.

We zullen zien hoe cruciaal het is voor Diltheys denken dat de uiteindelijke eenheid van het bewustzijn ``\textit{Erlebnis}'' wordt genoemd en niet ``sensatie'', zoals dat vanzelfsprekend was in het Kantianisme en in de positivistische epistemologie van de negentiende eeuw tot aan Ernst Mach. Dilthey omzeilt zo het ideaal van kennisopbouw uit atomen van sensatie en biedt in plaats daarvan een scherper gedefinieerde versie van het begrip van het gegeven. De eenheid van ervaring (en niet de psychische elementen waaruit deze kan worden geanalyseerd) vertegenwoordigt de ware eenheid van wat gegeven is. Zo vinden we in de epistemologie van de humanistische wetenschappen een levensbegrip dat het mechanistische model beperkt.

Dit levensbegrip is ideologisch bedoeld; leven is voor Dilthey productiviteit. Omdat leven zich objectifieert in betekenissenstructuren, bestaat al het begrip van betekenis uit ``het terugvertalen van de objectificaties van het leven in het geestelijke leven waaruit zij zijn voortgekomen''. Zo is het begrip ervaring de epistemologische basis voor alle kennis van het objectieve.

De epistemologische functie van het begrip ervaring in Husserls fenomenologie is eveneens universeel. In het vijfde deel van de \textit{Logische Onderzoeken} (Hoofdstuk 2) wordt het fenomenologische begrip van ervaring uitdrukkelijk onderscheiden van het volksbegrip. De eenheid van ervaring wordt niet begrepen als een stuk van de werkelijke stroom van ervaringen van een ``ik'', maar als een intentionele relatie. Ook hier is \textit{Erlebnis}, als eenheid van betekenis, teleologisch. Ervaringen bestaan alleen voor zover er iets in wordt ervaren en bedoeld. Het is waar dat Husserl ook niet-intentionele ervaringen erkent, maar deze zijn slechts materiaal voor eenheden van betekenis, intentionele ervaringen. Zo wordt ervaring voor Husserl de omvattende naam voor alle bewustzijnshandelingen wier wezen intentionaliteit is.

Zowel bij Dilthey als bij Husserl, zowel in de levensfilosofie als in de fenomenologie, is het begrip \textit{Erlebnis} in de eerste plaats louter epistemologisch. Zijn teleologische betekenis wordt in aanmerking genomen, maar is niet conceptueel bepaald. Dat leven (\textit{Leben}) zich manifesteert in ervaring (\textit{Erlebnis}) betekent eenvoudigweg dat leven de uiteindelijke grondslag is. De geschiedenis van het woord bood een zekere rechtvaardiging om het op te vatten als een prestatie (\textit{Leistung}). We hebben immers gezien dat de vorming \textit{Erlebnis} een verdichtende, intensiverende betekenis heeft. Als iets een \textit{Erlebnis} wordt genoemd of beschouwd, betekent dat het is afgerond tot de eenheid van een betekenisvol geheel. Een ervaring onderscheidt zich zowel van andere ervaringen — waarin andere dingen worden ervaren — als van de rest van het leven waarin ``niets'' wordt ervaren. Een ervaring is niet langer iets dat snel voorbijstroomt in de stroom van het bewuste leven; het wordt bedoeld als een eenheid en verkrijgt zo een nieuwe wijze van zijn. Het is dan ook heel begrijpelijk dat het woord opduikt in biografische literatuur en uiteindelijk voortkomt uit het gebruik in autobiografieën. Wat een ervaring kan worden genoemd, constituteert zich in de herinnering. Door het zo te noemen, verwijzen we naar de blijvende betekenis die een ervaring heeft voor de persoon die deze heeft. Dit is de reden om te spreken over een intentionele ervaring en de ideologische structuur van het bewustzijn.

Aan de andere kant impliceert het begrip ervaring echter ook een contrast tussen leven en louter begrip. Ervaring heeft een bepaalde onmiddellijkheid die elke mening over de betekenis ervan ontgaat. Alles wat wordt ervaren, wordt door jezelf ervaren, en deel van de betekenis ervan is dat het behoort tot de eenheid van dit zelf en dus een onmiskenbare en onvervangbare relatie tot het geheel van dit ene leven bevat. Wezenlijk voor een ervaring is daarom dat deze niet kan worden uitgeput in wat erover gezegd kan worden of begrepen kan worden als de betekenis ervan. Als deze door autobiografische of biografische reflectie wordt bepaald, blijft de betekenis ervan verbonden met de hele beweging van het leven en begeleidt deze constant. De wijze van zijn van ervaring is juist zo bepalend dat men er nooit mee klaar is. Nietzsche zegt: ``alle ervaringen duren lang bij diepzinnige mensen''. Hij bedoelt dat ze niet snel worden vergeten, dat het lang duurt om ze te verwerken, en dat dit (veel meer dan hun oorspronkelijke inhoud als zodanig) hun specifieke zijn en betekenis uitmaakt. Wat we een \textit{Erlebnis} noemen in deze nadrukkelijke zin betekent dus iets onvergetelijks en onvervangbaars, iets waarvan de betekenis niet kan worden uitgeput door conceptuele bepaling.

Filosofisch gezien betekent de dubbelzinnigheid die we in het begrip \textit{Erlebnis} hebben opgemerkt, dat dit begrip niet volledig wordt uitgeput door zijn zijn als uiteindelijke gegevenheid en grondslag van alle kennis. Er is nog iets heel anders dat in het begrip ``ervaring'' moet worden herkend, en het onthult een reeks problemen die nog moeten worden behandeld: de innerlijke relatie tot het leven.

Er waren twee uitgangspunten voor dit verstrekkende thema — de relatie tussen leven en ervaring — en we zullen hieronder zien hoe Dilthey, en in het bijzonder Husserl, verstrikt raakten in deze reeks problemen. Hier zien we het cruciale belang van Kants kritiek op elke substantiële leer van de ziel en, anders dan deze, het belang van de transcendentale eenheid van het zelfbewustzijn, de synthetische eenheid van de apperceptie. Deze kritiek op de rationalistische psychologie gaf aanleiding tot het idee van een psychologie gebaseerd op Kants kritische methode, zoals Paul Natorp in 1888 ondernam en waar Richard Honigswald later het begrip \textit{Denkpsychologie} op baseerde. Natorp wees \textit{Bewusstheit}, dat de onmiddellijkheid van de ervaring uitdrukt, aan als het object van de kritische psychologie, en ontwikkelde universele subjectivering als de onderzoeksmethode van de reconstructieve psychologie. Natorp steunde en werkte zijn basisidee later verder uit door een grondige kritiek op de begrippen van het tijdgenootschappelijke psychologische onderzoek, maar al in 1888 was het basisidee er: de concretie van de oorspronkelijke ervaring — dat wil zeggen, de totaliteit van het bewustzijn — vertegenwoordigt een ongedifferentieerde eenheid, die wordt gedifferentieerd en bepaald door de objectiverende methode van kennis. ``Maar bewustzijn betekent leven — dat wil zeggen, een ondeelbare onderlinge samenhang.'' Dit zien we met name in de relatie tussen bewustzijn en tijd: ``Bewustzijn is niet gegeven als een gebeurtenis in de tijd, maar tijd als een vorm van bewustzijn.''

In hetzelfde jaar, 1888, waarin Natorp zich zo tegen de heersende psychologie verzette, verscheen Henri Bergsons eerste boek, \textit{Les données immédiates de la conscience}, een kritische aanval op de tijdgenootschappelijke psychofysica, die het idee van leven net zo vastberaden gebruikte als Natorp deed tegen de objectiverende en ruimtelijke neiging van psychologische begrippen. Hier vinden we uitspraken over ``bewustzijn'' en zijn ongedeelde concretie, net als bij Natorp. Bergson bedacht daarvoor de inmiddels beroemde naam \textit{durée}, die de absolute continuïteit van het psychische uitdrukt. Bergson begrijpt deze als ``organisatie'' — dat wil zeggen, hij definieert deze door te verwijzen naar de wijze van zijn van levende wezens (\textit{être vivant}), een wijze waarin elk element representatief is voor het geheel (\textit{représentatif du tout}). Hij vergelijkt de innerlijke doordringing van alle elementen in het bewustzijn met de manier waarop alle noten vermengen wanneer we naar een melodie luisteren. Ook Bergson verdedigt dus het anti-Cartesiaanse element van het levensbegrip tegen de objectiverende wetenschap.

Als we nauwkeuriger kijken naar wat hier ``leven'' wordt genoemd en welke aspecten daarvan het begrip ervaring beïnvloeden, zien we dat de relatie tussen leven en ervaring niet die is van een universeel tot een bijzonder. In plaats daarvan staat de eenheid van de ervaring, bepaald door de intentionele inhoud, in een onmiddellijke relatie tot het geheel, tot de totaliteit van het leven. Bergson spreekt over de representatie van het geheel, en op soortgelijke wijze is Natorps begrip van onderlinge samenhang een uitdrukking van de ``organische'' relatie tussen deel en geheel die hier plaatsvindt. Het was met name Georg Simmel die het begrip leven in dit opzicht analyseerde als ``het uitsteken van het leven voorbij zichzelf''.

De representatie van het geheel in de momentane \textit{Erlebnis} gaat duidelijk veel verder dan het feit dat deze wordt bepaald door zijn object. Elke ervaring is, in Schleiermachers woorden, ``een element van oneindig leven''. Georg Simmel, die grotendeels verantwoordelijk was voor het feit dat het woord \textit{Erlebnis} zo populair werd, ziet het belangrijke aan het begrip ervaring als dit: ``het objectieve wordt niet alleen een beeld en idee, zoals in het kennen, maar een element in het levensproces zelf''. Hij zegt zelfs dat elke ervaring iets van een avontuur heeft. Maar wat is een avontuur? Een avontuur is beslist niet zomaar een episode. Episoden zijn een opeenvolging van details die geen innerlijke samenhang hebben en juist daardoor geen blijvende betekenis. Een avontuur daartegen onderbreekt de gebruikelijke gang van zaken, maar staat positief en betekenisvol in relatie tot de context die het onderbreekt. Zo laat een avontuur het leven als een geheel voelen, in zijn breedte en in zijn kracht. Hierin ligt de fascinatie van een avontuur. Het verwijdert de voorwaarden en verplichtingen van het alledaagse leven. Het waagt zich in het onzekere.

Maar tegelijkertijd weet het dat het, als avontuur, uitzonderlijk is en daardoor gerelateerd blijft aan de terugkeer naar het alledaagse, waartoe het avontuur niet kan behoren. Zo wordt het avontuur ``ondergaan'', als een beproeving of test waaruit men rijker en volwassener tevoorschijn komt.

Er zit een element hiervan in elke \textit{Erlebnis}. Elke ervaring wordt uit de continuïteit van het leven genomen en is tegelijkertijd gerelateerd aan het geheel van het leven. Het is niet zo dat een ervaring alleen levend blijft zolang deze niet volledig is geïntegreerd in het bewustzijn van het leven; de manier waarop deze ``bewaard en opgeheven'' (\textit{aufgehoben}) wordt door verwerkt te worden in het geheel van het levensbewustzijn, gaat veel verder dan elke ``betekenis'' die men erin zou kunnen zien. Omdat het zelf binnen het geheel van het leven is, is het geheel van het leven ook in het ervaren aanwezig.

Aan het einde van onze conceptuele analyse van ervaring kunnen we de verwantschap zien tussen de structuur van \textit{Erlebnis} als zodanig en de wijze van zijn van het esthetische. Esthetische ervaring is niet zomaar één soort ervaring onder andere, maar vertegenwoordigt de essentie van ervaring \textit{an sich}. Net zoals het kunstwerk als zodanig een wereld op zich is, zo is ook wat esthetisch wordt ervaren, als \textit{Erlebnis}, losgemaakt van alle verbindingen met de werkelijkheid. Het kunstwerk zou bijna per definitie een esthetische ervaring zijn: dat betekent echter dat de kracht van het kunstwerk de persoon die het ervaart plotseling uit de context van zijn leven rukt, en toch weer terugverbindt met het geheel van zijn bestaan. In de ervaring van kunst is een volheid van betekenis aanwezig die niet alleen tot deze specifieke inhoud of dit object behoort, maar die staat voor het betekenisvolle geheel van het leven. Een esthetische \textit{Erlebnis} bevat altijd de ervaring van een oneindig geheel. Juist omdat het niet combineert met andere ervaringen om één open ervaringsstroom te vormen, maar onmiddellijk het geheel vertegenwoordigt, is de betekenis ervan oneindig.

Aangezien esthetische ervaring, zoals hierboven gezegd, een voorbeeldig geval is van de betekenis van het begrip \textit{Erlebnis}, is het duidelijk dat het begrip \textit{Erlebnis} een bepalend kenmerk is van de grondslag van kunst. Het kunstwerk wordt begrepen als de voltooing van de symbolische representatie van het leven, en naar deze voltooing streeft elke ervaring al. Vandaar dat het zelf is aangewezen als het object van esthetische ervaring. Voor de esthetiek volgt hieruit de conclusie dat de zogenaamde \textit{Erlebniskunst} (kunst gebaseerd op ervaring) kunst \textit{per se} is.

\subsubsection{De grenzen van \textit{Erlebniskunst} en de rehabilitatie van de allegorie}

Het begrip \textit{Erlebniskunst} bevat een belangrijke dubbelzinnigheid. Oorspronkelijk betekende \textit{Erlebniskunst} duidelijk dat kunst voortkomt uit ervaring en een uitdrukking van ervaring is. Maar in afgeleide zin wordt het begrip \textit{Erlebniskunst} ook gebruikt voor kunst die bedoeld is om esthetisch ervaren te worden. Beide betekenissen hangen duidelijk met elkaar samen. De betekenis van datgene waarvan het wezen erin bestaat een ervaring uit te drukken, kan niet worden begrepen anders dan door een ervaring.

Zoals altijd in dergelijke gevallen, wordt het begrip \textit{Erlebniskunst} beïnvloed door de ervaring van de grenzen die eraan worden gesteld. Pas wanneer het niet meer vanzelfsprekend is dat een kunstwerk bestaat uit de transformatie van ervaringen — en wanneer het niet meer vanzelfsprekend is dat deze transformatie gebaseerd is op de ervaring van een geïnspireerd genie, dat met de zekerheid van een slaapwandelaar het kunstwerk schept, dat vervolgens een ervaring wordt voor de persoon die eraan wordt blootgesteld — wordt men zich bewust van het begrip \textit{Erlebniskunst} in zijn omtrek. De eeuw van Goethe komt ons opvallend voor door de vanzelfsprekendheid van deze aannames, een eeuw die een heel tijdperk, een epoche, vormt. Alleen omdat deze periode voor ons afgerond is en we er voorbij kunnen kijken, zijn we in staat om deze binnen haar eigen grenzen te zien en er een begrip van te vormen.

Langzaam realiseren we ons dat deze periode slechts een episode is in de totale geschiedenis van kunst en literatuur. Het monumentale werk van Curtius over middeleeuwse literatuuresthetiek geeft ons een goed beeld hiervan. Als we beginnen voorbij de grenzen van \textit{Erlebniskunst} te kijken en andere criteria hanteren, openen zich nieuwe gezichtspunten binnen de Europese kunst: we ontdekken dat van de klassieke oudheid tot aan de barok kunst werd gedomineerd door heel andere waardenormen dan die van het ervaren, en zo openen onze ogen zich voor volslagen onbekende artistieke werelden.

Natuurlijk kunnen ook deze voor ons ``ervaringen'' worden. Zo’n esthetisch zelfbegrip is altijd beschikbaar. Maar het kan niet worden ontkend dat het kunstwerk dat op deze manier een ervaring voor ons wordt, zelf niet zo bedoeld was om zo te worden begrepen. Genie en ervaren, onze waardenormen, zijn hier niet toereikend. We kunnen ons ook heel andere criteria herinneren en bijvoorbeeld zeggen dat het niet de echtheid van de ervaring of de intensiteit van de uitdrukking is, maar de ingenieuze hantering van vaste vormen en uitdrukkingswijzen die iets tot een kunstwerk maakt. Dit verschil in criteria geldt voor alle soorten kunst, maar is met name opvallend in de literatuur. Nog in de achttiende eeuw vinden we poëzie en retorica naast elkaar op een manier die voor het moderne bewustzijn verrassend is. Kant ziet in beide ``een vrij spel van de verbeelding en een serieuze zaak van het verstand''. Voor hem zijn zowel poëzie als retorica schone kunsten en zijn ze ``vrij'' voor zover beide de onbedoelde harmonie van beide kenvermogens, de zintuigen en het verstand, tentoonspreiden.

Tegen deze traditie in voert de waardering van ervaren en geïnspireerd genie onvermijdelijk een heel andere opvatting van ``vrije'' kunst in, waartoe poëzie alleen behoort voor zover deze alles louter toevalligs elimineert en retorica volledig verbant.

Zo volgt de devaluatie van de retorica in de negentiende eeuw noodzakelijkerwijs uit de leer dat genie onbewust schept. Laten we één bijzonder voorbeeld van deze devaluatie volgen: de geschiedenis van de begrippen symbool en allegorie, en de veranderende relatie tussen beide in de moderne tijd.

Zelfs geleerden die geïnteresseerd zijn in taalkundige geschiedenis, houden vaak onvoldoende rekening met het feit dat de esthetische tegenstelling tussen allegorie en symbool — die voor ons vanzelfsprekend lijkt — pas in de afgelopen twee eeuwen filosofisch is uitgewerkt, en zo weinig te verwachten was in eerdere tijden dat de vraag die gesteld moet worden veeleer is hoe de behoefte aan dit onderscheid en deze tegenstelling is ontstaan. Het mag niet worden vergeten dat Winckelmann, wiens invloed op de esthetiek en de filosofie van de geschiedenis van zijn tijd zeer groot was, beide begrippen synoniem gebruikte; en hetzelfde geldt voor de esthetiek van de achttiende eeuw in het algemeen. De betekenissen van de twee woorden hebben inderdaad iets gemeenschappelijks. Beide woorden verwijzen naar iets waarvan de betekenis niet ligt in de uiterlijke verschijning of klank, maar in een betekenis die daarbuiten ligt. Gemeenschappelijk is dat in beide gevallen het ene voor het andere staat. Deze betekenisrelatie, waarbij het niet-zintuiglijke zintuiglijk wordt gemaakt, komt voor in de poëzie en de beeldende kunsten, alsmede in de religieuze en sacramentale sfeer.

Een gedetailleerder onderzoek zou nodig zijn om te ontdekken in hoeverre het klassieke gebruik van de woorden ``symbool'' en ``allegorie'' de weg heeft geëffend voor de latere tegenstelling tussen beide waar we nu mee vertrouwd zijn. Hier kunnen we slechts een paar basislijnen schetsen. Natuurlijk hadden de twee woorden oorspronkelijk niets met elkaar te maken. ``Allegorie'' behoorde oorspronkelijk tot de sfeer van het spreken, van de \textit{logos}, en is derhalve een retorische of hermeneutische figuur. In plaats van wat eigenlijk wordt bedoeld, wordt iets anders, tastbaarders, gezegd, maar op zo’n manier dat het eerste wordt begrepen. ``Symbool'', daarentegen, is niet beperkt tot de sfeer van de \textit{logos}, want een symbool is niet gerelateerd aan een andere betekenis door zijn eigen betekenis, maar zijn eigen zintuiglijke bestaan heeft ``betekenis''. Als iets dat getoond wordt, stelt het iemand in staat iets anders te herkennen, zoals bij de \textit{tessera hospitalis} en dergelijke. Duidelijk is dat een symbool iets is dat niet alleen waarde heeft vanwege zijn inhoud, maar omdat het ``geproduceerd'' kan worden — dat wil zeggen, omdat het een document is waaraan de leden van een gemeenschap elkaar herkennen; of het nu een religieus symbool is of in een seculiere context verschijnt — als een insigne, een pas of een wachtwoord — in elk geval hangt de betekenis van het \textit{symbolon} af van zijn fysieke aanwezigheid en verkrijgt het een representatieve functie alleen door getoond of uitgesproken te worden.

Hoewel de twee begrippen, allegorie en symbool, tot verschillende sferen behoren, staan ze niet alleen door hun gemeenschappelijke structuur — het representeren van het ene door het andere — dicht bij elkaar, maar ook omdat beide hun hoofdtoepassing vinden in de religieuze sfeer. Allegorie ontstaat uit de theologische behoefte om stuitend materiaal uit een religieuze tekst — oorspronkelijk uit Homerus — te elimineren en geldige waarheden daarachter te herkennen. Het verkrijgt een correlatieve functie in de retorica waar omfloersing en indirecte uitspraken passender lijken. Het begrip symbool benadert nu dit retorisch-hermeneutische begrip van allegorie (symbool, in de zin van allegorie, lijkt voor het eerst voor te komen bij Chrysippus), met name door de christelijke transformatie van het neoplatonisme. Pseudo-Dionysius verdedigt aan het begin van zijn \textit{magnum opus} de noodzaak om symbolisch (\textit{symbolikos}) te werk te gaan door te verwijzen naar de onmeetbaarheid van het bovenzintuiglijke zijn van God met onze geest, die gewend is aan de wereld van de zintuigen. Zo verkrijgt \textit{symbolon} hier een anagogische functie; het leidt tot de kennis van het goddelijke — net zoals allegorische spraak leidt tot een ``hogere'' betekenis. De allegorische procedure van interpretatie en de symbolische procedure van kennis zijn beide noodzakelijk om dezelfde reden: het is alleen mogelijk het goddelijke te kennen door uit te gaan van de wereld van de zintuigen.

Maar het begrip symbool heeft een metafysische achtergrond die volledig ontbreekt in het retorische gebruik van allegorie. Het is mogelijk om voorbij het zintuiglijke naar het goddelijke geleid te worden. Want de wereld van de zintuigen is niet louter niets en duisternis, maar de uitstorting en weerspiegeling van de waarheid. Het moderne begrip symbool kan niet worden begrepen los van deze gnostische functie en zijn metafysische achtergrond. De enige reden dat het woord ``symbool'' verheven kan worden boven zijn oorspronkelijke gebruik (als document, teken of pas) naar het filosofische idee van een mysterieus teken, en daardoor vergelijkbaar wordt met een hiëroglief dat alleen door een ingewijde kan worden geïnterpreteerd, is dat het symbool geen willekeurig gekozen of gecreëerd teken is, maar een metafysische verbinding tussen zichtbaar en onzichtbaar veronderstelt. De onscheidbaarheid van zichtbare verschijning en onzichtbare betekenis, deze ``samenvalling'' van twee sferen, ligt ten grondslag aan alle vormen van religieuze verering. Het is gemakkelijk in te zien hoe de term werd uitgebreid naar de esthetische sfeer. Volgens Solger verwijst het symbolische naar een ``zijnde waarin de idee op de een of andere manier wordt herkend'' — dat wil zeggen, de innerlijke eenheid van ideaal en verschijning die specifiek is voor het kunstwerk. Allegorie daartegen creëert deze betekenisvolle eenheid alleen door te verwijzen naar iets anders.

Maar ook het begrip allegorie heeft een aanzienlijke uitbreiding ondergaan, voor zover allegorie niet alleen verwijst naar de stijlfiguur en de geïnterpreteerde zin (\textit{sensus allegoricus}), maar correlatief ook naar abstracte begrippen die kunstzinnig in beelden worden voorgesteld. Het is duidelijk dat de begrippen uit de retorica en de poëtica als model hebben gediend voor de ontwikkeling van esthetische begrippen in de sfeer van de beeldende kunsten. Het retorische element in het begrip allegorie draagt bij aan deze betekenisontwikkeling voor zover allegorie niet de soort oorspronkelijke metafysische verwantschap claimt die een symbool wel claimt, maar veeleer een coördinatie die door conventie en dogmatische overeenkomst tot stand is gekomen, en die het mogelijk maakt om iets beeldloos in beelden voor te stellen.

Kortom, de semantische tendensen aan het einde van de achttiende eeuw leidden ertoe dat het symbolische (opgevat als iets inherent en wezenlijk betekenisvols) werd gecontrasteerd met het allegorische, dat een externe en kunstmatige betekenis heeft. Het symbool is de samenvalling van het zintuiglijke en het niet-zintuiglijke; allegorie is de betekenisvolle relatie van het zintuiglijke tot het niet-zintuiglijke.

Nu, onder invloed van het begrip genie en de subjectivering van ``uitdrukking'', werd dit betekenisverschil een waardecontrast. Het symbool (dat onuitputtelijk kan worden geïnterpreteerd omdat het onbepaald is) wordt tegenover allegorie gesteld (begrepen als staand in een preciezere relatie tot betekenis en daardoor uitgeput) zoals kunst tegenover niet-kunst. Juist de onbepaaldheid van zijn betekenis bezorgde het woord en begrip ``symbolisch'' de overwinning toen de rationalistische esthetiek van de Verlichting bezweek voor de kritische filosofie en de esthetiek van het genie. Deze verbinding verdient nader onderzoek.

Kants logische analyse van het begrip symbool in §59 van de \textit{Kritiek van het oordeelsvermogen} wierp het duidelijkste licht op dit punt en was doorslaggevend: hij stelt symbolische en schematische voorstelling tegenover elkaar. Het symbolische is voorstelling (en niet louter notatie, zoals in het zogenaamde logische ``symbolisme''); maar symbolische voorstelling stelt een begrip niet direct voor (zoals het transcendentale schematisme in Kants filosofie doet), maar alleen op een indirecte manier, ``waardoor de uitdrukking niet het eigenlijke schema voor het begrip bevat, maar slechts een symbool voor de reflectie''. Dit begrip van symbolische voorstelling is een van de meest briljante resultaten van Kants denken. Hij doet zo recht aan de theologische waarheid die zijn scholastieke vorm had gevonden in de \textit{analogia entis} en houdt menselijke begrippen gescheiden van God. Daarbovenop ontdekt hij — met name verwijzend naar het feit dat ``deze zaak een diepgaander onderzoek vereist'' — de symbolische werkwijze van taal (haar consistente metaforiciteit); en ten slotte gebruikt hij het begrip analogie, in het bijzonder, om de relatie tussen het schone en het moreel goede te beschrijven, een relatie die noch ondergeschiktheid noch gelijkwaardigheid kan zijn. ``Het schone is het symbool van het moreel goede.'' In deze formule, voorzichtig als ze diepzinnig is, combineert Kant de eis van volledige vrijheid van reflectie in het esthetische oordeel met zijn menselijke betekenis — een idee dat van de grootste historische consequentie zou blijken te zijn.

Schiller volgde hem hierin. Toen hij het idee van een esthetische opvoe\-ding van de mensheid baseerde op de analogie tussen schoonheid en moraliteit die Kant had geformuleerd, kon Schiller een lijn volgen die Kant zelf had uitgezet: ``Smaak maakt de overgang van zintuiglijke aantrekkingskracht naar gewoontegetrouwe morele interesse mogelijk zonder, zo men zou kunnen zeggen, een te grote sprong.''

De vraag is: hoe zijn symbool en allegorie in de nu bekende tegenstelling gekomen? In eerste instantie vinden we niets van deze tegenstelling bij Schiller, ook al deelt hij de kritiek op de koude en kunstmatige allegorie die Klopstock, Lessing, de jonge Goethe, Karl Philipp Moritz en anderen in die tijd tegen Winckelmann richtten. Pas in de correspondentie tussen Schiller en Goethe vinden we de beginselen van het nieuwe begrip symbool. In zijn bekende brief van 17 augustus 1797 beschrijft Goethe de sentimentele stemming die bij hem is opgekomen door zijn indrukken van Frankfurt, en zegt over de objecten die deze stemming oproepen: ``dat zij eigenlijk symbolisch zijn — dat wil zeggen, zoals ik nauwelijks hoef uit te leggen, het zijn uitstekende voorbeelden die in een kenmerkende veelheid staan, als vertegenwoordigers van vele anderen, en een zekere totaliteit omvatten...'' Hij hecht waarde aan deze ervaring omdat deze bedoeld is om hem te helpen ontsnappen aan de ``miljoenkoppige hydra van het empirisme''. Schiller steunt hem hierin en vindt deze sentimentele wijze van voelen volledig in overeenstemming met ``wat wij op dit gebied hebben afgesproken''. Maar bij Goethe is het, zoals we weten, niet zozeer een esthetische ervaring als wel een ervaring van de werkelijkheid, en om deze te beschrijven schijnt hij het begrip ``symbolisch'' te ontlenen aan het vroege protestantse gebruik.

Schiller brengt idealistische bezwaren naar voren tegen het opvatten van de werkelijkheid als symbolisch, en duwt zo de betekenis van ``symbool'' in de richting van het esthetische. Goethes kunstvriend Meyer past het begrip symbool ook op het esthetische toe om het ware kunstwerk van allegorie te onderscheiden. Maar voor Goethe zelf is het contrast tussen symbool en allegorie in de kunsttheorie slechts een bijzonder geval van de algemene neiging naar betekenis die hij in alle verschijnselen zoekt. Zo past hij het begrip symbool toe op kleuren, omdat daar ``de ware relatie tegelijkertijd de betekenis uitdrukt''. Hier is de invloed van het traditionele hermeneutische schema van \textit{allegorice, symbolice, mystice} zo duidelijk dat hij ten slotte de zin schrijft die zo typisch voor hem is: ``Alles wat gebeurt is een symbool, en door zichzelf volledig voor te stellen, wijst het naar al het andere.''

In de filosofische esthetiek moet dit gebruik van het woord symbool zich hebben gevestigd via de Griekse ``godsdienst van de kunst''. Dit wordt duidelijk door Schellings ontwikkeling van de filosofie van de kunst uit de mythologie. In zijn \textit{Götterlehre} had Karl Philipp Moritz, op wie Schelling zich beriep, het ``oplossen'' van mythologische poëzie ``in louter allegorie'' afgewezen, maar toch gebruikte hij het woord ``symbool'' niet voor deze ``taal van de fantasie''. Schelling schrijft echter: ``Mythologie in het algemeen, en elk afzonderlijk stuk mythologische literatuur in het bijzonder, moet niet schematisch of allegorisch worden begrepen, maar symbolisch. Want de eis van absolute artistieke voorstelling is: een voorstelling met \textit{volledige onverschilligheid}, zodanig dat het algemene geheel het bijzondere \textit{is}, en het bijzondere tegelijkertijd geheel het algemene is, en het niet slechts betekent.''  Wanneer Schelling in zijn kritiek op Heines opvatting van Homerus aldus de ware relatie tussen mythologie en allegorie vaststelt, geeft hij tegelijkertijd het begrip symbool een centrale plaats binnen de filosofie van de kunst. Op dezelfde manier vinden we bij Solger de opvatting dat alle kunst symbolisch is. Solger is van mening dat het kunstwerk het bestaan van de ``idee'' zelf is — de betekenis ervan is geen ``idee die los van het werkelijke kunstwerk wordt gezocht''. Want dit is wat kenmerkend is voor het kunstwerk, de schepping van het genie: dat de betekenis ervan in het verschijnsel zelf ligt en niet willekeurig daarin wordt gelezen. Met verwijzing naar de Duitse vertaling van het woord ``symbool'' als \textit{Sinnbild} (zinnebeeld) beschrijft Schelling het als ``concreet, lijkend op niets anders dan op zichzelf, zoals een beeld, en toch universeel en vol betekenis als een begrip''. Inderdaad, wat het symbool, zelfs zoals Goethe het opvat, onderscheidt, is dat daarin de idee zelf zich bestaan geeft. Alleen omdat het begrip symbool de innerlijke eenheid van symbool en datgene wat wordt gesymboliseerd impliceert, was het mogelijk dat het symbool een universeel esthetisch grondbegrip werd. Een symbool is de samenvalling van zintuiglijke verschijning en bovenzintuiglijke betekenis, en deze samenvalling is, zoals de oorspronkelijke betekenis van het Griekse \textit{symbolon} en de voortzetting daarvan in de terminologie van verschillende religieuze denominaties, geen achteraf gecoördineerde relatie, zoals bij het gebruik van tekens, maar de vereniging van twee dingen die bij elkaar horen: alle symboliek, waardoor ``het priesterschap hogere kennis weerspiegelt'', berust veeleer op de ``oorspronkelijke verbinding'' tussen goden en mensen, schrijft Friedrich Creuzer, wiens \textit{Symbolik} de omstreden taak op zich nam om de raadselachtige symboliek van de oudheid te interpreteren.

Maar het begrip symbool werd niet zonder moeite uitgebreid tot een universeel esthetisch principe. Want de innerlijke eenheid van beeld en betekenis die het symbool kenmerkt, is niet eenvoudig. Het symbool lost de spanning tussen de wereld van de ideeën en de wereld van de zintuigen niet zomaar op: het wijst juist op een disproportie tussen vorm en wezen, uitdrukking en inhoud. Met name de religieuze functie van het symbool leeft van deze spanning. De mogelijkheid van de onmiddellijke en totale samenvalling van het schijnbare met het oneindige in een religieuze ceremonie veronderstelt dat het eindige en oneindige werkelijk bij elkaar horen. Zo correspondeert de religieuze vorm van het symbool precies met de oorspronkelijke aard van \textit{symbolon}, het verdelen van wat één is en het weer verenigen.

De disproportie tussen vorm en wezen is wezenlijk voor het symbool voor zover de betekenis van symbolen voorbij hun zintuiglijke verschijning wijst. Hier ontstaat die aarzeling, die onbeslistheid tussen vorm en wezen die kenmerkend is voor het symbool. Deze disproportie is kennelijk groter naarmate het symbool duisterder en betekenisvoller is — en kleiner naarmate de betekenis de vorm doordringt: dat was Creuzers idee. Hegel beperkt de term ``symbolisch'' tot de symbolische kunst van het Oosten juist vanwege deze disproportie tussen beeld en betekenis. Voor hem is een overschot aan betekenis kenmerkend voor een bepaalde kunstvorm, die verschilt van de klassieke kunst doordat deze voorbij deze disproportie is gekomen. Maar dit zeggen is kennelijk bewust het begrip vastleggen en kunstmatig beperken — een begrip dat, zoals we zagen, minder de disproportie dan de samenvalling van beeld en betekenis tracht uit te drukken. Het moet ook worden toegegeven dat Hegel, wanneer hij het begrip van het symbolische beperkt (ondanks zijn vele volgelingen), ingaat tegen de neiging van de moderne esthetiek, die (sinds Schelling) juist de eenheid van verschijning en betekenis in het symbolische heeft benadrukt om daardoor de esthetische autonomie te rechtvaardigen tegen de eisen van het begrip.

Laten we nu de overeenkomstige devaluatie van allegorie volgen. In het begin kan een factor zijn geweest het verlaten van het Franse classicisme in de Duitse esthetiek vanaf de tijd van Lessing en Herder. Toch gebruikt Solger de term ``allegorisch'' in een verheven zin voor de hele christelijke kunst, en Friedrich Schlegel gaat nog verder. Hij zegt: alle schoonheid is allegorie (\textit{Gespräch über Poesie}). Hegels gebruik van het begrip ``symbolisch'' (net als dat van Creuzer) staat nog heel dicht bij dit begrip van het allegorische. Maar het gebruik door de filosofen, gebaseerd op een romantische opvatting van de relatie tussen het onuitsprekelijke en de taal en op de ontdekking van de allegorische poëzie van het Oosten, werd niet behouden door de negentiende-eeuwse culturele humanisten. Men beriep zich op het classicisme van Weimar, en inderdaad was de afwaardering van allegorie de dominante zorg van het Duitse classicisme; die zorg was onvermijdelijk als gevolg van de opkomst van het begrip genie en de bevrijding van de kunst van de boeien van het rationalisme. Allegorie is zeker niet uitsluitend het product van genialiteit. Zij berust op vaste tradities en heeft altijd een vaste, onder woorden te brengen betekenis die zich niet verzet tegen rationeel begrip via het concept — integendeel, het begrip allegorie is nauw verbonden met dogmatiek: met de rationalisering van het mythologische (zoals in de Griekse Verlichting), of met de christelijke interpretatie van de Schrift in termen van leerstellige eenheid (zoals in de patristiek), en ten slotte met de verzoening van de christelijke traditie en de klassieke cultuur, die de basis vormt van de kunst en literatuur van het moderne Europa en wier laatste universele vorm de barok was. Met de afbraak van deze traditie was ook de allegorie ten einde. Zodra de kunst zich bevrijdde van alle dogmatische banden en kon worden gedefinieerd als de onbewuste schepping van genie, werd allegorie onvermijdelijk esthetisch verdacht.

Zo heeft Goethes werk in de esthetiek een sterke invloed gehad op het maken van het symbolische tot een positief en het allegorische tot een negatief artistiek begrip. Met name zijn eigen poëzie had dezelfde uitwerking, omdat deze werd gezien als de belijdenis van zijn leven, een poëtische vormgeving van ervaring (\textit{Erlebnis}). In de negentiende eeuw werd het criterium van ervaring, dat hij zelf had opgesteld, de hoogste waardenorm. In overeenstemming met de realistische geest van die eeuw werd alles in Goethes werk wat niet aan dit criterium volde — zoals de poëzie van zijn oude dag — afgedaan als allegorisch ``overladen''.

Uiteindelijk beïnvloedt dit ook de ontwikkeling van de filosofische esthetiek, die het begrip symbool in de universele, Goethiaanse zin aanvaardt, maar wier denken gebaseerd is op de tegenstelling tussen werkelijkheid en kunst — dat wil zeggen, dat het de dingen bekijkt vanaf het ``standpunt van de kunst'' en de negentiende-eeuwse esthetische religie van de cultuur. R. T. Vischer is typisch voor deze opvatting; hoe verder hij zich van Hegel verwijdert, des te meer hij Hegels begrip van symbool uitbreidt en het symbool ziet als een van de fundamentele prestaties van de subjectiviteit. De ``donkere symboliek van de geest'' geeft ziel en betekenis aan wat op zichzelf geen ziel heeft (natuur of verschijnselen). Omdat het esthetische bewustzijn — in tegenstelling tot het mythisch-religieuze — weet dat het vrij is, is de symboliek die het aan alles verleent ook ``vrij''. Hoe dubbelzinnig en onbepaald het symbool ook blijft, het kan niet langer worden gekarakteriseerd door zijn privatieve verhouding tot het begrip. Integendeel, het heeft zijn eigen positiviteit als schepping van de menselijke geest. Het is de volmaakte harmonie van verschijning en idee die nu — met Schelling — in het begrip symbool wordt benadrukt, terwijl dissonantie is voorbehouden aan allegorie of mythisch bewustzijn. Op dezelfde manier vinden we bij Cassirer dat esthetisch symbolisme zich onderscheidt van mythisch symbolisme door het feit dat in het esthetische symbool de spanning tussen beeld en betekenis in evenwicht is gebracht — een laatste echo van het classicistische begrip van de ``religie van de kunst''.

Uit dit overzicht van de taalkundige geschiedenis van symbool en allegorie trek ik een feitelijke conclusie. Het vaste contrast tussen de twee begrippen — het symbool dat ``organisch'' is ontstaan, en de koude, rationele allegorie — verliest aan kracht wanneer we de verbinding zien met de esthetiek van het genie en de ervaring (\textit{Erlebnis}). Als de herontdekking van barokkunst (die duidelijk zichtbaar is in de antiekmarkt) en, met name in de afgelopen decennia, de herontdekking van barokpoëzie, samen met modern esthetisch onderzoek, hebben geleid tot een zekere rehabilitatie van de allegorie, kunnen we nu de theoretische reden hiervoor inzien. De negentiende-eeuwse esthetiek was gebaseerd op de vrijheid van de symbolische activiteit van de geest. Maar is dat een voldoende grondslag? Is deze symbolische activiteit niet in feite ook beperkt door de voortdurende bestaan van een mythische, allegorische traditie?

Zodra dit wordt erkend, wordt het contrast tussen symbool en allegorie opnieuw relatief, terwijl het vooroordeel van de esthetiek van \textit{Erlebnis} het absolute leek te maken. Evenzo kan het verschil tussen esthetisch bewustzijn en mythisch bewustzijn nauwelijks als absoluut worden beschouwd.

We moeten erkennen dat het stellen van dergelijke vragen een fundamentele herziening van de grondbegrippen van de esthetiek noodzakelijk maakt. Kennelijk gaat het hier om meer dan slechts een verdere verandering in smaak en esthetische waarden. Integendeel, het begrip esthetisch bewustzijn zelf wordt twijfelachtig, en daardoor ook het standpunt van de kunst waartoe het behoort. Is de esthetische benadering van een kunstwerk de juiste? Of is wat we ``esthetisch bewustzijn'' noemen een abstractie? De herwaardering van de allegorie waar we over hebben gesproken, wijst erop dat er ook een dogmatisch element in het esthetische bewustzijn zit. En als het verschil tussen mythisch en esthetisch bewustzijn niet absoluut is, wordt dan niet het begrip kunst zelf problematisch? Want het is, zoals we hebben gezien, een product van het esthetische bewustzijn. In elk geval kan niet worden ontkend dat de grote periodes in de geschiedenis van de kunst die waren waarin mensen zonder enig esthetisch bewustzijn en zonder ons begrip van ``kunst'' zich omringden met scheppingen wier functie in het religieuze of seculiere leven door iedereen kon worden begrepen en die niemand louter esthetisch genoegen verschafte. Kan het begrip van de esthetische \textit{Erlebnis} op deze scheppingen worden toegepast zonder hun ware wezen geweld aan te doen?

\section{Het hernemen van de vraag naar artistieke waarheid}

\subsection{De twijfelachtigheid van het begrip esthetische vorming (Bildung)}

Om de omvang van deze vraag correct in te schatten, zullen we eerst een historische verkenning ondernemen om de specifieke, historisch ontwikkelde betekenis van het begrip ``esthetisch bewustzijn'' bloot te leggen. Het is duidelijk dat we vandaag de dag met ``esthetisch'' niet meer bedoelen wat Kant er nog mee associeerde, toen hij de leer van ruimte en tijd ``transcendentaal esthetiek'' noemde en de leer van het schone en het verhevene in natuur en kunst een ``kritiek van het esthetische oordeel''. Het keerpunt lijkt te zijn gekomen met Schiller, die het transcendentale idee van smaak omvormde tot een morele eis en het formuleerde als een imperatief: Leef esthetisch! In zijn esthetische geschriften nam Schiller de radicale subjectivering over, waardoor Kant het oordeel van smaak en diens claim op universele geldigheid transcendentale rechtvaardiging had gegeven, en veranderde deze van een methodologische vooronderstelling in een inhoudelijke.

Het is waar dat hij hierin Kant zelf kon volgen, voor zover Kant de smaak al de betekenis had toegekend van een overgang van zintuiglijk genot naar moreel gevoel. Maar toen Schiller verkondigde dat kunst de beoefening van vrijheid is, doelde hij meer op Fichte dan op Kant. Kant baseerde het a priori van smaak en genie op het vrije spel van de kennisvermogens. Schiller herinterpreteerde dit antropologisch in termen van Fichtes theorie van driften: de speeldrift moest de vormdrift en de stofdrift in harmonie brengen. Het cultiveren van de speeldrift is het doel van de esthetische opvoeding.

Dit had verstrekkende gevolgen. Want nu werd kunst, als de kunst van de schone schijn, gecontrasteerd met de praktische werkelijkheid en in termen van dit contrast begrepen. In plaats dat kunst en natuur elkaar aanvulden, zoals lang leek te gelden, werden ze tegenover elkaar gezet als schijn en werkelijkheid. Traditioneel was het doel van kunst — waartoe ook alle bewuste transformatie van de natuur voor menselijk gebruik behoort — om de leemtes die de natuur openliet aan te vullen en te vervolmaken. En de schone kunsten, zolang ze binnen dit kader werden gezien, waren een vervolmaking van de werkelijkheid, geen schijn die deze maskeert, verbergt of veredelt. Maar als het begrip kunst wordt gedefinieerd als schijn in contrast met werkelijkheid, dan vormt de natuur niet langer een alomvattend kader. Kunst wordt een eigen standpunt en vestigt een autonome claim op suprematie.

Waar kunst heerst, gelden de wetten van de schoonheid en worden de grenzen van de werkelijkheid overschreden. Dit ideale rijk moet worden verdedigd tegen alle inbreuken, zelfs tegen de moralistische voogdij van staat en samenleving. Het hoort waarschijnlijk bij de innerlijke verschuiving in de ontologische grondslag van Schillers esthetiek dat zijn grote plan in de Brieven over esthetische opvoeding tijdens de uitwerking verandert. Zoals bekend wordt een opvoeding door kunst een opvoeding tot kunst. In plaats dat kunst ons voorbereidt op ware morele en politieke vrijheid, hebben we de cultuur van een esthetische staat, een gecultiveerde samenleving die interesse toont in kunst. Maar dit roept een nieuw obstakel op voor het overwinnen van Kants dualisme tussen de wereld van de zintuigen en de wereld van de moraliteit, zoals deze worden overwonnen in de vrijheid van het esthetische spel en de harmonie van het kunstwerk. De verzoening van ideaal en leven door de kunst is slechts een lokale en tijdelijke verzoening. Schoonheid en kunst geven de werkelijkheid slechts een vluchtige, veredelende glans. De geestelijke vrijheid waartoe ze ons verheffen, is vrijheid in een esthetische toestand, niet in de werkelijkheid. Zo opent zich onder het dualisme van zijn en behoren te zijn, dat Kant esthetisch verzoent, een dieper, onopgelost dualisme. De poëzie van de esthetische verzoening moet haar eigen zelfbewustzijn zoeken tegenover de proza van de vervreemde werkelijkheid.

Het werkelijkheidsbegrip waartegen Schiller de poëzie stelt, is ongetwijfeld niet meer Kantiaans. Kant vertrekt immers, zoals we hebben gezien, altijd van de natuurlijke schoonheid. Maar doordat Kant, met het oog op de kritiek op de dogmatische metafysica, zijn kennisbegrip volledig beperkte tot de mogelijkheid van zuivere natuurwetenschap en daardoor de geldigheid van het nominalistische werkelijkheidsbegrip niet betwistte, gaat de ontologische moeilijkheid waarin de negentiende-eeuwse esthetiek zich bevond uiteindelijk op Kant zelf terug. Onder de heerschappij van nominalistische vooroordelen kan het esthetische zijn slechts onvolledig en onvolmaakt worden begrepen.

In wezen danken we onze bevrijding van de begrippen die een passend begrip van het esthetische zijn in de weg stonden, aan de fenomenologische kritiek op de negentiende-eeuwse psychologie en kennistheorie. Deze kritiek heeft de onjuistheid aangetoond van alle pogingen om de wijze van zijn van het esthetische te begrijpen in termen van de ervaring van de werkelijkheid en als een modificatie daarvan. Alle dergelijke ideeën als nabootsing, schijn, onwerkelijkheid, illusie, magie, droom, gaan ervan uit dat kunst betrokken is op iets anders dan zichzelf: de werkelijke werkelijkheid. Maar de fenomenologische terugkeer naar de esthetische ervaring leert ons dat deze laatste niet denkt in termen van deze relatie, maar integendeel datgene wat ze ervaart als echte waarheid beschouwt. Diens aard is derhalve zodanig dat deze niet teleurgesteld kan worden door een meer echte ervaring van de werkelijkheid. In tegenstelling daartoe hoort bij al de bovenvermelde modificaties van de werkelijkheidservaring wel een ervaring van teleurstelling. Wat slechts schijn was, onthult zich; wat aan werkelijkheid ontbeerde, verkrijgt deze; wat magisch was, verliest zijn magie; wat illusie was, wordt doorzien; en uit wat een droom was, ontwaken we. Als het esthetische slechts schijn in deze zin was, dan zou zijn kracht — zoals de angst van dromen — slechts kunnen duren zolang er geen twijfel bestond over zijn werkelijkheid, en zou het zijn waarheid verliezen bij het ontwaken.

De verschuiving in de ontologische definitie van het esthetische naar het begrip esthetische schijn heeft zijn theoretische basis in het feit dat de dominantie van het wetenschappelijke model van de kennistheorie ertoe leidt dat alle kennis mogelijkheden die buiten deze nieuwe methodologie vallen, in diskrediet worden gebracht.

Laten we ons herinneren dat Helmholtz in het bekende citaat waar we mee begonnen, geen betere manier wist om de kwaliteit te karakteriseren die het werk in de geesteswetenschappen onderscheidt van dat in de natuurwetenschappen dan door het te beschrijven als kunstzinnig. Positief correspondenteert hiermee wat we esthetisch bewustzijn kunnen noemen. Het is gegeven met het standpunt van de kunst, dat Schiller als eerste heeft gegrondvest. Want net zoals de kunst van de schone schijn zich verzet tegen de werkelijkheid, zo omvat het esthetische bewustzijn een vervreemding van de werkelijkheid — het is een vorm van de vervreemde geest, zoals Hegel de cultuur (Bildung) begreep. Het vermogen om een esthetische houding aan te nemen, maakt deel uit van het gecultiveerde bewustzijn. Want in het esthetische bewustzijn vinden we de kenmerken die het gecultiveerde bewustzijn onderscheiden: het opstijgen naar het universele, het distantiëren van de bijzonderheid van onmiddellijke aanvaarding of afwijzing, het respecteren van wat niet beantwoordt aan de eigen verwachting of voorkeur.

We hebben hierboven al de betekenis van het begrip smaak in deze context besproken. De eenheid van een smaakideaal dat een samenleving kenmerkt en haar leden verbindt, verschilt echter van wat de gestalte van de esthetische cultuur uitmaakt. Smaak gehoorzaamt nog steeds aan een inhoudelijk criterium. Wat in een samenleving als geldig wordt beschouwd, haar heersende smaak, ontleent zijn stempel aan de gemeenschappelijke kenmerken van het sociale leven. Een dergelijke samenleving kiest en weet wat bij haar hoort en wat niet. Zelfs haar artistieke interesses zijn niet willekeurig of in principe universeel, maar wat kunstenaars scheppen en wat de samenleving waardeert, behoren bij elkaar in de eenheid van een levensstijl en een smaakideaal.

In tegenstelling hiertoe bestaat het idee van esthetische cultivering — zoals we het uit Schiller hebben afgeleid — juist hierin dat elk inhoudelijk criterium wordt uitgesloten en het kunstwerk losgemaakt wordt van zijn wereld. Een uiting van deze losmaking is dat het domein waartoe het esthetisch gecultiveerde bewustzijn aanspraak maakt, universeel wordt uitgebreid. Alles waaraan het kwaliteit toekent, behoort ertoe. Het kiest niet meer, omdat het zelf niets is, noch streeft ernaar iets te zijn waarop een keuze gebaseerd zou kunnen zijn. Door reflectie is het esthetische bewustzijn voorbij elke bepalende en bepaalde smaak gegaan en vertegenwoordigt het zelf een totale onbepaaldheid. Het aanvaardt niet langer dat het kunstwerk en zijn wereld bij elkaar horen, maar integendeel: het esthetische bewustzijn is het ervarende centrum waarvandaan alles wat als kunst wordt beschouwd, wordt gemeten.

Wat we een kunstwerk noemen en esthetisch ervaren, hangt af van een proces van abstractie. Door alles te negeren waaraan een werk geworteld is — zijn oorspronkelijke levenscontext en de religieuze of seculiere functie die het betekenis gaf — wordt het zichtbaar als het zuivere kunstwerk. Door deze abstractie uit te voeren, verricht het esthetische bewustzijn een taak die op zichzelf positief is. Het toont wat een zuiver kunstwerk is en laat het in zijn eigen recht bestaan. Ik noem dit esthetische differentiatie.

Terwijl een bepaalde smaak differentieert — dat wil zeggen: selecteert en afwijst — op basis van een bepaalde inhoud, is esthetische differentiatie een abstractie die slechts selecteert op basis van esthetische kwaliteit als zodanig. Deze wordt uitgevoerd in het zelfbewustzijn van esthetische ervaringen. De esthetische ervaring is gericht op wat het eigenlijke werk zou moeten zijn — wat het negeert, zijn de extra-esthetische elementen die eraan kleven, zoals doel, functie, de betekenis van de inhoud. Deze elementen kunnen voldoende betekenis hebben voor zover ze het werk in zijn wereld situeren en aldus de hele zinvolheid bepalen die het oorspronkelijk bezat. Maar als kunst moet het werk van al deze elementen worden onderscheiden. Het is praktisch de definitie van esthetisch bewustzijn om te zeggen dat het datgene wat esthetisch bedoeld is, scheidt van alles wat buiten de esthetische sfeer valt. Het abstracteert van alle voorwaarden die de toegankelijkheid van een werk bepalen. Dit is dus een specifiek esthetische vorm van differentiatie. Het onderscheidt de esthetische kwaliteit van een werk van alle inhoudelijke elementen die ons ertoe aanzetten een morele of religieuze houding aan te nemen, en stelt het louter op zichzelf in zijn esthetische zijn.

Op dezelfde manier scheidt het in de uitvoerende kunsten tussen het origineel (toneelstuk of muziekcompositie) en de uitvoering ervan, en wel op zodanige wijze dat zowel het origineel (in tegenstelling tot de reproductie) als de reproductie op zichzelf (in tegenstelling tot het origineel of andere mogelijke interpretaties) als esthetisch kunnen worden gesteld. De soevereiniteit van het esthetische bewustzijn bestaat in zijn vermogen om deze esthetische differentiatie overal toe te passen en alles esthetisch te zien.

Omdat het esthetische bewustzijn claimt alles van artistieke waarde te omvatten, heeft het het karakter van simultaneïteit. Als esthetisch is zijn vorm van reflectie waarin het zich beweegt daarom niet alleen aanwezig. Want voor zover het esthetische bewustzijn alles wat het waardeert simultaan maakt, constituteert het zichzelf tegelijkertijd als historisch. Het is niet zo dat het historische kennis omvat en deze als onderscheidend kenmerk gebruikt; de opheffing van elke inhoudelijk bepaalde smaak, zoals eigen is aan de esthetische smaak, wordt ook expliciet zichtbaar in het scheppende werk van kunstenaars die zich tot het historische wenden. Het historische schilderij dat niet ontstaat uit een hedendaagse behoefte aan afbeelden, maar een voorstelling is in historische terugblik; de historische roman; en bovenal de historiserende vormen van de negentiende-eeuwse architectuur die zich overgaf aan voortdurende stijlenherinnering, tonen aan hoe nauw het esthetische en het historische bij een gecultiveerd bewustzijn bij elkaar horen.

Men zou kunnen tegenwerpen dat gelijktijdigheid niet voortkomt uit esthetische differentiatie, maar altijd al het resultaat is geweest van het samenbindende proces van het historische leven. Vooral grote architectonische werken blijven in het heden voortbestaan als levende getuigen van het verleden. Ook het bewaren van overgeleverde gewoonten, gedragingen, beelden en versieringen doet hetzelfde: het brengt een oudere manier van leven over naar die van het heden. Maar het esthetisch gecultiveerde bewustzijn verschilt hiervan. Het ziet zichzelf niet als deze vorm van integratie van de tijden; de simultaneïteit die er kenmerkend voor is, berust op het bewustzijn van de historische relativiteit van smaak. Feitelijke gelijktijdigheid wordt pas in principe simultaneïteit wanneer men fundamenteel bereid is om elke smaak die afwijkt van de eigen goede smaak niet te verwerpen. In plaats van de eenheid van een smaak hebben we nu een mobiel gevoel voor kwaliteit.

De esthetische differentiatie die het esthetische bewustzijn uitvoert, schept ook een externe existentie voor zichzelf. Het bewijst zijn productiviteit door speciale locaties voor simultaneïteit te reserveren: de universele bibliotheek op het gebied van literatuur, het museum, het theater, de concertzaal, enzovoort. Het is belangrijk om te zien hoe dit verschilt van wat ervoor kwam. Het museum is bijvoorbeeld niet simpelweg een openbare verzameling. Integendeel, de oudere verzamelingen (van hoven niet minder dan van steden) weerspiegelden de keuze van een bepaalde smaak en bevatten voornamelijk werken van dezelfde school, die als voorbeeldig werd beschouwd. Een museum daarentegen is een verzameling van dergelijke verzamelingen en vindt zijn volmaaktheid kenmerkend in het verbergen van het feit dat het uit dergelijke verzamelingen is voortgekomen, hetzij door de hele collectie historisch te herordenen, hetzij door deze zo uitgebreid mogelijk te maken. Op dezelfde manier zou men kunnen aantonen hoe de programma’s van permanent gevestigde theaters of concertzalen in de afgelopen eeuw steeds verder zijn afgedwaald van hedendaags werk en zich hebben aangepast aan de behoefte aan zelfbevestiging die kenmerkend is voor de gecultiveerde samenleving die deze instellingen ondersteunt. Zelfs kunstvormen als architectuur, die hiertegen lijken in te gaan, worden meegesleept in de simultaneïteit van de esthetische ervaring, hetzij door moderne reproductietechnieken, die gebouwen in beelden veranderen, hetzij door modern toerisme, dat reizen verandert in het bladeren door prentenboeken.

Zo verliest het werk door esthetische differentiatie zijn plaats en de wereld waartoe het behoort, voor zover het in plaats daarvan aan het esthetische bewustzijn toebehoort. Daardoor verliest ook de kunstenaar zijn plaats in de wereld. Dit blijkt uit de diskredietering van wat opdrachtkunst wordt genoemd. In een tijdperk waarin het openbare bewustzijn wordt gedomineerd door het idee dat kunst berust op ervaring (\textit{Erlebnis}), is het noodzakelijk om te herinneren dat scheppen uit vrije inspiratie — zonder opdracht, een gegeven thema en een gegeven aanleiding — vroeger de uitzondering was in het artistieke werk, terwijl we tegenwoordig het gevoel hebben dat een architect iemand sui generis is, omdat hij, in tegenstelling tot de dichter, schilder of componist, niet onafhankelijk is van opdracht en aanleiding. De vrije kunstenaar schept zonder opdracht. Hij lijkt onderscheiden door de volledige onafhankelijkheid van zijn creativiteit en verkrijgt daardoor de kenmerkende sociale trekken van een buitenstaander wiens levensstijl niet gemeten kan worden aan de normen van de openbare moraliteit. Het begrip boheem, dat in de negentiende eeuw ontstond, weerspiegelt dit proces. De woonplaats van de zigeuners werd het algemene woord voor de levensstijl van de kunstenaar. Maar tegelijkertijd draagt de kunstenaar, die vrij is als een vogel of een vis, de last van een roeping die hem tot een ambivalente figuur maakt. Want een gecultiveerde samenleving die losgemaakt is van haar religieuze tradities verwacht meer van de kunst dan het esthetische bewustzijn en het standpunt van de kunst kunnen leveren. De romantische eis voor een nieuwe mythologie — zoals deze door F. Schlegel, Schelling, Hölderlin en de jonge Hegel tot uiting werd gebracht, maar ook in de schilderijen en reflecties van Runge te vinden is — geeft de kunstenaar en zijn taak in de wereld het bewustzijn van een nieuwe wijding. Hij is iets als een wereldlijke verlosser (Immermann), want van zijn scheppingen wordt verwacht dat ze op kleine schaal de verzoening van de ramp bewerkstelligen waarnaar een onverloste wereld hoopt. Deze eis heeft sindsdien de tragedie van de kunstenaar in de wereld gedefinieerd, want elke vervulling ervan is altijd slechts een lokale, en in feite betekent dat deze wordt weerlegd. De experimentele zoektocht naar nieuwe symbolen of een nieuwe mythe die iedereen zal verenigen, kan zeker een publiek verzamelen en een gemeenschap scheppen, maar omdat elke kunstenaar zijn eigen gemeenschap vindt, getuigt de bijzonderheid van dergelijke gemeenschappen slechts van de desintegratie die plaatsvindt. Wat iedereen verenigt, is slechts de universele vorm van de esthetische cultuur. Hier is het werkelijke proces van cultivering — dat wil zeggen: de verheffing naar het universele — als het ware in zichzelf gedesintegreerd. De bereidheid van de intellectuele reflectie om in algemene begrippen te bewegen, om iets vanaf welk standpunt dan ook te bezien en het aldus met ideeën te bekleden, is, volgens Hegel, de manier om niet betrokken te raken bij de werkelijke inhoud van ideeën. Immermann noemt deze vrije zelfoverschrijding van de geest binnen zichzelf extravagant zelfgenot. Hij beschrijft hiermee de situatie die is ontstaan door de klassieke literatuur en filosofie van de tijd van Goethe, toen de epigonen alle vormen van de geest al bestaand vonden en derhalve het genot van cultuur in de plaats stelden van de echte prestatie ervan: het wegpoetsen van het vreemde en het ruwe. Het was gemakkelijk geworden om een goed gedicht te schrijven, en juist daardoor moeilijk om dichter te zijn.

\subsection{Kritiek op de abstractie die inherent is aan het esthetische bewustzijn}

Nu we de vorm hebben beschreven die het aannam als cultivering (\textit{Bildung}), zullen we het begrip esthetische differentiatie nader bekijken en de theoretische moeilijkheden bespreken die verbonden zijn met het \textit{begrip van het esthetische}. Abstractie tot het louter esthetische schaft het immers af. Dit wordt het meest duidelijk in Hamanns poging om een systematische esthetiek te ontwikkelen op basis van Kants onderscheidingen. Hamanns werk is opmerkelijk omdat hij daadwerkelijk teruggaat naar Kants transcendentale bedoeling en zo de eenzijdige toepassing van \textit{Erlebnis} als enige criterium voor kunst afbreekt. Door de implicaties van het esthetische element overal waar het zich voordoet te volgen, doet hij recht aan ook die bijzondere vormen van het esthetische die verbonden zijn met een doel, zoals de kunst van monumenten en affiches. Maar zelfs hier blijft Hamann trouw aan de taak van esthetische differentiatie. Want ook in deze vormen scheidt hij het esthetische van de niet-esthetische relaties waarin het staat, net zoals we buiten de kunstervaring kunnen zeggen dat iemand zich esthetisch gedraagt.
Zo krijgt het probleem van de esthetiek opnieuw zijn volle omvang, en wordt het transcendentale onderzoek hersteld dat was verlaten door het standpunt van de kunst en diens onderscheid tussen schone schijn en harde werkelijkheid. Esthetische ervaring is onverschillig voor de vraag of haar object werkelijk is, of de scène nu het toneel is of het echte leven. Het esthetische bewustzijn heeft onbeperkte soevereiniteit over alles.

Maar Hamanns poging mislukt aan de andere kant: in het begrip van kunst, dat hij, met perfecte consistentie, zo ver buiten het domein van het esthetische drijft dat het samenvloeit met virtuositeit. Hier wordt esthetische differentiatie tot het uiterste doorgevoerd. Het abstraheert zelfs van de kunst zelf.

Het fundamentele esthetische begrip waar Hamann van uitgaat, is dat waarneming in zichzelf betekenisvol is (\textit{Eigenbedeutsamkeit der Wahrnehmung}). Dit begrip betekent duidelijk hetzelfde als Kants theorie van de doelmatige harmonie met de toestand van onze kenvermogens. Net als bij Kant wordt voor Hamann het criterium van het begrip of de betekenis, dat essentieel is voor kennis, opgeschort. Taalkundig bezien is het woord \textit{Bedeutsamkeit} (de kwaliteit van betekenis of significantie) een afleiding van \textit{Bedeutung} en verplaatst het de associatie met een specifieke betekenis naar de sfeer van het onzekere. Iets is \textit{bedeutsam} als zijn betekenis (\textit{Bedeutung}) onuitgesproken of onbekend is. \textit{Eigenbedeutsamkeit} gaat echter nog een stap verder. Als iets \textit{eigenbedeutsam} (in zichzelf betekenisvol) is in plaats van \textit{fremdbedeutsam} (betekenisvol in relatie tot iets anders), dan koppelt het zich los van alles wat zijn betekenis zou kunnen bepalen. Kan zo’n begrip een stevige grondslag voor de esthetiek vormen? Kan men het begrip in zichzelf betekenisvol wel toepassen op een waarneming? Moeten we niet ook voor esthetische ervaring toestaan wat we over waarneming zeggen, namelijk dat deze waarheid waarneemt — dat wil zeggen: gerelateerd blijft aan kennis?

Het is de moeite waard om hier Aristoteles te herinneren. Hij toonde aan dat alle \textit{aisthesis} neigt naar het universele, ook al heeft elk zintuig zijn eigen specifieke domein en is daardoor wat onmiddellijk in de waarneming gegeven is niet universeel. Maar de specifieke zintuiglijke waarneming van iets als zodanig is een abstractie. Het feit is dat we zintuiglijke bijzonderheden zien in relatie tot iets universeels. Bijvoorbeeld: we herkennen een wit verschijnsel als een mens.

Nu wordt esthetisch zien zeker gekenmerkt doordat men niet haastig probeert wat men ziet te relateren aan het universele, de bekende betekenis, het beoogde doel, enzovoort, maar erbij stilstaat als iets esthetisch. Maar dat belet ons nog niet om relaties te zien — bijvoorbeeld om te herkennen dat dit witte verschijnsel dat we esthetisch bewonderen in feite een mens is. Zo is onze waarneming nooit een eenvoudige weerspiegeling van wat aan de zintuigen wordt gegeven.

Integendeel, we hebben van de moderne psychologie geleerd — met name van de scherpe kritiek die Scheler, maar ook W. Koehler, E. Strauss, M. Wertheimer en anderen uitoefenden op het begrip van zuivere waarneming als ``reactie op een prikkel'' — dat dit begrip zijn oorsprong vindt in een epistemologisch dogmatisme. De ware zin ervan is louter normatief: ``reactie op een prikkel'' is het ideale eindresultaat van de vernietiging van alle instinctieve fantasieën, het gevolg van een groot ontnuchteringsproces dat uiteindelijk in staat stelt om te zien wat er is, in plaats van de verbeeldingen van het instinct. Maar dat betekent dat zuivere waarneming, gedefinieerd als de adequate reactie op een prikkel, slechts een ideaal grensgeval is.

Er is echter een tweede punt. Zelfs waarneming, opgevat als een adequate reactie op een prikkel, zou nooit een loutere weerspiegeling zijn van wat er is. Want het zou altijd een begrip blijven van iets als iets. Alle begrijpen-als is een articulatie van wat er is, doordat het wegkijkt-van, kijkt-naar, ziet-samen-als. Al deze elementen kunnen het centrum van een observatie innemen of slechts ``meegaan'' met het zien, aan de rand of op de achtergrond. Zo is het onbetwist dat het zien, als een articulerende lezing van wat er is, veel van wat er is negeert, zodat het voor het zien simpelweg niet meer bestaat. Ook verwachtingen leiden ertoe dat men inleest wat er helemaal niet is. Laten we ook de neiging tot invariantie binnen het zien zelf niet vergeten, zodat men zoveel mogelijk altijd dezelfde dingen ziet.

Deze kritiek op de theorie van de zuivere waarneming, ondernomen op basis van pragmatische ervaring, werd vervolgens door Heidegger tot op de grondvesten onderzocht. Dit betekent echter dat deze kritiek ook opgaat voor het esthetische bewustzijn, ook al kijkt men hier niet simpelweg voorbij wat men ziet — bijvoorbeeld naar het algemene gebruik voor een doel — maar erbij stilstaat. Langdurig kijken en verwerken is geen eenvoudige waarneming van wat er is, maar is zelf begrijpen-als. De wijze van zijn van wat esthetisch wordt waargenomen, is niet voorhanden-zijn. In het geval van betekenisvolle representatie — bijvoorbeeld in werken van beeldende kunst, mits deze niet non-figuratief en abstract zijn — bepaalt het feit van hun betekenis duidelijk de manier waarop wat wordt gezien, wordt gelezen. Alleen als we ``herkennen'' wat wordt voorgesteld, zijn we in staat een schilderij te ``lezen''; in feite is dat uiteindelijk wat het tot een schilderij maakt. Zien betekent articuleren. Zolang we nog verschillende manieren proberen om wat we zien te ordenen of aarzelen tussen deze manieren, zoals bij bepaalde optische illusies, zien we nog niet wat er is. De optische illusie is als het ware de kunstmatige verlenging van deze aarzeling, de kwelling van het zien. Hetzelfde geldt voor het literair werk. Alleen wanneer we een tekst begrijpen — dat wil zeggen: ten minste zijn taal beheersen — kan deze voor ons een literair kunstwerk zijn. Zelfs bij het luisteren naar absolute muziek moeten we deze begrijpen. En alleen wanneer we deze begrijpen, wanneer deze duidelijk voor ons is, bestaat deze als een artistieke schepping voor ons. Hoewel absolute muziek een zuivere beweging van vorm op zich is, een soort auditieve wiskunde waarbij geen inhoud met een objectieve betekenis te onderscheiden is, omvat het begrijpen ervan niettemin het betreden van een relatie met wat betekenisvol is. Het is de onbepaaldheid van deze relatie die de specifieke relatie van dergelijke muziek tot betekenis markeert.

Zuiver zien en zuiver horen zijn dogmatische abstracties die verschijnselen kunstmatig reduceren. Waarneming omvat altijd betekenis. Daarom is het een doorgeschoten formalisme om de eenheid van het kunstwerk uitsluitend te zoeken in zijn vorm, tegenover zijn inhoud — een formalisme dat zich bovendien ten onrechte op Kant beroept.
Kant bedoelde met zijn begrip van vorm namelijk iets heel anders. Voor hem verwijst het begrip vorm naar de structuur van het esthetische object, niet als tegenstelling tot de betekenisvolle inhoud van een kunstwerk, maar als tegenstelling tot de louter zintuiglijke aantrekkingskracht van het materiaal. De zogenoemde objectieve inhoud is geen ruwe stof die pas later vorm krijgt, maar is in het kunstwerk al verbonden met de eenheid van vorm en betekenis.

Het woord motief, gangbaar in de taal van schilders, illustreert dit. Het kan figuratief zijn, maar ook abstract; maar in beide gevallen is het, ontologisch bezien, immaterieel (\textit{aneu hules}). Dat betekent geenszins dat het zonder inhoud is. Integendeel, wat een motief tot motief maakt, is dat het op overtuigende wijze eenheid bezit en dat de kunstenaar deze eenheid als de eenheid van een betekenis heeft doorgevoerd, net zoals de toeschouwer deze als eenheid begrijpt. In dit verband spreekt Kant van esthetische ideeën, waaraan veel onbenoembaars is toegevoegd. Dat is zijn manier om voorbij de transcendentale zuiverheid van het esthetische te gaan en de wijze van zijn van de kunst te erkennen. Zoals we hierboven hebben laten zien, was hij ver verwijderd van het vermijden van de intellectualisering van zuiver esthetisch genot. De arabesk is beslist niet zijn esthetische ideaal, maar slechts een geliefd methodologisch voorbeeld. Om recht te doen aan de kunst, moet de esthetiek voorbij zichzelf gaan en de zuiverheid van het esthetische opgeven. Maar zou dit haar werkelijk een stevige positie geven? Bij Kant had het begrip genie een transcendentale functie, en het begrip kunst was daarin gefundeerd. We zagen hoe dit begrip genie door zijn opvolgers werd uitgebreid tot de universele grondslag van de esthetiek. Maar is het begrip genie hiertoe wel geschikt?

Het moderne artistieke bewustzijn lijkt te suggereren van niet. Er lijkt een soort schemering van het genie te zijn ingetreden. Het idee van de slaapwandelende onbewustheid waarmee het genie schept — een idee dat overigens wel gerechtvaardigd kan worden door Goethes beschrijving van zijn eigen wijze van poëzie schrijven — lijkt vandaag de dag valse romantiek. Een dichter als Paul Valéry heeft daartegen het criterium gesteld van een kunstenaar en ingenieur als Leonardo da Vinci, in wiens totale genie vakmanschap, mechanische uitvinding en artistiek genie nog ongescheiden één waren. Het populaire bewustzijn wordt echter nog steeds beïnvloed door de achttiende-eeuwse verering van het genie en door de sacralisering van kunst, die — zoals we hebben gezien — kenmerkend werd voor de burgerlijke samenleving van de negentiende eeuw. Dit wordt bevestigd door het feit dat het begrip genie nu fundamenteel wordt begrepen vanuit het standpunt van de waarnemer. Dit oude begrip lijkt niet overtuigend voor de scheppende geest, maar wel voor de kritische geest. Het feit dat het werk voor de waarnemer een wonder lijkt, iets ondenkbaars voor iemand om te maken, weerspiegelt zich als een wonderbaarlijkheid van de schepping door geïnspireerd genie. Degenen die scheppen, passen vervolgens deze zelfde categorieën op zichzelf toe, en zo werd de geniecultus van de achttiende eeuw zeker ook door kunstenaars zelf gevoed. Maar zij zijn nooit zo ver gegaan in zelfverheffing als de burgerlijke samenleving hun zou hebben toegestaan. Het zelfbegrip van de kunstenaar blijft veel nuchterder. Hij ziet mogelijkheden van maken en doen, en vragen van techniek, waar de waarnemer inspiratie, mysterie en diepere betekenis zoekt.

Als men rekening wil houden met deze kritiek op de theorie van de onbewuste productiviteit van het genie, wordt men opnieuw geconfronteerd met het probleem dat Kant oploste door de transcendentale functie die hij aan het begrip genie toekende. Wat is een kunstwerk en hoe onderscheidt het zich van het product van een ambachtsman of zelfs van een slechte roman — dat wil zeggen: iets van mindere esthetische waarde? Voor Kant en het idealisme was het kunstwerk, per definitie, het werk van het genie. Zijn distinctiviteit — zijn volledige geslaagdheid en voorbeeldigheid — werd bewezen door het feit dat het aan genot en contemplatie een onuitputtelijk object van blijvende aandacht en interpretatie bood. Dat het genie van de schepping wordt geëvenaard door genie in de waardering, maakte al deel uit van Kants theorie van smaak en genie, en K.P. Moritz en Goethe onderwezen dit nog explicieter.

Maar hoe kan de aard van het artistieke genot en het verschil tussen wat een ambachtsman maakt en wat een kunstenaar schept, worden begrepen zonder het begrip genie? 

Hoe kan zelfs de volledigheid van een kunstwerk, het feit dat het af is, worden opgevat? De volledigheid van alles wat gemaakt of geproduceerd wordt, wordt gemeten aan de hand van het criterium van zijn doel — dat wil zeggen: deze wordt bepaald door het gebruik dat ervan gemaakt zal worden. Het werk is af wanneer het beantwoordt aan het doel waartoe het bestemd is. Maar hoe moet men het criterium bedenken om de volledigheid van een kunstwerk te meten? Hoe rationeel en nuchter men ook over artistieke productie mag denken, veel van wat we kunst noemen is niet bedoeld om gebruikt te worden, en geen ervan ontleent de norm voor zijn volledigheid aan zo’n doel. Schijnt het bestaan van het werk dan niet de afbreking van een scheppingsproces te zijn dat in wezen daarvoor uitsteekt? Misschien kan het in zichzelf helemaal niet worden voltooid?

Paul Valéry dacht inderdaad dat dit het geval was. Maar hij werkte de consequentie niet uit die hieruit voortvloeit voor iemand die een kunstwerk tegenkomt en poogt het te begrijpen. Als het waar is dat een kunstwerk in zichzelf niet voltooid kan zijn, wat is dan het criterium voor een passende ontvangst en begrip? Een willekeurig en arbitrair afgebroken scheppingsproces kan niets verplichtends inhouden. Hieruit volgt dat het aan de ontvanger moet worden overgelaten om iets van het werk te maken. De ene manier om een werk te begrijpen is dan niet minder legitiem dan de andere. Er is geen criterium voor een gepaste reactie. Niet alleen bezit de kunstenaar zelf dit niet — de esthetiek van het genie zou hierin instemmen; elke ontmoeting met het werk heeft de status en rechten van een nieuwe schepping. Dit lijkt mij een onhoudbaar hermeneutisch nihilisme. Als Valéry soms dergelijke conclusies trok voor zijn eigen werk om de mythe van de onbewuste productiviteit van het genie te vermijden, dan is hij, naar mijn mening, juist daarin verstrikt geraakt, want nu draagt hij de autoriteit van absolute schepping, die hij zelf niet langer wil uitoefenen, over aan de lezer en interpretator. Maar genie in het begrip is inderdaad niet meer waard dan genie in de schepping.

Dezelfde aporie doet zich voor als men uitgaat van het begrip esthetische ervaring in plaats van dat van genie. Over dit onderwerp onthult het fundamentele essay van Georg von Lukács, \textit{De subject-objectrelatie in de esthetiek}, het probleem. Hij schrijft een Heraclitische structuur toe aan de esthetische sfeer, waaraan hij bedoelt dat de eenheid van het esthetische object in werkelijkheid niet gegeven is. Het kunstwerk is slechts een lege vorm, een louter knooppunt in de mogelijke verscheidenheid van esthetische ervaringen (\textit{Erlebnisse}), en het esthetische object bestaat alleen in deze ervaringen. Zoals duidelijk is, is absolute discontinuïteit — dat wil zeggen: de desintegratie van de eenheid van het esthetische object in de veelvoudigheid van ervaringen — de noodzakelijke consequentie van een esthetiek van de \textit{Erlebnis}. Volgens de gedachten van Lukács heeft Oskar Becker ronduit gesteld dat ``het werk in tijdelijke zin slechts in een moment bestaat (namelijk: nu); het is nu dit werk en nu is het dit werk niet meer!'' Feitelijk is dat logisch. Een esthetiek die op ervaring is gebaseerd, leidt tot een absolute reeks punten, die de eenheid van het kunstwerk, de identiteit van de kunstenaar met zichzelf en de identiteit van de persoon die het kunstwerk begrijpt of ervan geniet, vernietigt.

Door de destructieve gevolgen van het subjectivisme te erkennen en de zelfvernietiging van de esthetische onmiddellijkheid te beschrijven, lijkt Kierkegaard mij de eerste te zijn geweest die de onhoudbaarheid van deze positie aantoonde. Zijn leer van het esthetische stadium van het bestaan is ontwikkeld vanuit het standpunt van de moralist die heeft gezien hoe wanhopig en onhoudbaar het bestaan in zuivere onmiddellijkheid en discontinuïteit is. Daarom is zijn kritiek op het esthetische bewustzijn van fundamenteel belang, omdat hij de innerlijke tegenspraken van het esthetische bestaan blootlegt, zodat het gedwongen wordt voorbij zichzelf te gaan. Omdat de esthetische bestaanswijze uiteindelijk onhoudbaar blijkt, beseffen we dat zelfs het verschijnsel van de kunst het menselijk bestaan een onontkoombare opgave stelt: het bereiken van een voortdurende samenhang in het zelfbegrip. Alleen zo’n continu zelfverstaan kan het menselijk bestaan dragen, ondanks de aanspraak van de alles opslorpende intensiteit van de vluchtige esthetische ervaring.

Als men toch nog het wezen van het esthetische bestaan op zo’n manier zou willen definieren dat het buiten de hermeneutische continuïteit van het menselijk bestaan wordt geconstrueerd, dan zou men, naar ik meen, het punt van Kierkegaards kritiek voorbij zijn. Toegegeven, het natuurlijke, als gezamenlijke voorwaarde van ons geestelijk leven, beperkt ons zelfbegrip en doet dat door zich op diverse manieren in het geestelijke te projecteren — als mythe, als droom, als de onbewuste voorvorming van het bewuste leven. En men moet toegeven dat esthetische verschijnselen op dezelfde manier de grenzen van het historische zelfbegrip van het \textit{Dasein} manifesteren. Maar we hebben geen standpunt dat ons in staat stelt deze grenzen en voorwaarden op zichzelf te zien of onszelf van buitenaf te zien als op deze manier beperkt en voorwaardelijk. Zelfs wat voor ons begrip gesloten is, ervaren wij zelf als beperkend, en bijgevolg behoort het nog steeds tot de continuïteit van het zelfbegrip waarin het menselijk bestaan zich beweegt. We herkennen ``de kwetsbaarheid van het schone en de avontuurlijkheid van de kunstenaar.'' Maar dat betekent niet dat we ons buiten een ``hermeneutische fenomenologie'' van het \textit{Dasein} bevinden. Integendeel, het stelt de taak om de hermeneutische continuïteit die ons zijn uitmaakt te bewaren, ondanks de discontinuïteit die inherent is aan het esthetische zijn en de esthetische ervaring.

Het pantheon van de kunst is geen tijdloos heden dat zich aan een zuiver esthetisch bewustzijn presenteert, maar de daad van een geest en een ziel die zich historisch heeft verzameld en gebundeld. Ook onze ervaring van het esthetische is een wijze van zelfbegrip. Zelfbegrip geschiedt altijd door het begrip van iets anders dan het zelf, en omvat de eenheid en integriteit van dat andere. Omdat we het kunstwerk in de wereld aantreffen en in het individuele kunstwerk een wereld tegenkomen, is het kunstwerk geen vreemde kosmos waartoe we voor even magisch worden getransporteerd. Integendeel, we leren onszelf in en door het kunstwerk begrijpen, en dat betekent dat we de discontinuïteit en het atomisme van geïsoleerde ervaringen opheffen in de continuïteit van ons eigen bestaan. Om deze reden moeten we een standpunt innemen ten opzichte van kunst en het schone dat niet doet alsof het onmiddellijk is, maar beantwoordt aan de historische aard van de menselijke conditie. De beroep op onmiddellijkheid, op de plotselinge flits van het genie, op de betekenis van ervaringen (\textit{Erlebnisse}), kan de eis van het menselijk bestaan naar continuïteit en eenheid van zelfbegrip niet weerstaan. De bindende kwaliteit van de ervaring (\textit{Erfahrung}) van kunst mag niet worden gedesintegreerd door het esthetische bewustzijn.

Dit negatieve inzicht, positief geformuleerd, luidt dat kunst kennis is en het ervaren van een kunstwerk betekent deelnemen aan die kennis.

Dit roept de vraag op hoe men recht kan doen aan de waarheid van de esthetische ervaring (\textit{Erfahrung}) en de radicale subjectivering van het esthetische, die begon met Kants \textit{Kritiek van het Oordeelsvermogen}, kan overwinnen. We hebben laten zien dat het een methodologische abstractie was, die correspondeerde met een heel specifieke transcendentale taak van fundering, die Kant ertoe bracht het esthetische oordeel volledig te relateren aan de toestand van het subject. Als deze esthetische abstractie echter vervolgens als inhoud werd begrepen en veranderde in de eis dat kunst louter esthetisch moest worden begrepen, dan kunnen we nu zien hoe deze eis tot abstractie in onoplosbare tegenspraak kwam met de ware ervaring van kunst.

Moet er dan in de kunst geen kennis zijn? Bevat de ervaring van kunst dan geen claim op waarheid die zeker anders is dan die van de wetenschap, maar die even zeker niet onderdoet voor die van de wetenschap? En is de taak van de esthetiek niet juist om te funderen dat de ervaring (\textit{Erfahrung}) van kunst een wijze van kennis is van een unieke aard, zeker anders dan die zintuiglijke kennis die de wetenschap de uiteindelijke gegevens verschaft waaruit zij de kennis van de natuur construeert, en zeker anders dan alle morele rationele kennis, en inderdaad anders dan alle begrippelijke kennis — maar toch kennis, dat wil zeggen: waarheid overbrengend?

Dit kan nauwelijks worden erkend als men, met Kant, de waarheid van kennis meet aan het wetenschappelijke begrip van kennis en het wetenschappelijke begrip van werkelijkheid. Het is noodzakelijk het begrip ervaring (\textit{Erfahrung}) ruimer te nemen dan Kant deed, zodat de ervaring van het kunstwerk als ervaring kan worden begrepen. Hiervoor kunnen we een beroep doen op Hegels bewonderenswaardige colleges over esthetiek. Hierin wordt de waarheid die in elke artistieke ervaring ligt opgesloten, erkend en tegelijkertijd bemiddeld met het historische bewustzijn. Derhalve wordt de esthetiek een geschiedenis van wereldbeelden — dat wil zeggen: een geschiedenis van waarheid, zoals deze zich manifesteert in de spiegel van de kunst. Het is ook een fundamentele erkenning van de taak die ik als volgt formuleerde: de kennis van waarheid die plaatsvindt in de ervaring van kunst zelf, te legitimeren.

Het bekende begrip wereldbeeld — dat voor het eerst bij Hegel in de \textit{Fenomenologie van de geest} opduikt als term voor Kants en Fichtes postulatorische uitbreiding van de fundamentele morele ervaring tot een morele wereldorde — verkrijgt zijn bijzondere stempel eerst in de esthetiek. Het is de veelvoudigheid en de mogelijke verandering van wereldbeelden die het begrip wereldbeeld zijn vertrouwde klank heeft gegeven. Maar de geschiedenis van de kunst is het beste voorbeeld hiervan, omdat deze historische veelvoudigheid niet kan worden overtroffen door vooruitgang naar de ene, ware kunst. Toegegeven, Hegel kon de waarheid van de kunst slechts erkennen door deze ondergeschikt te maken aan de alomvattende kennis van de filosofie en door de geschiedenis van wereldbeelden, net als de wereldgeschiedenis en de geschiedenis van de filosofie, te construeren vanuit het standpunt van het volledige zelfbewustzijn van het heden. Maar dit kan niet eenvoudigweg als een verkeerde afslag worden beschouwd, want de sfeer van het subjectieve bewustzijn is ver overtroffen. Hegels stap daarbovenuit blijft een blijvend element van waarheid in zijn denken. Zeker, voor zover het de begrippelijke waarheid almachtig maakt, omdat het begrip alle ervaring overstijgt, verloochent Hegels filosofie tegelijkertijd de weg naar waarheid die zij in de ervaring van kunst heeft erkend. Als we de kunst als een eigen weg naar waarheid willen rechtvaardigen, dan moeten we volledig realiseren wat waarheid hier betekent.

Het is in de geesteswetenschappen als geheel dat een antwoord op deze vraag moet worden gevonden. Want deze streven er niet naar om de veelvoudigheid van ervaringen — of deze nu esthetisch, historisch, religieus of politiek van aard zijn — te overstijgen, maar om deze te begrijpen, en dat betekent dat zij er waarheid in verwachten te vinden. We zullen moeten ingaan op de relatie tussen Hegel en het zelfbegrip van de geesteswetenschappen, vertegenwoordigd door de ``historische school'', en ook op de manier waarop beiden verschillen over wat het mogelijk maakt om correct te begrijpen wat waarheid in de geesteswetenschappen betekent. In ieder geval zullen we het probleem van de kunst niet vanuit het standpunt van het esthetische bewustzijn recht kunnen doen, maar alleen binnen dit ruimere kader.

We hebben slechts één stap in deze richting gezet door de zelfinterpretatie van het esthetische bewustzijn te corrigeren en de vraag naar de waarheid van de kunst, waartoe de esthetische ervaring getuigt, te herstellen. Onze zorg is dan ook om de ervaring van kunst op zo’n manier te bezien dat deze als ervaring (\textit{Erfahrung}) wordt begrepen. De ervaring van kunst mag niet worden vervalst door deze om te vormen tot een bezit van de esthetische cultuur, waardoor haar bijzondere claim wordt geneutraliseerd. We zullen zien dat dit een verstrekkende hermeneutische consequentie inhoudt, want elke ontmoeting met de taal van de kunst is een ontmoeting met een onafgemaakt gebeuren en is zelf deel van dit gebeuren. Dit is wat benadrukt moet worden tegenover het esthetische bewustzijn en diens neutralisatie van de vraag naar waarheid.

% t/m p. 85
Als het speculatieve idealisme probeerde het esthetische subjectivisme en het op Kant gebaseerde agnosticisme te overwinnen door zich te verheffen tot het standpunt van oneindige kennis, dan bracht deze gnostische zelfverlossing van de eindigheid — zoals we hebben gezien — met zich mee dat de kunst werd opgeheven in de filosofie.
Wij daarentegen zullen vast moeten houden aan het standpunt van de eindigheid. Het lijkt me dat het productieve aan Heideggers kritiek op het moderne subjectivisme is dat zijn tijdsinterpretatie van het zijn nieuwe mogelijkheden heeft geopend.
Het zijn interpreteren vanaf de horizon van de tijd betekent niet, zoals constant verkeerd wordt begrepen, dat het \textit{Dasein} radicaal tijdelijk is, zodat het niet langer als eeuwig of tijdloos kan worden beschouwd, maar alleen begrepen kan worden in relatie tot zijn eigen tijd en toekomst. Als dit de betekenis was, zou het geen kritiek en overwinnng van het subjectivisme zijn, maar een existentialistische radicalisering ervan, waarvan men gemakkelijk zou kunnen voorspellen dat deze een collectivistische toekomst zou hebben. De filosofische vraag die hier echter aan de orde is, is juist gericht op dit subjectivisme zelf. Dit wordt tot zijn uiterste punt gedreven alleen om het in vraag te stellen. De filosofische vraag luidt: wat is het zijn van het zelfbegrip? Met deze vraag overschrijdt het fundamenteel de horizon van dit zelfbegrip. Door de tijd te onthullen als de grondslag die voor het zelfbegrip verborgen is, predikt het geen blinde toewijding uit nihilistische wanhoop, maar opent het zich voor een tot nu toe verborgen ervaring die het denken vanaf de positie van de subjectiviteit te boven gaat, een ervaring die Heidegger \textit{zijn} noemt.

Om recht te doen aan de ervaring (\textit{Erfahrung}) van de kunst zijn we begonnen met een kritiek op het esthetische bewustzijn. De ervaring van kunst erkent dat zij de volledige waarheid van wat zij ervaart niet in termen van definitieve kennis kan presenteren. Er is geen absolute vooruitgang en geen definitieve uitputting van wat in een kunstwerk ligt opgesloten. De ervaring van kunst weet dit van zichzelf.
Tegelijkertijd kunnen we niet eenvoudigweg accepteren wat het esthetische bewustzijn als zijn ervaring beschouwt. Want zoals we hebben gezien, beschouwt het uiteindelijk zijn ervaring als de discontinuïteit van ervaringen (\textit{Erlebnisse}). Maar deze conclusie hebben we onaanvaardbaar gevonden.

We vragen dus niet aan de ervaring van kunst hoe deze zichzelf opvat, maar wat deze werkelijk is en wat haar waarheid is, ook al weet zij niet wat zij is en niet kan zeggen wat zij weet — net zoals Heidegger heeft gevraagd wat metafysica is, in tegenstelling tot wat deze zichzelf denkt te zijn. In de ervaring van kunst zien we een echte ervaring (\textit{Erfahrung}) die door het werk wordt opgeroepen en die degene die deze ervaring heeft, niet onveranderd laat, en we onderzoeken de wijze van zijn van wat op deze manier wordt ervaren. Zo hopen we beter te begrijpen wat voor soort waarheid ons daar tegemoet treedt.

We zullen zien dat dit de dimensie opent waarin, in het begrip zoals dat door de geesteswetenschappen wordt beoefend, de vraag naar de waarheid op een nieuwe manier wordt gesteld.

Als we willen weten wat waarheid is op het gebied van de geesteswetenschappen, zullen we de filosofische vraag naar de hele procedure van de geesteswetenschappen op dezelfde manier moeten stellen als Heidegger deze aan de metafysica stelde en wij aan het esthetische bewustzijn. Maar we zullen niet eenvoudigweg de eigen opvatting van de geesteswetenschappen over zichzelf kunnen accepteren, maar moeten vragen wat hun wijze van begrip in waarheid is. De vraag naar de waarheid van de kunst in het bijzonder kan dienen om de weg te bereiden voor deze ruimere vraag, omdat de ervaring van het kunstwerk begrip inhoudt en dus zelf een hermeneutisch fenomeen vertegenwoordigt — maar geenszins in de zin van een wetenschappelijke methode. Integendeel, begrip behoort tot de ontmoeting met het kunstwerk zelf, en deze toebehoren kan alleen worden verhelderd op basis van de wijze van zijn van het kunstwerk zelf.





